% Generated human-review packet template.  Source records, Lean statements,
% and raw-source-to-Spec screening fields are substituted by
% scripts/review_dashboard_packet.py.
% Do not hand-edit generated packets; annotate the PDF or reconcile notes into
% the paper-local human-review log.
\documentclass[10pt]{article}
\usepackage[margin=0.43in]{geometry}
\usepackage{fontspec}
\setmainfont{Latin Modern Roman}
\setmonofont{DejaVu Sans Mono}
\usepackage{hyperref}
\usepackage{amssymb}
\usepackage{fvextra}
\usepackage{enumitem}
\setlength{\parindent}{0pt}
\setlength{\parskip}{0.15em}
\linespread{0.92}
\hypersetup{colorlinks=true,linkcolor=blue,urlcolor=blue}
\DefineVerbatimEnvironment{ReviewVerbatim}{Verbatim}{breaklines=true,breakanywhere=true,fontsize=\scriptsize}
\newcommand{\reviewerbox}{%
  \par\noindent\fbox{\parbox{0.96\linewidth}{\rule{0pt}{0.38in}}}\par
}
\newcommand{\reviewmatch}[1]{%
  \par\noindent\textbf{Reviewer check:} $\square$\; #1\par
  \noindent\emph{If left unchecked, explain the discrepancy or uncertainty below.}\par
}

\begin{document}
\fontsize{9}{10}\selectfont

\begin{center}
{\LARGE Human review packet: GJ19OptimalBinaryRatingSystems}\\[0.35em]
\small Generated from byte-pinned source excerpts, the expanded PaperInterface
specifications, and hash-bound raw-source-to-Spec screening records.
\end{center}

\noindent\textbf{Paper:} GJ19OptimalBinaryRatingSystems\\
\textbf{Paper id:} \texttt{GJ19OptimalBinaryRatingSystems}\\
\textbf{Source version:} not recorded\\
\textbf{Source record:} \texttt{audit/paper\_statement\_map.json}\\
\textbf{Rows:} 59\\
\textbf{Generated:} 2026-08-18\\
\textbf{Source URL:} \href{https://proceedings.mlr.press/v89/garg19a.html}{official source}

\noindent\fbox{\parbox{0.96\linewidth}{\textbf{Review status:} 0/59 source-claim screens; 0/62 library prerequisite screens. This is a review worksheet while those direct checks remain pending; it is not a final validation receipt.}}\par



\section*{How to use this packet}
For each source claim, compare the exact source excerpts with the expanded
PaperInterface specification. Mark up this PDF or fill the blank reviewer box in the TeX source;
return the annotated file for reconciliation into the paper-local human-review
log.

\section*{Open the interactive dashboard}
The interactive dashboard is an optional alternative to using this PDF; it
presents the same review surface rather than an additional review requirement.
After forking and cloning (or pulling) EconCSLib, run the following from the
repository root. The cache command is needed only the first time in a new clone.
\begin{ReviewVerbatim}
cd <your-EconCSLib-checkout>
lake exe cache get
python3 scripts/review_dashboard.py --paper GJ19OptimalBinaryRatingSystems --serve
\end{ReviewVerbatim}
Open \url{http://127.0.0.1:8765/} in a browser and keep that terminal running
while reviewing. The dashboard records saved annotations in this paper's
reviewer-owned review record.

\section*{Contents}
\begin{itemize}[leftmargin=1.2em]
\item \hyperlink{library-section-a85f9574c44a156c}{Semantic library prerequisites (62; not paper claims)}
\begin{itemize}[leftmargin=1.2em]
\item \hyperlink{library-prerequisite-29ba02126b109d91}{EconCSLib.Optimization.IsLexicographicMaximizerOn}
\item \hyperlink{library-prerequisite-0a2e53c8e507cb74}{EconCSLib.Optimization.IsMaximizerOn}
\item \hyperlink{library-prerequisite-9e10d8e3c1b32191}{EconCSLib.Optimization.IsMinimizerOn}
\item \hyperlink{library-prerequisite-8d013dcbf8315291}{EconCSLib.Optimization.nestedBisectionStepBound}
\item \hyperlink{library-prerequisite-33a658a76ab86d8b}{EconCSLib.Optimization.realBisectionMidpoint}
\item \hyperlink{library-prerequisite-f7e825d2da547274}{EconCSLib.Optimization.realBisectionStep}
\item \hyperlink{library-prerequisite-e24bfb229e6e0f0c}{EconCSLib.Optimization.realBisectionStepFn}
\item \hyperlink{library-prerequisite-5b79c0723b67e9a9}{EconCSLib.Optimization.realBisectionRun}
\item \hyperlink{library-prerequisite-e4ddb15be8735d63}{EconCSLib.Probability.ExponentialRateCertificate}
\item \hyperlink{library-prerequisite-8386068971d3a1c1}{EconCSLib.Probability.FiniteMeasurableSetPartition}
\item \hyperlink{library-prerequisite-94654614a042c92d}{EconCSLib.Probability.FiniteMeasurableSetPartition.mk}
\item \hyperlink{library-prerequisite-6f62f2ee03548cd2}{EconCSLib.Probability.FiniteMeasurableSetPartition.orderedRealIocNatCutpoints}
\item \hyperlink{library-prerequisite-aab7fae9d893ec88}{EconCSLib.Probability.FiniteMeasurableSetPartition.orderedRealIocNatCutpointsMidpoint}
\item \hyperlink{library-prerequisite-0b261505488b443f}{EconCSLib.Probability.FiniteMeasurableSetPartition.orderedRealIocNatCutpointsProdMidpoint}
\item \hyperlink{library-prerequisite-5833c2529d9caeeb}{EconCSLib.Probability.FiniteMeasurableSetPartition.pieceSet}
\item \hyperlink{library-prerequisite-a2c998741ee4703f}{EconCSLib.Probability.FiniteMeasurableSetPartition.piecewiseConstKernel}
\item \hyperlink{library-prerequisite-a9f78bbf3e21bb13}{EconCSLib.Probability.FiniteMeasurableSetPartition.support}
\item \hyperlink{library-prerequisite-3161aa1be2e11c39}{EconCSLib.Probability.FiniteMeasurableSetPartition.prod}
\item \hyperlink{library-prerequisite-e99ffb1d85bad07e}{EconCSLib.Probability.FiniteMeasurableSetPartition.subpartition}
\item \hyperlink{library-prerequisite-cb9e54451ad9288c}{EconCSLib.Probability.FiniteMeasurableSetPartition.selectedProduct}
\item \hyperlink{library-prerequisite-943d19d91afe6d63}{EconCSLib.Probability.FiniteRatingLDPModel}
\item \hyperlink{library-prerequisite-c3203b699a19e900}{EconCSLib.Probability.FiniteRatingLDPModel.score}
\item \hyperlink{library-prerequisite-fc314aabfaabee42}{EconCSLib.Probability.FiniteRatingLDPModel.typeLaw}
\item \hyperlink{library-prerequisite-6258b3da52471ab5}{EconCSLib.Probability.finiteMGF}
\item \hyperlink{library-prerequisite-20dbd208a81d9e80}{EconCSLib.Probability.finiteLogMGF}
\item \hyperlink{library-prerequisite-c959d6e7b586c2bc}{EconCSLib.Probability.FiniteRatingLDPModel.logMGF}
\item \hyperlink{library-prerequisite-bd97fc0d50344aa7}{EconCSLib.Probability.FiniteRatingLDPModel.mk}
\item \hyperlink{library-prerequisite-d169ba913523821f}{EconCSLib.Probability.HasAEEssentialInfimum}
\item \hyperlink{library-prerequisite-94743506c5de049b}{EconCSLib.Probability.logDecay}
\item \hyperlink{library-prerequisite-380666720005662c}{EconCSLib.Probability.HasExponentialRate}
\item \hyperlink{library-prerequisite-e43b036a6f31f38f}{EconCSLib.Probability.HasPositiveWeightNearAEEssentialInfimum}
\item \hyperlink{library-prerequisite-2f4464ce60bae569}{EconCSLib.Probability.PairwiseThresholdRateTopLdpCertificate}
\item \hyperlink{library-prerequisite-479fa520bc58d797}{EconCSLib.Probability.bernoulliFirstEndpointEqualSplit}
\item \hyperlink{library-prerequisite-c4d922c1dbf3d0a6}{EconCSLib.Probability.bernoulliInteriorFailureSplitNumerator}
\item \hyperlink{library-prerequisite-48bd8c17f399b1fc}{EconCSLib.Probability.bernoulliInteriorSuccessSplitNumerator}
\item \hyperlink{library-prerequisite-17d704ff7a24219b}{EconCSLib.Probability.bernoulliInteriorSplitDenominator}
\item \hyperlink{library-prerequisite-57b352998b6e4a14}{EconCSLib.Probability.bernoulliInteriorEqualSplit}
\item \hyperlink{library-prerequisite-ca060c260b48fe49}{EconCSLib.Probability.bernoulliKL}
\item \hyperlink{library-prerequisite-2735ba3861b2f309}{EconCSLib.Probability.bernoulliLastEndpointEqualSplit}
\item \hyperlink{library-prerequisite-4cea2ad3a0ac0028}{EconCSLib.Probability.binaryRatingScore}
\item \hyperlink{library-prerequisite-e0e20858ab2c39b5}{EconCSLib.Probability.finiteIidScoreSum}
\item \hyperlink{library-prerequisite-f584f3b3a668a101}{EconCSLib.Probability.floorSampleCount}
\item \hyperlink{library-prerequisite-b6ce97a274cbdb86}{EconCSLib.Probability.realBernoulliPMF}
\item \hyperlink{library-prerequisite-342012c3b44e591d}{EconCSLib.Probability.realBinaryRatingLDPModel}
\item \hyperlink{library-prerequisite-10ffb019e9ff740e}{EconCSLib.Probability.twoBernoulliThresholdRate}
\item \hyperlink{library-prerequisite-f185c8e6394c9317}{EconCSLib.pmfProd}
\item \hyperlink{library-prerequisite-38eb9a31a47bfbd9}{EconCSLib.pmfProduct}
\item \hyperlink{library-prerequisite-59131a14024b99ca}{EconCSLib.Probability.twoSampleRatingLaw}
\item \hyperlink{library-prerequisite-0777cfce9fa4aa84}{EconCSLib.Probability.twoSampleScoreGapSum}
\item \hyperlink{library-prerequisite-e5c7eedc40c11b7a}{EconCSLib.pmfExp}
\item \hyperlink{library-prerequisite-7fec9dedef722fc8}{EconCSLib.pmfProb}
\item \hyperlink{library-prerequisite-7977d3effede269a}{EconCSLib.Probability.twoSampleScoreGapStrictLeftProb}
\item \hyperlink{library-prerequisite-04fdd75df8db4e26}{EconCSLib.Probability.twoSampleScoreGapTieProb}
\item \hyperlink{library-prerequisite-ba4c2212912ca97d}{EconCSLib.Probability.twoSamplePkComplementErrorProb}
\item \hyperlink{library-prerequisite-94bd3b0a19bf1178}{EconCSLib.Probability.twoSampleFloorPkComplementErrorProb}
\item \hyperlink{library-prerequisite-11866f137653a75d}{EconCSLib.Probability.weightedBernoulliFailureBase}
\item \hyperlink{library-prerequisite-ffca3d1e44838990}{EconCSLib.Probability.weightedBernoulliSuccessBase}
\item \hyperlink{library-prerequisite-2d3d970688daf0d7}{EconCSLib.Probability.weightedBernoulliClosedRateBase}
\item \hyperlink{library-prerequisite-c3cc92a498459481}{EconCSLib.Probability.weightedBernoulliClosedThresholdRate}
\item \hyperlink{library-prerequisite-ecc26a3133fdd38b}{EconCSLib.finiteMax}
\item \hyperlink{library-prerequisite-087e3287a5c76915}{EconCSLib.finiteMin}
\item \hyperlink{library-prerequisite-dae3217c424b40b2}{EconCSLib.strictUpperPairSet}
\end{itemize}
\item \hyperlink{source-section-f10b5f68e4acf3dc}{Source claims (59)}
\begin{itemize}[leftmargin=1.2em]
\item \hyperlink{source-claim-d4f7ad79732f053c}{source\_quality\_domainSpec}
\item \hyperlink{source-claim-7979c322a62ef2c2}{source\_matching\_countSpec}
\item \hyperlink{source-claim-62613dce8c1c3874}{source\_matching\_functionSpec}
\item \hyperlink{source-claim-ad2abdefec264024}{source\_empirical\_reputation\_scoreSpec}
\item \hyperlink{source-claim-ab291bf70e6ff37c}{source\_system\_stateSpec}
\item \hyperlink{source-claim-b8da8ea96b79a0b8}{source\_definition\_pairwise\_accuracy\_eq1Spec}
\item \hyperlink{source-claim-88670d8e30834161}{source\_definition\_weighted\_objective\_eq2Spec}
\item \hyperlink{source-claim-55d3b1e0bb545f86}{source\_definition\_large\_deviation\_rateSpec}
\item \hyperlink{source-claim-9c0b6bf527ad55ca}{source\_definition\_step\_rule\_partition\_levelsSpec}
\item \hyperlink{source-claim-b2f4ad11d62b1527}{source\_definition\_lexicographic\_optimalitySpec}
\item \hyperlink{source-claim-84eec684fc1865ab}{theorem31\_source\_matching\_function\_unique\_value\_argmax\_lexicographicSpec}
\item \hyperlink{source-claim-8a7eb113492e3a82}{source\_theorem31\_adjacent\_rate\_eq3Spec}
\item \hyperlink{source-claim-50a2d9627ea4a9d0}{theorem31\_source\_cell\_matching\_rate\_eq\_lower\_cutpointSpec}
\item \hyperlink{source-claim-74286919e0e8cecc}{definition\_bernoulli\_kl\_formulaSpec}
\item \hyperlink{source-claim-d2fd3442bfa83879}{source\_lemma31\_equalization\_eq4Spec}
\item \hyperlink{source-claim-e679d8fd634fddb9}{lemma31\_closed\_adjacent\_rate\_formulaSpec}
\item \hyperlink{source-claim-71d943c2f56bac62}{theorem32\_weighted\_nested\_bisection\_outputSpec}
\item \hyperlink{source-claim-9c02a19f4b41b0c2}{theorem32\_weighted\_nested\_bisection\_loss\_and\_runtimeSpec}
\item \hyperlink{source-claim-7159be549a8d27ef}{section4\_question\_distributionSpec}
\item \hyperlink{source-claim-a5f9ed46c3edb5c1}{section4\_induced\_binary\_responseSpec}
\item \hyperlink{source-claim-3fed172202fb05fc}{section4\_l1\_design\_objectiveSpec}
\item \hyperlink{source-claim-642a63ff09a11d8b}{section4\_question\_design\_solutionSpec}
\item \hyperlink{source-claim-ffa7fbb9f1c16e4b}{appendixB1\_active\_setSpec}
\item \hyperlink{source-claim-9053ffafcfb8db42}{appendixB1\_transition\_kernelSpec}
\item \hyperlink{source-claim-6724e8b8f279b968}{appendixB1\_state\_updateSpec}
\item \hyperlink{source-claim-52ebae870dd67c82}{lemmaB1\_matching\_rate\_shiftSpec}
\item \hyperlink{source-claim-f84fba16d2553222}{source\_theoremB1\_quantile\_representationSpec}
\item \hyperlink{source-claim-5249ffde05e3b08f}{paper\_theoremB1\_uniform\_subsequence\_principle\_to\_of\_quantile\_floor\_tendstoUniformlyOn\_geometric\_meshSpec}
\item \hyperlink{source-claim-c34337b3884bb5d6}{appendixB5\_known\_type\_experimentSpec}
\item \hyperlink{source-claim-b26a46506065f742}{appendixB5\_random\_question\_experimentSpec}
\item \hyperlink{source-claim-8ccc8a8f737a5e7c}{appendixB5\_empirical\_question\_responseSpec}
\item \hyperlink{source-claim-276fd5feab006dec}{paper\_lemmaB2\_knownTypeExperiment\_uniform\_convergence\_of\_lipschitz\_trackingSpec}
\item \hyperlink{source-claim-76b42fc11ca5198e}{lemmaB2\_knownTypeExperiment\_random\_question\_sllnSpec}
\item \hyperlink{source-claim-e38b113829a05c8d}{appendixB5\_unknown\_type\_experimentSpec}
\item \hyperlink{source-claim-045f11cd1b489b98}{paper\_lemmaB3\_unknownTypeExperiment\_uniform\_convergence\_of\_ranking\_trackingSpec}
\item \hyperlink{source-claim-99b64aa08841b772}{lemmaB3\_unknownTypeExperiment\_random\_question\_response\_and\_rankSpec}
\item \hyperlink{source-claim-7ea3c8ab498f92dc}{source\_lemmaC1\_pairwise\_error\_rate\_from\_derivativesSpec}
\item \hyperlink{source-claim-5122f84bd17b3a2c}{source\_lemmaC2\_binary\_complement\_rateSpec}
\item \hyperlink{source-claim-d95b8ab740a8ceda}{source\_theoremC1\_laplace\_principleSpec}
\item \hyperlink{source-claim-16ee341855124303}{source\_remarkC1\_full\_square\_application\_under\_explicit\_conditionsSpec}
\item \hyperlink{source-claim-7529d1dedcffdbe3}{paper\_lemmaC3\_partition\_integral\_hasExponentialRate\_of\_min\_component\_uniform\_logRateSpec}
\item \hyperlink{source-claim-043c064a62132f60}{paper\_lemmaC4\_finite\_selected\_pullback\_positive\_rate\_and\_nonfiniteStep\_source\_zero\_rateSpec}
\item \hyperlink{source-claim-15e577a1127ae1d6}{source\_remarkC2\_weighted\_kl\_separation\_monotonicitySpec}
\item \hyperlink{source-claim-f9871b4133903309}{source\_appendixC5\_rate\_notation\_eq23Spec}
\item \hyperlink{source-claim-f8aedb48228e5085}{source\_lemmaC5\_uniform\_doubled\_chainSpec}
\item \hyperlink{source-claim-428f19cd38ba7b0a}{source\_lemmaC5\_refinement\_equations24\_25Spec}
\item \hyperlink{source-claim-4cea7b91f971b40f}{paper\_corollaryC2\_uniform\_equalized\_last\_rate\_tendsto\_zeroSpec}
\item \hyperlink{source-claim-510e870d600e0e8e}{paper\_corollaryC2\_uniform\_equalized\_adjacent\_mesh\_tendsto\_zeroSpec}
\item \hyperlink{source-claim-6e85310ac4deebe8}{lemmaC6\_monotone\_matching\_penultimate\_level\_boundSpec}
\item \hyperlink{source-claim-2d8e50cd70ac2429}{lemmaC7\_uniform\_doubled\_objective\_rate\_ge\_one\_fifth\_old\_objectiveSpec}
\item \hyperlink{source-claim-905297a638ebd8e0}{source\_lemmaC8\_uniform\_first\_level\_polynomial\_lower\_boundSpec}
\item \hyperlink{source-claim-d7f37d83d608e4fa}{corollaryC3\_monotone\_scaled\_first\_level\_ge\_half\_inv\_adjacent\_count\_sqSpec}
\item \hyperlink{source-claim-8cc85076f8ab0b2b}{source\_lemmaC9\_nested\_bisection\_runtime\_log\_squaredSpec}
\item \hyperlink{source-claim-51c42bc84b84361b}{source\_theoremB1\_proof\_selector\_nestingSpec}
\item \hyperlink{source-claim-b7febbacfb3801fa}{source\_corollaryC4\_kendall\_spearman\_subsequenceSpec}
\item \hyperlink{source-claim-196d59db7138fcbd}{definitionC1\_kendall\_spearman\_population\_objectivesSpec}
\item \hyperlink{source-claim-3b01d7ad0b5d48f9}{paper\_lemmaC10\_spearman\_linear\_weight\_ordered\_pair\_interval\_objective\_eq\_gap\_objectiveSpec}
\item \hyperlink{source-claim-75f334c9ed5f4c67}{paper\_lemmaC11\_kendall\_constant\_weight\_ordered\_pair\_interval\_objective\_le\_equispacedSpec}
\item \hyperlink{source-claim-9e5abbd5563b8516}{paper\_lemmaC12\_spearman\_linear\_weight\_ordered\_pair\_interval\_objective\_le\_equispacedSpec}
\end{itemize}
\end{itemize}

\clearpage
\hypertarget{library-section-a85f9574c44a156c}{}
\section*{Material library prerequisites}
\hypertarget{library-prerequisite-29ba02126b109d91}{}
\subsection*{\small\ttfamily\raggedright EconCSLib.\allowbreak{}Optimization.\allowbreak{}IsLexicographicMaximizerOn}
\textbf{Source connection:} no explicit byte-pinned paper source connection is registered.
\paragraph{Lean-expanded library semantic target}
\begin{ReviewVerbatim}
∀ {α : Type u_1} (feasible : α → Prop) (primary secondary : α → ℝ) (x : α),
  feasible x ∧ ∀ (y : α), feasible y → primary y < primary x ∨ primary y = primary x ∧ secondary y ≤ secondary x
\end{ReviewVerbatim}

\paragraph{Exact Lean library declaration}
\begin{ReviewVerbatim}
def IsLexicographicMaximizerOn {α : Type*} (feasible : α → Prop)
    (primary secondary : α → ℝ) (x : α) : Prop :=
  feasible x ∧ ∀ y, feasible y →
    primary y < primary x ∨
      (primary y = primary x ∧ secondary y ≤ secondary x)
\end{ReviewVerbatim}

\paragraph{Recorded source-to-library screening}
\textbf{Verdict:} \texttt{not recorded}
\reviewmatch{Matches source input}
\noindent\textbf{Reviewer annotation}\par
\reviewerbox
\clearpage
\hypertarget{library-prerequisite-0a2e53c8e507cb74}{}
\subsection*{\small\ttfamily\raggedright EconCSLib.\allowbreak{}Optimization.\allowbreak{}IsMaximizerOn}
\textbf{Source connection:} no explicit byte-pinned paper source connection is registered.
\paragraph{Lean-expanded library semantic target}
\begin{ReviewVerbatim}
∀ {α : Type u_1} (feasible : α → Prop) (objective : α → ℝ) (x : α),
  feasible x ∧ ∀ (y : α), feasible y → objective y ≤ objective x
\end{ReviewVerbatim}

\paragraph{Exact Lean library declaration}
\begin{ReviewVerbatim}
def IsMaximizerOn {α : Type*} (feasible : α → Prop)
    (objective : α → ℝ) (x : α) : Prop :=
  feasible x ∧ ∀ y, feasible y → objective y ≤ objective x
\end{ReviewVerbatim}

\paragraph{Recorded source-to-library screening}
\textbf{Verdict:} \texttt{not recorded}
\reviewmatch{Matches source input}
\noindent\textbf{Reviewer annotation}\par
\reviewerbox
\clearpage
\hypertarget{library-prerequisite-9e10d8e3c1b32191}{}
\subsection*{\small\ttfamily\raggedright EconCSLib.\allowbreak{}Optimization.\allowbreak{}IsMinimizerOn}
\textbf{Source connection:} no explicit byte-pinned paper source connection is registered.
\paragraph{Lean-expanded library semantic target}
\begin{ReviewVerbatim}
∀ {α : Type u_1} (feasible : α → Prop) (objective : α → ℝ) (x : α),
  feasible x ∧ ∀ (y : α), feasible y → objective x ≤ objective y
\end{ReviewVerbatim}

\paragraph{Exact Lean library declaration}
\begin{ReviewVerbatim}
def IsMinimizerOn {α : Type*} (feasible : α → Prop)
    (objective : α → ℝ) (x : α) : Prop :=
  feasible x ∧ ∀ y, feasible y → objective x ≤ objective y
\end{ReviewVerbatim}

\paragraph{Recorded source-to-library screening}
\textbf{Verdict:} \texttt{not recorded}
\reviewmatch{Matches source input}
\noindent\textbf{Reviewer annotation}\par
\reviewerbox
\clearpage
\hypertarget{library-prerequisite-8d013dcbf8315291}{}
\subsection*{\small\ttfamily\raggedright EconCSLib.\allowbreak{}Optimization.\allowbreak{}nestedBisectionStepBound}
\textbf{Source connection:} no explicit byte-pinned paper source connection is registered.
\paragraph{Lean-expanded library semantic target}
\begin{ReviewVerbatim}
(M L : ℕ) → (L + 1) * (1 + (M - 3) * L)
\end{ReviewVerbatim}

\paragraph{Exact Lean library declaration}
\begin{ReviewVerbatim}
def nestedBisectionStepBound (M L : ℕ) : ℕ :=
  (L + 1) * (1 + (M - 3) * L)
\end{ReviewVerbatim}

\paragraph{Recorded source-to-library screening}
\textbf{Verdict:} \texttt{not recorded}
\reviewmatch{Matches source input}
\noindent\textbf{Reviewer annotation}\par
\reviewerbox
\clearpage
\hypertarget{library-prerequisite-33a658a76ab86d8b}{}
\subsection*{\small\ttfamily\raggedright EconCSLib.\allowbreak{}Optimization.\allowbreak{}realBisectionMidpoint}
\textbf{Source connection:} no explicit byte-pinned paper source connection is registered.
\paragraph{Lean-expanded library semantic target}
\begin{ReviewVerbatim}
(lower upper : ℝ) → (lower + upper) / 2
\end{ReviewVerbatim}

\paragraph{Exact Lean library declaration}
\begin{ReviewVerbatim}
noncomputable def realBisectionMidpoint (lower upper : ℝ) : ℝ :=
  (lower + upper) / 2
\end{ReviewVerbatim}

\paragraph{Recorded source-to-library screening}
\textbf{Verdict:} \texttt{not recorded}
\reviewmatch{Matches source input}
\noindent\textbf{Reviewer annotation}\par
\reviewerbox
\clearpage
\hypertarget{library-prerequisite-f7e825d2da547274}{}
\subsection*{\small\ttfamily\raggedright EconCSLib.\allowbreak{}Optimization.\allowbreak{}realBisectionStep}
\textbf{Source connection:} no explicit byte-pinned paper source connection is registered.
\paragraph{Lean-expanded library semantic target}
\begin{ReviewVerbatim}
(above : ℝ → Bool) →
  (lower upper : ℝ) →
    have mid := EconCSLib.Optimization.realBisectionMidpoint lower upper;
    if above mid = true then (lower, mid) else (mid, upper)
\end{ReviewVerbatim}

\paragraph{Exact Lean library declaration}
\begin{ReviewVerbatim}
noncomputable def realBisectionStep
    (above : ℝ → Bool) (lower upper : ℝ) : ℝ × ℝ :=
  let mid := realBisectionMidpoint lower upper
  if above mid then (lower, mid) else (mid, upper)
\end{ReviewVerbatim}

\paragraph{Recorded source-to-library screening}
\textbf{Verdict:} \texttt{not recorded}
\reviewmatch{Matches source input}
\noindent\textbf{Reviewer annotation}\par
\reviewerbox
\clearpage
\hypertarget{library-prerequisite-e24bfb229e6e0f0c}{}
\subsection*{\small\ttfamily\raggedright EconCSLib.\allowbreak{}Optimization.\allowbreak{}realBisectionStepFn}
\textbf{Source connection:} no explicit byte-pinned paper source connection is registered.
\paragraph{Lean-expanded library semantic target}
\begin{ReviewVerbatim}
(above : ℝ → Bool) → (p : ℝ × ℝ) → EconCSLib.Optimization.realBisectionStep above p.1 p.2
\end{ReviewVerbatim}

\paragraph{Exact Lean library declaration}
\begin{ReviewVerbatim}
noncomputable def realBisectionStepFn
    (above : ℝ → Bool) (p : ℝ × ℝ) : ℝ × ℝ :=
  realBisectionStep above p.1 p.2
\end{ReviewVerbatim}

\paragraph{Recorded source-to-library screening}
\textbf{Verdict:} \texttt{not recorded}
\reviewmatch{Matches source input}
\noindent\textbf{Reviewer annotation}\par
\reviewerbox
\clearpage
\hypertarget{library-prerequisite-5b79c0723b67e9a9}{}
\subsection*{\small\ttfamily\raggedright EconCSLib.\allowbreak{}Optimization.\allowbreak{}realBisectionRun}
\textbf{Source connection:} no explicit byte-pinned paper source connection is registered.
\paragraph{Lean-expanded library semantic target}
\begin{ReviewVerbatim}
(above : ℝ → Bool) → (n : ℕ) → (lower upper : ℝ) → (EconCSLib.Optimization.realBisectionStepFn above)^[n] (lower, upper)
\end{ReviewVerbatim}

\paragraph{Exact Lean library declaration}
\begin{ReviewVerbatim}
noncomputable def realBisectionRun
    (above : ℝ → Bool) (n : ℕ) (lower upper : ℝ) : ℝ × ℝ :=
  Nat.iterate (realBisectionStepFn above) n (lower, upper)
\end{ReviewVerbatim}

\paragraph{Recorded source-to-library screening}
\textbf{Verdict:} \texttt{not recorded}
\reviewmatch{Matches source input}
\noindent\textbf{Reviewer annotation}\par
\reviewerbox
\clearpage
\hypertarget{library-prerequisite-e4ddb15be8735d63}{}
\subsection*{\small\ttfamily\raggedright EconCSLib.\allowbreak{}Probability.\allowbreak{}ExponentialRateCertificate}
\textbf{Source connection:} no explicit byte-pinned paper source connection is registered.
\paragraph{Lean declaration type (metadata)}
\begin{ReviewVerbatim}
(ℕ → ℝ) → ℝ → Prop
\end{ReviewVerbatim}

\paragraph{Exact Lean library declaration}
\begin{ReviewVerbatim}
structure ExponentialRateCertificate (p : ℕ → ℝ) (rate : ℝ) where
  eventually_pos : ∀ᶠ n in atTop, 0 < p n
  has_rate : HasExponentialRate p rate
\end{ReviewVerbatim}

\paragraph{Recorded source-to-library screening}
\textbf{Verdict:} \texttt{not recorded}
\reviewmatch{Matches source input}
\noindent\textbf{Reviewer annotation}\par
\reviewerbox
\clearpage
\hypertarget{library-prerequisite-8386068971d3a1c1}{}
\subsection*{\small\ttfamily\raggedright EconCSLib.\allowbreak{}Probability.\allowbreak{}FiniteMeasurableSetPartition}
\textbf{Source connection:} no explicit byte-pinned paper source connection is registered.
\paragraph{Lean declaration type (metadata)}
\begin{ReviewVerbatim}
{α : Type u_1} →
  [inst : MeasurableSpace α] → MeasureTheory.Measure α → (Piece : Type u_2) → [Fintype Piece] → Type (max u_1 u_2)
\end{ReviewVerbatim}

\paragraph{Exact Lean library declaration}
\begin{ReviewVerbatim}
structure FiniteMeasurableSetPartition
    {α : Type*} [MeasurableSpace α]
    (μ : Measure α) (Piece : Type*) [Fintype Piece] where
  support : Set α
  pieceSet : Piece → Set α
  measurable_piece : ∀ piece, MeasurableSet (pieceSet piece)
  pairwiseDisjoint :
    Pairwise (fun piece₁ piece₂ => Disjoint (pieceSet piece₁) (pieceSet piece₂))
  cover : Set.iUnion pieceSet = support
\end{ReviewVerbatim}

\paragraph{Recorded source-to-library screening}
\textbf{Verdict:} \texttt{not recorded}
\reviewmatch{Matches source input}
\noindent\textbf{Reviewer annotation}\par
\reviewerbox
\clearpage
\hypertarget{library-prerequisite-94654614a042c92d}{}
\subsection*{\small\ttfamily\raggedright EconCSLib.\allowbreak{}Probability.\allowbreak{}FiniteMeasurableSetPartition.\allowbreak{}mk}
\textbf{Source connection:} no explicit byte-pinned paper source connection is registered.
\paragraph{Lean declaration type (metadata)}
\begin{ReviewVerbatim}
{α : Type u_1} →
  [inst : MeasurableSpace α] →
    {μ : MeasureTheory.Measure α} →
      {Piece : Type u_2} →
        [inst_1 : Fintype Piece] →
          (support : Set α) →
            (pieceSet : Piece → Set α) →
              (∀ (piece : Piece), MeasurableSet (pieceSet piece)) →
                (Pairwise fun piece₁ piece₂ => Disjoint (pieceSet piece₁) (pieceSet piece₂)) →
                  Set.iUnion pieceSet = support → EconCSLib.Probability.FiniteMeasurableSetPartition μ Piece
\end{ReviewVerbatim}

\paragraph{Exact Lean library declaration}
\begin{ReviewVerbatim}
library declaration is not registered or discoverable
\end{ReviewVerbatim}

\paragraph{Recorded source-to-library screening}
\textbf{Verdict:} \texttt{not recorded}
\reviewmatch{Matches source input}
\noindent\textbf{Reviewer annotation}\par
\reviewerbox
\clearpage
\hypertarget{library-prerequisite-6f62f2ee03548cd2}{}
\subsection*{\small\ttfamily\raggedright EconCSLib.\allowbreak{}Probability.\allowbreak{}FiniteMeasurableSetPartition.\allowbreak{}orderedRealIocNatCutpoints}
\textbf{Source connection:} no explicit byte-pinned paper source connection is registered.
\paragraph{Lean-expanded library semantic target}
\begin{ReviewVerbatim}
(μ : MeasureTheory.Measure ℝ) →
  (n : ℕ) →
    (cut : ℕ → ℝ) →
      (hmono : Monotone cut) →
        { support := Set.Ioc (cut 0) (cut n), pieceSet := fun piece => Set.Ioc (cut ↑piece) (cut (↑piece + 1)),
          measurable_piece := ⋯, pairwiseDisjoint := ⋯, cover := ⋯ }
\end{ReviewVerbatim}

\paragraph{Exact Lean library declaration}
\begin{ReviewVerbatim}
def orderedRealIocNatCutpoints
    (μ : Measure ℝ) (n : ℕ) (cut : ℕ → ℝ) (hmono : Monotone cut) :
    FiniteMeasurableSetPartition μ (Fin n) where
  support := Set.Ioc (cut 0) (cut n)
  pieceSet := fun piece => Set.Ioc (cut piece.val) (cut (piece.val + 1))
  measurable_piece := by
    intro _piece
    exact measurableSet_Ioc
  pairwiseDisjoint := by
    intro piece₁ piece₂ hne
    by_cases hlt : piece₁.val < piece₂.val
    · exact Set.Ioc_disjoint_Ioc_of_le
        (hmono (Nat.succ_le_of_lt hlt))
    · have hval_ne : piece₁.val ≠ piece₂.val := by
        intro hval
        exact hne (Fin.ext hval)
      have hgt : piece₂.val < piece₁.val := by
        exact lt_of_le_of_ne (Nat.le_of_not_gt hlt) hval_ne.symm
      exact (Set.Ioc_disjoint_Ioc_of_le
        (hmono (Nat.succ_le_of_lt hgt))).symm
  cover := by
    ext x
    constructor
    · intro hx
      rcases Set.mem_iUnion.mp hx with ⟨piece, hxpiece⟩
      exact
        ⟨lt_of_le_of_lt (hmono (Nat.zero_le piece.val)) hxpiece.1,
          le_trans hxpiece.2 (hmono (Nat.succ_le_of_lt piece.isLt))⟩
    · intro hx
      have hxUnion :
          x ∈ ⋃ i ∈ Finset.range n,
            Set.Ioc (cut i) (cut (i + 1)) :=
        Ioc_subset_biUnion_Ioc n cut hx
      rcases Set.mem_iUnion.mp hxUnion with ⟨i, hxi⟩
      rcases Set.mem_iUnion.mp hxi with ⟨hirange, hxpiece⟩
      have hi : i < n := by
        simpa using hirange
      exact Set.mem_iUnion.mpr ⟨⟨i, hi⟩, by simpa using hxpiece⟩
\end{ReviewVerbatim}

\paragraph{Recorded source-to-library screening}
\textbf{Verdict:} \texttt{not recorded}
\reviewmatch{Matches source input}
\noindent\textbf{Reviewer annotation}\par
\reviewerbox
\clearpage
\hypertarget{library-prerequisite-aab7fae9d893ec88}{}
\subsection*{\small\ttfamily\raggedright EconCSLib.\allowbreak{}Probability.\allowbreak{}FiniteMeasurableSetPartition.\allowbreak{}orderedRealIocNatCutpointsMidpoint}
\textbf{Source connection:} no explicit byte-pinned paper source connection is registered.
\paragraph{Lean-expanded library semantic target}
\begin{ReviewVerbatim}
(cut : ℕ → ℝ) → {n : ℕ} → (piece : Fin n) → (cut ↑piece + cut (↑piece + 1)) / 2
\end{ReviewVerbatim}

\paragraph{Exact Lean library declaration}
\begin{ReviewVerbatim}
def orderedRealIocNatCutpointsMidpoint
    (cut : ℕ → ℝ) {n : ℕ} (piece : Fin n) : ℝ :=
  (cut piece.val + cut (piece.val + 1)) / 2
\end{ReviewVerbatim}

\paragraph{Recorded source-to-library screening}
\textbf{Verdict:} \texttt{not recorded}
\reviewmatch{Matches source input}
\noindent\textbf{Reviewer annotation}\par
\reviewerbox
\clearpage
\hypertarget{library-prerequisite-0b261505488b443f}{}
\subsection*{\small\ttfamily\raggedright EconCSLib.\allowbreak{}Probability.\allowbreak{}FiniteMeasurableSetPartition.\allowbreak{}orderedRealIocNatCutpointsProdMidpoint}
\textbf{Source connection:} no explicit byte-pinned paper source connection is registered.
\paragraph{Lean-expanded library semantic target}
\begin{ReviewVerbatim}
(cut₁ cut₂ : ℕ → ℝ) →
  {n m : ℕ} →
    (piece : Fin n × Fin m) →
      (EconCSLib.Probability.FiniteMeasurableSetPartition.orderedRealIocNatCutpointsMidpoint cut₁ piece.1,
        EconCSLib.Probability.FiniteMeasurableSetPartition.orderedRealIocNatCutpointsMidpoint cut₂ piece.2)
\end{ReviewVerbatim}

\paragraph{Exact Lean library declaration}
\begin{ReviewVerbatim}
def orderedRealIocNatCutpointsProdMidpoint
    (cut₁ cut₂ : ℕ → ℝ) {n m : ℕ} (piece : Fin n × Fin m) :
    ℝ × ℝ :=
  (orderedRealIocNatCutpointsMidpoint cut₁ piece.1,
    orderedRealIocNatCutpointsMidpoint cut₂ piece.2)
\end{ReviewVerbatim}

\paragraph{Recorded source-to-library screening}
\textbf{Verdict:} \texttt{not recorded}
\reviewmatch{Matches source input}
\noindent\textbf{Reviewer annotation}\par
\reviewerbox
\clearpage
\hypertarget{library-prerequisite-5833c2529d9caeeb}{}
\subsection*{\small\ttfamily\raggedright EconCSLib.\allowbreak{}Probability.\allowbreak{}FiniteMeasurableSetPartition.\allowbreak{}pieceSet}
\textbf{Source connection:} no explicit byte-pinned paper source connection is registered.
\paragraph{Lean-expanded library semantic target}
\begin{ReviewVerbatim}
∀ {α : Type u_1} [inst : MeasurableSpace α] {μ : MeasureTheory.Measure α} {Piece : Type u_2} [inst_1 : Fintype Piece]
  (self : EconCSLib.Probability.FiniteMeasurableSetPartition μ Piece) (a : Piece) (a_1 : α), self.2 a a_1
\end{ReviewVerbatim}

\paragraph{Exact Lean library declaration}
\begin{ReviewVerbatim}
library declaration is not registered or discoverable
\end{ReviewVerbatim}

\paragraph{Recorded source-to-library screening}
\textbf{Verdict:} \texttt{not recorded}
\reviewmatch{Matches source input}
\noindent\textbf{Reviewer annotation}\par
\reviewerbox
\clearpage
\hypertarget{library-prerequisite-a2c998741ee4703f}{}
\subsection*{\small\ttfamily\raggedright EconCSLib.\allowbreak{}Probability.\allowbreak{}FiniteMeasurableSetPartition.\allowbreak{}piecewiseConstKernel}
\textbf{Source connection:} no explicit byte-pinned paper source connection is registered.
\paragraph{Lean-expanded library semantic target}
\begin{ReviewVerbatim}
{α : Type u_1} →
  {Piece : Type u_2} →
    [inst : MeasurableSpace α] →
      [inst_1 : Fintype Piece] →
        {μ : MeasureTheory.Measure α} →
          (P : EconCSLib.Probability.FiniteMeasurableSetPartition μ Piece) →
            (p : Piece → ℕ → ℝ) → (n : ℕ) → (x : α) → ∑ piece, (P.pieceSet piece).indicator (fun x => p piece n) x
\end{ReviewVerbatim}

\paragraph{Exact Lean library declaration}
\begin{ReviewVerbatim}
def piecewiseConstKernel
    (P : FiniteMeasurableSetPartition μ Piece)
    (p : Piece → ℕ → ℝ) (n : ℕ) (x : α) : ℝ :=
  ∑ piece : Piece, (P.pieceSet piece).indicator (fun _ : α => p piece n) x
\end{ReviewVerbatim}

\paragraph{Recorded source-to-library screening}
\textbf{Verdict:} \texttt{not recorded}
\reviewmatch{Matches source input}
\noindent\textbf{Reviewer annotation}\par
\reviewerbox
\clearpage
\hypertarget{library-prerequisite-a9f78bbf3e21bb13}{}
\subsection*{\small\ttfamily\raggedright EconCSLib.\allowbreak{}Probability.\allowbreak{}FiniteMeasurableSetPartition.\allowbreak{}support}
\textbf{Source connection:} no explicit byte-pinned paper source connection is registered.
\paragraph{Lean-expanded library semantic target}
\begin{ReviewVerbatim}
∀ {α : Type u_1} [inst : MeasurableSpace α] {μ : MeasureTheory.Measure α} {Piece : Type u_2} [inst_1 : Fintype Piece]
  (self : EconCSLib.Probability.FiniteMeasurableSetPartition μ Piece) (a : α), self.1 a
\end{ReviewVerbatim}

\paragraph{Exact Lean library declaration}
\begin{ReviewVerbatim}
library declaration is not registered or discoverable
\end{ReviewVerbatim}

\paragraph{Recorded source-to-library screening}
\textbf{Verdict:} \texttt{not recorded}
\reviewmatch{Matches source input}
\noindent\textbf{Reviewer annotation}\par
\reviewerbox
\clearpage
\hypertarget{library-prerequisite-3161aa1be2e11c39}{}
\subsection*{\small\ttfamily\raggedright EconCSLib.\allowbreak{}Probability.\allowbreak{}FiniteMeasurableSetPartition.\allowbreak{}prod}
\textbf{Source connection:} no explicit byte-pinned paper source connection is registered.
\paragraph{Lean-expanded library semantic target}
\begin{ReviewVerbatim}
{α : Type u_1} →
  {Piece : Type u_2} →
    [inst : MeasurableSpace α] →
      [inst_1 : Fintype Piece] →
        {μ : MeasureTheory.Measure α} →
          {β : Type u_4} →
            {Pieceβ : Type u_5} →
              [inst_2 : MeasurableSpace β] →
                [inst_3 : Fintype Pieceβ] →
                  {ν : MeasureTheory.Measure β} →
                    (P : EconCSLib.Probability.FiniteMeasurableSetPartition μ Piece) →
                      (Q : EconCSLib.Probability.FiniteMeasurableSetPartition ν Pieceβ) →
                        { support := P.support ×ˢ Q.support,
                          pieceSet := fun piece => P.pieceSet piece.1 ×ˢ Q.pieceSet piece.2, measurable_piece := ⋯,
                          pairwiseDisjoint := ⋯, cover := ⋯ }
\end{ReviewVerbatim}

\paragraph{Exact Lean library declaration}
\begin{ReviewVerbatim}
def prod {β Pieceβ : Type*} [MeasurableSpace β] [Fintype Pieceβ]
    {ν : Measure β}
    (P : FiniteMeasurableSetPartition μ Piece)
    (Q : FiniteMeasurableSetPartition ν Pieceβ) :
    FiniteMeasurableSetPartition (μ.prod ν) (Piece × Pieceβ) where
  support := P.support ×ˢ Q.support
  pieceSet := fun piece => P.pieceSet piece.1 ×ˢ Q.pieceSet piece.2
  measurable_piece := by
    intro piece
    exact (P.measurable_piece piece.1).prod (Q.measurable_piece piece.2)
  pairwiseDisjoint := by
    intro piece₁ piece₂ hne
    rw [Set.disjoint_left]
    intro x hx₁ hx₂
    rcases x with ⟨x₁, x₂⟩
    by_cases hfirst : piece₁.1 = piece₂.1
    · have hsecond : piece₁.2 ≠ piece₂.2 := by
        intro hsecond
        exact hne (Prod.ext hfirst hsecond)
      exact
        (Set.disjoint_left.mp
          (Q.pairwiseDisjoint hsecond)) hx₁.2 hx₂.2
    · exact
        (Set.disjoint_left.mp
          (P.pairwiseDisjoint hfirst)) hx₁.1 hx₂.1
  cover := by
    ext x
    rcases x with ⟨x₁, x₂⟩
    constructor
    · intro hx
      rcases Set.mem_iUnion.mp hx with ⟨piece, hxpiece⟩
      have hx₁ : x₁ ∈ Set.iUnion P.pieceSet :=
        Set.mem_iUnion.mpr ⟨piece.1, hxpiece.1⟩
      have hx₂ : x₂ ∈ Set.iUnion Q.pieceSet :=
        Set.mem_iUnion.mpr ⟨piece.2, hxpiece.2⟩
      exact ⟨by simpa [P.cover] using hx₁, by simpa [Q.cover] using hx₂⟩
    · intro hx
      have hx₁ : x₁ ∈ Set.iUnion P.pieceSet := by
        simpa [P.cover] using hx.1
      have hx₂ : x₂ ∈ Set.iUnion Q.pieceSet := by
        simpa [Q.cover] using hx.2
      rcases Set.mem_iUnion.mp hx₁ with ⟨piece₁, hxpiece₁⟩
      rcases Set.mem_iUnion.mp hx₂ with ⟨piece₂, hxpiece₂⟩
      exact Set.mem_iUnion.mpr ⟨(piece₁, piece₂), ⟨hxpiece₁, hxpiece₂⟩⟩
\end{ReviewVerbatim}

\paragraph{Recorded source-to-library screening}
\textbf{Verdict:} \texttt{not recorded}
\reviewmatch{Matches source input}
\noindent\textbf{Reviewer annotation}\par
\reviewerbox
\clearpage
\hypertarget{library-prerequisite-e99ffb1d85bad07e}{}
\subsection*{\small\ttfamily\raggedright EconCSLib.\allowbreak{}Probability.\allowbreak{}FiniteMeasurableSetPartition.\allowbreak{}subpartition}
\textbf{Source connection:} no explicit byte-pinned paper source connection is registered.
\paragraph{Lean-expanded library semantic target}
\begin{ReviewVerbatim}
{α : Type u_1} →
  {Piece : Type u_2} →
    [inst : MeasurableSpace α] →
      [inst_1 : Fintype Piece] →
        {μ : MeasureTheory.Measure α} →
          (P : EconCSLib.Probability.FiniteMeasurableSetPartition μ Piece) →
            (selected : Piece → Prop) →
              [inst_2 : DecidablePred selected] →
                { support := ⋃ piece, P.pieceSet ↑piece, pieceSet := fun piece => P.pieceSet ↑piece,
                  measurable_piece := ⋯, pairwiseDisjoint := ⋯, cover := ⋯ }
\end{ReviewVerbatim}

\paragraph{Exact Lean library declaration}
\begin{ReviewVerbatim}
def subpartition
    (P : FiniteMeasurableSetPartition μ Piece)
    (selected : Piece → Prop) [DecidablePred selected] :
    FiniteMeasurableSetPartition μ {piece : Piece // selected piece} where
  support := Set.iUnion (fun piece : {piece : Piece // selected piece} =>
    P.pieceSet piece.val)
  pieceSet := fun piece => P.pieceSet piece.val
  measurable_piece := by
    intro piece
    exact P.measurable_piece piece.val
  pairwiseDisjoint := by
    intro piece₁ piece₂ hne
    have hval : piece₁.val ≠ piece₂.val := by
      intro hval
      exact hne (Subtype.ext hval)
    exact P.pairwiseDisjoint hval
  cover := rfl
\end{ReviewVerbatim}

\paragraph{Recorded source-to-library screening}
\textbf{Verdict:} \texttt{not recorded}
\reviewmatch{Matches source input}
\noindent\textbf{Reviewer annotation}\par
\reviewerbox
\clearpage
\hypertarget{library-prerequisite-cb9e54451ad9288c}{}
\subsection*{\small\ttfamily\raggedright EconCSLib.\allowbreak{}Probability.\allowbreak{}FiniteMeasurableSetPartition.\allowbreak{}selectedProduct}
\textbf{Source connection:} no explicit byte-pinned paper source connection is registered.
\paragraph{Lean-expanded library semantic target}
\begin{ReviewVerbatim}
{α : Type u_1} →
  {Piece : Type u_2} →
    [inst : MeasurableSpace α] →
      [inst_1 : Fintype Piece] →
        {μ : MeasureTheory.Measure α} →
          {β : Type u_4} →
            {Pieceβ : Type u_5} →
              [inst_2 : MeasurableSpace β] →
                [inst_3 : Fintype Pieceβ] →
                  {ν : MeasureTheory.Measure β} →
                    (P : EconCSLib.Probability.FiniteMeasurableSetPartition μ Piece) →
                      (Q : EconCSLib.Probability.FiniteMeasurableSetPartition ν Pieceβ) →
                        (selected : Piece × Pieceβ → Prop) →
                          [inst_4 : DecidablePred selected] → (P.prod Q).subpartition selected
\end{ReviewVerbatim}

\paragraph{Exact Lean library declaration}
\begin{ReviewVerbatim}
def selectedProduct {β Pieceβ : Type*} [MeasurableSpace β] [Fintype Pieceβ]
    {ν : Measure β}
    (P : FiniteMeasurableSetPartition μ Piece)
    (Q : FiniteMeasurableSetPartition ν Pieceβ)
    (selected : Piece × Pieceβ → Prop) [DecidablePred selected] :
    FiniteMeasurableSetPartition (μ.prod ν)
      {piece : Piece × Pieceβ // selected piece} :=
  (P.prod Q).subpartition selected
\end{ReviewVerbatim}

\paragraph{Recorded source-to-library screening}
\textbf{Verdict:} \texttt{not recorded}
\reviewmatch{Matches source input}
\noindent\textbf{Reviewer annotation}\par
\reviewerbox
\clearpage
\hypertarget{library-prerequisite-943d19d91afe6d63}{}
\subsection*{\small\ttfamily\raggedright EconCSLib.\allowbreak{}Probability.\allowbreak{}FiniteRatingLDPModel}
\textbf{Source connection:} no explicit byte-pinned paper source connection is registered.
\paragraph{Lean declaration type (metadata)}
\begin{ReviewVerbatim}
Type u_2 → (Rating : Type u_3) → [Fintype Rating] → [DecidableEq Rating] → Type (max u_2 u_3)
\end{ReviewVerbatim}

\paragraph{Exact Lean library declaration}
\begin{ReviewVerbatim}
structure FiniteRatingLDPModel (θ Rating : Type*) [Fintype Rating]
    [DecidableEq Rating] where
  typeLaw : θ → PMF Rating
  score : Rating → ℝ
\end{ReviewVerbatim}

\paragraph{Recorded source-to-library screening}
\textbf{Verdict:} \texttt{not recorded}
\reviewmatch{Matches source input}
\noindent\textbf{Reviewer annotation}\par
\reviewerbox
\clearpage
\hypertarget{library-prerequisite-c3203b699a19e900}{}
\subsection*{\small\ttfamily\raggedright EconCSLib.\allowbreak{}Probability.\allowbreak{}FiniteRatingLDPModel.\allowbreak{}score}
\textbf{Source connection:} no explicit byte-pinned paper source connection is registered.
\paragraph{Lean-expanded library semantic target}
\begin{ReviewVerbatim}
{θ : Type u_2} →
  {Rating : Type u_3} →
    [inst : Fintype Rating] →
      [inst_1 : DecidableEq Rating] →
        (self : EconCSLib.Probability.FiniteRatingLDPModel θ Rating) → (a : Rating) → self.2 a
\end{ReviewVerbatim}

\paragraph{Exact Lean library declaration}
\begin{ReviewVerbatim}
library declaration is not registered or discoverable
\end{ReviewVerbatim}

\paragraph{Recorded source-to-library screening}
\textbf{Verdict:} \texttt{not recorded}
\reviewmatch{Matches source input}
\noindent\textbf{Reviewer annotation}\par
\reviewerbox
\clearpage
\hypertarget{library-prerequisite-fc314aabfaabee42}{}
\subsection*{\small\ttfamily\raggedright EconCSLib.\allowbreak{}Probability.\allowbreak{}FiniteRatingLDPModel.\allowbreak{}typeLaw}
\textbf{Source connection:} no explicit byte-pinned paper source connection is registered.
\paragraph{Lean-expanded library semantic target}
\begin{ReviewVerbatim}
{θ : Type u_2} →
  {Rating : Type u_3} →
    [inst : Fintype Rating] →
      [inst_1 : DecidableEq Rating] → (self : EconCSLib.Probability.FiniteRatingLDPModel θ Rating) → (a : θ) → self.1 a
\end{ReviewVerbatim}

\paragraph{Exact Lean library declaration}
\begin{ReviewVerbatim}
library declaration is not registered or discoverable
\end{ReviewVerbatim}

\paragraph{Recorded source-to-library screening}
\textbf{Verdict:} \texttt{not recorded}
\reviewmatch{Matches source input}
\noindent\textbf{Reviewer annotation}\par
\reviewerbox
\clearpage
\hypertarget{library-prerequisite-6258b3da52471ab5}{}
\subsection*{\small\ttfamily\raggedright EconCSLib.\allowbreak{}Probability.\allowbreak{}finiteMGF}
\textbf{Source connection:} no explicit byte-pinned paper source connection is registered.
\paragraph{Lean-expanded library semantic target}
\begin{ReviewVerbatim}
{α : Type u_1} →
  [inst : Fintype α] → (μ : PMF α) → (score : α → ℝ) → (z : ℝ) → ∑ a, (μ a).toReal * Real.exp (z * score a)
\end{ReviewVerbatim}

\paragraph{Exact Lean library declaration}
\begin{ReviewVerbatim}
def finiteMGF (μ : PMF α) (score : α → ℝ) (z : ℝ) : ℝ :=
  ∑ a : α, (μ a).toReal * Real.exp (z * score a)
\end{ReviewVerbatim}

\paragraph{Recorded source-to-library screening}
\textbf{Verdict:} \texttt{not recorded}
\reviewmatch{Matches source input}
\noindent\textbf{Reviewer annotation}\par
\reviewerbox
\clearpage
\hypertarget{library-prerequisite-20dbd208a81d9e80}{}
\subsection*{\small\ttfamily\raggedright EconCSLib.\allowbreak{}Probability.\allowbreak{}finiteLogMGF}
\textbf{Source connection:} no explicit byte-pinned paper source connection is registered.
\paragraph{Lean-expanded library semantic target}
\begin{ReviewVerbatim}
{α : Type u_1} →
  [inst : Fintype α] → (μ : PMF α) → (score : α → ℝ) → (z : ℝ) → Real.log (EconCSLib.Probability.finiteMGF μ score z)
\end{ReviewVerbatim}

\paragraph{Exact Lean library declaration}
\begin{ReviewVerbatim}
def finiteLogMGF (μ : PMF α) (score : α → ℝ) (z : ℝ) : ℝ :=
  Real.log (finiteMGF μ score z)
\end{ReviewVerbatim}

\paragraph{Recorded source-to-library screening}
\textbf{Verdict:} \texttt{not recorded}
\reviewmatch{Matches source input}
\noindent\textbf{Reviewer annotation}\par
\reviewerbox
\clearpage
\hypertarget{library-prerequisite-c959d6e7b586c2bc}{}
\subsection*{\small\ttfamily\raggedright EconCSLib.\allowbreak{}Probability.\allowbreak{}FiniteRatingLDPModel.\allowbreak{}logMGF}
\textbf{Source connection:} no explicit byte-pinned paper source connection is registered.
\paragraph{Lean-expanded library semantic target}
\begin{ReviewVerbatim}
{θ : Type u_2} →
  {Rating : Type u_3} →
    [inst : Fintype Rating] →
      [inst_1 : DecidableEq Rating] →
        (M : EconCSLib.Probability.FiniteRatingLDPModel θ Rating) →
          (t : θ) → (z : ℝ) → EconCSLib.Probability.finiteLogMGF (M.typeLaw t) M.score z
\end{ReviewVerbatim}

\paragraph{Exact Lean library declaration}
\begin{ReviewVerbatim}
def logMGF (M : FiniteRatingLDPModel θ Rating) (t : θ) (z : ℝ) : ℝ :=
  finiteLogMGF (M.typeLaw t) M.score z
\end{ReviewVerbatim}

\paragraph{Recorded source-to-library screening}
\textbf{Verdict:} \texttt{not recorded}
\reviewmatch{Matches source input}
\noindent\textbf{Reviewer annotation}\par
\reviewerbox
\clearpage
\hypertarget{library-prerequisite-bd97fc0d50344aa7}{}
\subsection*{\small\ttfamily\raggedright EconCSLib.\allowbreak{}Probability.\allowbreak{}FiniteRatingLDPModel.\allowbreak{}mk}
\textbf{Source connection:} no explicit byte-pinned paper source connection is registered.
\paragraph{Lean declaration type (metadata)}
\begin{ReviewVerbatim}
{θ : Type u_2} →
  {Rating : Type u_3} →
    [inst : Fintype Rating] →
      [inst_1 : DecidableEq Rating] →
        (θ → PMF Rating) → (Rating → ℝ) → EconCSLib.Probability.FiniteRatingLDPModel θ Rating
\end{ReviewVerbatim}

\paragraph{Exact Lean library declaration}
\begin{ReviewVerbatim}
library declaration is not registered or discoverable
\end{ReviewVerbatim}

\paragraph{Recorded source-to-library screening}
\textbf{Verdict:} \texttt{not recorded}
\reviewmatch{Matches source input}
\noindent\textbf{Reviewer annotation}\par
\reviewerbox
\clearpage
\hypertarget{library-prerequisite-d169ba913523821f}{}
\subsection*{\small\ttfamily\raggedright EconCSLib.\allowbreak{}Probability.\allowbreak{}HasAEEssentialInfimum}
\textbf{Source connection:} no explicit byte-pinned paper source connection is registered.
\paragraph{Lean-expanded library semantic target}
\begin{ReviewVerbatim}
∀ {α : Type u_1} [inst : MeasurableSpace α] (μ : MeasureTheory.Measure α) (phi : α → ℝ) (rate : ℝ),
  (∀ᵐ (x : α) ∂μ, rate ≤ phi x) ∧ ∀ ε > 0, ∃ s, MeasurableSet s ∧ 0 < μ.real s ∧ ∀ x ∈ s, phi x ≤ rate + ε
\end{ReviewVerbatim}

\paragraph{Exact Lean library declaration}
\begin{ReviewVerbatim}
def HasAEEssentialInfimum
    {α : Type*} [MeasurableSpace α]
    (μ : Measure α) (phi : α → ℝ) (rate : ℝ) : Prop :=
  (∀ᵐ x ∂μ, rate ≤ phi x) ∧
    ∀ ε > 0, ∃ s : Set α,
      MeasurableSet s ∧
        0 < μ.real s ∧
          ∀ x, x ∈ s → phi x ≤ rate + ε
\end{ReviewVerbatim}

\paragraph{Recorded source-to-library screening}
\textbf{Verdict:} \texttt{not recorded}
\reviewmatch{Matches source input}
\noindent\textbf{Reviewer annotation}\par
\reviewerbox
\clearpage
\hypertarget{library-prerequisite-94743506c5de049b}{}
\subsection*{\small\ttfamily\raggedright EconCSLib.\allowbreak{}Probability.\allowbreak{}logDecay}
\textbf{Source connection:} no explicit byte-pinned paper source connection is registered.
\paragraph{Lean-expanded library semantic target}
\begin{ReviewVerbatim}
(p : ℕ → ℝ) → (n : ℕ) → -(↑n)⁻¹ * Real.log (p n)
\end{ReviewVerbatim}

\paragraph{Exact Lean library declaration}
\begin{ReviewVerbatim}
def logDecay (p : ℕ → ℝ) (n : ℕ) : ℝ :=
  -((n : ℝ)⁻¹) * Real.log (p n)
\end{ReviewVerbatim}

\paragraph{Recorded source-to-library screening}
\textbf{Verdict:} \texttt{not recorded}
\reviewmatch{Matches source input}
\noindent\textbf{Reviewer annotation}\par
\reviewerbox
\clearpage
\hypertarget{library-prerequisite-380666720005662c}{}
\subsection*{\small\ttfamily\raggedright EconCSLib.\allowbreak{}Probability.\allowbreak{}HasExponentialRate}
\textbf{Source connection:} no explicit byte-pinned paper source connection is registered.
\paragraph{Lean-expanded library semantic target}
\begin{ReviewVerbatim}
∀ (p : ℕ → ℝ) (rate : ℝ), Filter.Tendsto (EconCSLib.Probability.logDecay p) Filter.atTop (nhds rate)
\end{ReviewVerbatim}

\paragraph{Exact Lean library declaration}
\begin{ReviewVerbatim}
def HasExponentialRate (p : ℕ → ℝ) (rate : ℝ) : Prop :=
  Tendsto (logDecay p) atTop (nhds rate)
\end{ReviewVerbatim}

\paragraph{Recorded source-to-library screening}
\textbf{Verdict:} \texttt{not recorded}
\reviewmatch{Matches source input}
\noindent\textbf{Reviewer annotation}\par
\reviewerbox
\clearpage
\hypertarget{library-prerequisite-e43b036a6f31f38f}{}
\subsection*{\small\ttfamily\raggedright EconCSLib.\allowbreak{}Probability.\allowbreak{}HasPositiveWeightNearAEEssentialInfimum}
\textbf{Source connection:} no explicit byte-pinned paper source connection is registered.
\paragraph{Lean-expanded library semantic target}
\begin{ReviewVerbatim}
∀ {α : Type u_1} [inst : MeasurableSpace α] (μ : MeasureTheory.Measure α) (w phi : α → ℝ) (rate ε : ℝ),
  ε > 0 → ∃ s c, 0 < μ.real s ∧ 0 < c ∧ ∀ᵐ (x : α) ∂μ.restrict s, c ≤ w x ∧ phi x ≤ rate + ε
\end{ReviewVerbatim}

\paragraph{Exact Lean library declaration}
\begin{ReviewVerbatim}
def HasPositiveWeightNearAEEssentialInfimum
    {α : Type*} [MeasurableSpace α]
    (μ : Measure α) (w phi : α → ℝ) (rate : ℝ) : Prop :=
  ∀ ε > 0, ∃ s : Set α, ∃ c : ℝ,
    0 < μ.real s ∧ 0 < c ∧
      ∀ᵐ x ∂μ.restrict s, c ≤ w x ∧ phi x ≤ rate + ε
\end{ReviewVerbatim}

\paragraph{Recorded source-to-library screening}
\textbf{Verdict:} \texttt{not recorded}
\reviewmatch{Matches source input}
\noindent\textbf{Reviewer annotation}\par
\reviewerbox
\clearpage
\hypertarget{library-prerequisite-2f4464ce60bae569}{}
\subsection*{\small\ttfamily\raggedright EconCSLib.\allowbreak{}Probability.\allowbreak{}PairwiseThresholdRateTopLdpCertificate}
\textbf{Source connection:} no explicit byte-pinned paper source connection is registered.
\paragraph{Lean declaration type (metadata)}
\begin{ReviewVerbatim}
{Seller : Type u_1} →
  {Rating : Type u_2} →
    {Pair : Type u_3} →
      [inst : Fintype Rating] →
        [inst_1 : DecidableEq Rating] →
          EconCSLib.Probability.FiniteRatingLDPModel Seller Rating →
            (Seller → ℝ) → (Pair → Seller) → (Pair → Seller) → Type u_3
\end{ReviewVerbatim}

\paragraph{Exact Lean library declaration}
\begin{ReviewVerbatim}
structure PairwiseThresholdRateTopLdpCertificate
    {Seller Rating Pair : Type*} [Fintype Rating] [DecidableEq Rating]
    (M : FiniteRatingLDPModel Seller Rating) (sampleRate : Seller → ℝ)
    (pairHi pairLo : Pair → Seller) where
  rate : Pair → ℝ
  threshold_rate_top_eq :
    ∀ p : Pair,
      pairwiseSellerThresholdRateTop M sampleRate (pairHi p) (pairLo p) =
        (rate p : WithTop ℝ)
  leftTail :
    ∀ p : Pair,
      ExponentialRateCertificate
        (twoSampleFloorScoreGapLeftTailProb M sampleRate (pairHi p) (pairLo p))
        (rate p)
\end{ReviewVerbatim}

\paragraph{Recorded source-to-library screening}
\textbf{Verdict:} \texttt{not recorded}
\reviewmatch{Matches source input}
\noindent\textbf{Reviewer annotation}\par
\reviewerbox
\clearpage
\hypertarget{library-prerequisite-479fa520bc58d797}{}
\subsection*{\small\ttfamily\raggedright EconCSLib.\allowbreak{}Probability.\allowbreak{}bernoulliFirstEndpointEqualSplit}
\textbf{Source connection:} no explicit byte-pinned paper source connection is registered.
\paragraph{Lean-expanded library semantic target}
\begin{ReviewVerbatim}
(pHi : ℝ) → (1 - √(1 - pHi)) / 2
\end{ReviewVerbatim}

\paragraph{Exact Lean library declaration}
\begin{ReviewVerbatim}
def bernoulliFirstEndpointEqualSplit (pHi : ℝ) : ℝ :=
  (1 - Real.sqrt (1 - pHi)) / 2
\end{ReviewVerbatim}

\paragraph{Recorded source-to-library screening}
\textbf{Verdict:} \texttt{not recorded}
\reviewmatch{Matches source input}
\noindent\textbf{Reviewer annotation}\par
\reviewerbox
\clearpage
\hypertarget{library-prerequisite-c4d922c1dbf3d0a6}{}
\subsection*{\small\ttfamily\raggedright EconCSLib.\allowbreak{}Probability.\allowbreak{}bernoulliInteriorFailureSplitNumerator}
\textbf{Source connection:} no explicit byte-pinned paper source connection is registered.
\paragraph{Lean-expanded library semantic target}
\begin{ReviewVerbatim}
(pLo pHi : ℝ) → (√(1 - pLo) - √(1 - pHi)) ^ 2
\end{ReviewVerbatim}

\paragraph{Exact Lean library declaration}
\begin{ReviewVerbatim}
def bernoulliInteriorFailureSplitNumerator (pLo pHi : ℝ) : ℝ :=
  (Real.sqrt (1 - pLo) - Real.sqrt (1 - pHi)) ^ 2
\end{ReviewVerbatim}

\paragraph{Recorded source-to-library screening}
\textbf{Verdict:} \texttt{not recorded}
\reviewmatch{Matches source input}
\noindent\textbf{Reviewer annotation}\par
\reviewerbox
\clearpage
\hypertarget{library-prerequisite-48bd8c17f399b1fc}{}
\subsection*{\small\ttfamily\raggedright EconCSLib.\allowbreak{}Probability.\allowbreak{}bernoulliInteriorSuccessSplitNumerator}
\textbf{Source connection:} no explicit byte-pinned paper source connection is registered.
\paragraph{Lean-expanded library semantic target}
\begin{ReviewVerbatim}
(pLo pHi : ℝ) → (√pHi - √pLo) ^ 2
\end{ReviewVerbatim}

\paragraph{Exact Lean library declaration}
\begin{ReviewVerbatim}
def bernoulliInteriorSuccessSplitNumerator (pLo pHi : ℝ) : ℝ :=
  (Real.sqrt pHi - Real.sqrt pLo) ^ 2
\end{ReviewVerbatim}

\paragraph{Recorded source-to-library screening}
\textbf{Verdict:} \texttt{not recorded}
\reviewmatch{Matches source input}
\noindent\textbf{Reviewer annotation}\par
\reviewerbox
\clearpage
\hypertarget{library-prerequisite-17d704ff7a24219b}{}
\subsection*{\small\ttfamily\raggedright EconCSLib.\allowbreak{}Probability.\allowbreak{}bernoulliInteriorSplitDenominator}
\textbf{Source connection:} no explicit byte-pinned paper source connection is registered.
\paragraph{Lean-expanded library semantic target}
\begin{ReviewVerbatim}
(pLo pHi : ℝ) →
  EconCSLib.Probability.bernoulliInteriorFailureSplitNumerator pLo pHi +
    EconCSLib.Probability.bernoulliInteriorSuccessSplitNumerator pLo pHi
\end{ReviewVerbatim}

\paragraph{Exact Lean library declaration}
\begin{ReviewVerbatim}
def bernoulliInteriorSplitDenominator (pLo pHi : ℝ) : ℝ :=
  bernoulliInteriorFailureSplitNumerator pLo pHi +
    bernoulliInteriorSuccessSplitNumerator pLo pHi
\end{ReviewVerbatim}

\paragraph{Recorded source-to-library screening}
\textbf{Verdict:} \texttt{not recorded}
\reviewmatch{Matches source input}
\noindent\textbf{Reviewer annotation}\par
\reviewerbox
\clearpage
\hypertarget{library-prerequisite-57b352998b6e4a14}{}
\subsection*{\small\ttfamily\raggedright EconCSLib.\allowbreak{}Probability.\allowbreak{}bernoulliInteriorEqualSplit}
\textbf{Source connection:} no explicit byte-pinned paper source connection is registered.
\paragraph{Lean-expanded library semantic target}
\begin{ReviewVerbatim}
(pLo pHi : ℝ) →
  EconCSLib.Probability.bernoulliInteriorFailureSplitNumerator pLo pHi /
    EconCSLib.Probability.bernoulliInteriorSplitDenominator pLo pHi
\end{ReviewVerbatim}

\paragraph{Exact Lean library declaration}
\begin{ReviewVerbatim}
def bernoulliInteriorEqualSplit (pLo pHi : ℝ) : ℝ :=
  bernoulliInteriorFailureSplitNumerator pLo pHi /
    bernoulliInteriorSplitDenominator pLo pHi
\end{ReviewVerbatim}

\paragraph{Recorded source-to-library screening}
\textbf{Verdict:} \texttt{not recorded}
\reviewmatch{Matches source input}
\noindent\textbf{Reviewer annotation}\par
\reviewerbox
\clearpage
\hypertarget{library-prerequisite-ca060c260b48fe49}{}
\subsection*{\small\ttfamily\raggedright EconCSLib.\allowbreak{}Probability.\allowbreak{}bernoulliKL}
\textbf{Source connection:} no explicit byte-pinned paper source connection is registered.
\paragraph{Lean-expanded library semantic target}
\begin{ReviewVerbatim}
(a b : ℝ) → a * Real.log (a / b) + (1 - a) * Real.log ((1 - a) / (1 - b))
\end{ReviewVerbatim}

\paragraph{Exact Lean library declaration}
\begin{ReviewVerbatim}
def bernoulliKL (a b : ℝ) : ℝ :=
  a * Real.log (a / b) + (1 - a) * Real.log ((1 - a) / (1 - b))
\end{ReviewVerbatim}

\paragraph{Recorded source-to-library screening}
\textbf{Verdict:} \texttt{not recorded}
\reviewmatch{Matches source input}
\noindent\textbf{Reviewer annotation}\par
\reviewerbox
\clearpage
\hypertarget{library-prerequisite-2735ba3861b2f309}{}
\subsection*{\small\ttfamily\raggedright EconCSLib.\allowbreak{}Probability.\allowbreak{}bernoulliLastEndpointEqualSplit}
\textbf{Source connection:} no explicit byte-pinned paper source connection is registered.
\paragraph{Lean-expanded library semantic target}
\begin{ReviewVerbatim}
(pLo : ℝ) → (1 + √pLo) / 2
\end{ReviewVerbatim}

\paragraph{Exact Lean library declaration}
\begin{ReviewVerbatim}
def bernoulliLastEndpointEqualSplit (pLo : ℝ) : ℝ :=
  (1 + Real.sqrt pLo) / 2
\end{ReviewVerbatim}

\paragraph{Recorded source-to-library screening}
\textbf{Verdict:} \texttt{not recorded}
\reviewmatch{Matches source input}
\noindent\textbf{Reviewer annotation}\par
\reviewerbox
\clearpage
\hypertarget{library-prerequisite-4cea2ad3a0ac0028}{}
\subsection*{\small\ttfamily\raggedright EconCSLib.\allowbreak{}Probability.\allowbreak{}binaryRatingScore}
\textbf{Source connection:} no explicit byte-pinned paper source connection is registered.
\paragraph{Lean-expanded library semantic target}
\begin{ReviewVerbatim}
(b : Bool) → if b = true then 1 else 0
\end{ReviewVerbatim}

\paragraph{Exact Lean library declaration}
\begin{ReviewVerbatim}
def binaryRatingScore (b : Bool) : ℝ :=
  if b then 1 else 0
\end{ReviewVerbatim}

\paragraph{Recorded source-to-library screening}
\textbf{Verdict:} \texttt{not recorded}
\reviewmatch{Matches source input}
\noindent\textbf{Reviewer annotation}\par
\reviewerbox
\clearpage
\hypertarget{library-prerequisite-e0e20858ab2c39b5}{}
\subsection*{\small\ttfamily\raggedright EconCSLib.\allowbreak{}Probability.\allowbreak{}finiteIidScoreSum}
\textbf{Source connection:} no explicit byte-pinned paper source connection is registered.
\paragraph{Lean-expanded library semantic target}
\begin{ReviewVerbatim}
{α : Type u_1} → {ι : Type u_2} → [inst : Fintype ι] → (score : α → ℝ) → (sample : ι → α) → ∑ i, score (sample i)
\end{ReviewVerbatim}

\paragraph{Exact Lean library declaration}
\begin{ReviewVerbatim}
def finiteIidScoreSum {ι : Type*} [Fintype ι]
    (score : α → ℝ) (sample : ι → α) : ℝ :=
  ∑ i : ι, score (sample i)
\end{ReviewVerbatim}

\paragraph{Recorded source-to-library screening}
\textbf{Verdict:} \texttt{not recorded}
\reviewmatch{Matches source input}
\noindent\textbf{Reviewer annotation}\par
\reviewerbox
\clearpage
\hypertarget{library-prerequisite-f584f3b3a668a101}{}
\subsection*{\small\ttfamily\raggedright EconCSLib.\allowbreak{}Probability.\allowbreak{}floorSampleCount}
\textbf{Source connection:} no explicit byte-pinned paper source connection is registered.
\paragraph{Lean-expanded library semantic target}
\begin{ReviewVerbatim}
{Seller : Type u_1} → (sampleRate : Seller → ℝ) → (θ : Seller) → (k : ℕ) → ⌊↑k * sampleRate θ⌋₊
\end{ReviewVerbatim}

\paragraph{Exact Lean library declaration}
\begin{ReviewVerbatim}
def floorSampleCount {Seller : Type*} (sampleRate : Seller → ℝ)
    (θ : Seller) (k : ℕ) : ℕ :=
  Nat.floor ((k : ℝ) * sampleRate θ)
\end{ReviewVerbatim}

\paragraph{Recorded source-to-library screening}
\textbf{Verdict:} \texttt{not recorded}
\reviewmatch{Matches source input}
\noindent\textbf{Reviewer annotation}\par
\reviewerbox
\clearpage
\hypertarget{library-prerequisite-b6ce97a274cbdb86}{}
\subsection*{\small\ttfamily\raggedright EconCSLib.\allowbreak{}Probability.\allowbreak{}realBernoulliPMF}
\textbf{Source connection:} no explicit byte-pinned paper source connection is registered.
\paragraph{Lean-expanded library semantic target}
\begin{ReviewVerbatim}
(p : ℝ) → (hp0 : 0 ≤ p) → (hp1 : p ≤ 1) → PMF.bernoulli ⟨p, hp0⟩ ⋯
\end{ReviewVerbatim}

\paragraph{Exact Lean library declaration}
\begin{ReviewVerbatim}
def realBernoulliPMF (p : ℝ) (hp0 : 0 ≤ p) (hp1 : p ≤ 1) : PMF Bool :=
  PMF.bernoulli ⟨p, hp0⟩ (by
    change p ≤ (1 : ℝ)
    exact hp1)
\end{ReviewVerbatim}

\paragraph{Recorded source-to-library screening}
\textbf{Verdict:} \texttt{not recorded}
\reviewmatch{Matches source input}
\noindent\textbf{Reviewer annotation}\par
\reviewerbox
\clearpage
\hypertarget{library-prerequisite-342012c3b44e591d}{}
\subsection*{\small\ttfamily\raggedright EconCSLib.\allowbreak{}Probability.\allowbreak{}realBinaryRatingLDPModel}
\textbf{Source connection:} no explicit byte-pinned paper source connection is registered.
\paragraph{Lean-expanded library semantic target}
\begin{ReviewVerbatim}
{θ : Type u_1} →
  (successProb : θ → ℝ) →
    (hprob0 : ∀ (t : θ), 0 ≤ successProb t) →
      (hprob1 : ∀ (t : θ), successProb t ≤ 1) →
        { typeLaw := fun t => EconCSLib.Probability.realBernoulliPMF (successProb t) ⋯ ⋯,
          score := EconCSLib.Probability.binaryRatingScore }
\end{ReviewVerbatim}

\paragraph{Exact Lean library declaration}
\begin{ReviewVerbatim}
def realBinaryRatingLDPModel {θ : Type*}
    (successProb : θ → ℝ)
    (hprob0 : ∀ t, 0 ≤ successProb t)
    (hprob1 : ∀ t, successProb t ≤ 1) :
    FiniteRatingLDPModel θ Bool where
  typeLaw := fun t => realBernoulliPMF (successProb t) (hprob0 t) (hprob1 t)
  score := binaryRatingScore
\end{ReviewVerbatim}

\paragraph{Recorded source-to-library screening}
\textbf{Verdict:} \texttt{not recorded}
\reviewmatch{Matches source input}
\noindent\textbf{Reviewer annotation}\par
\reviewerbox
\clearpage
\hypertarget{library-prerequisite-10ffb019e9ff740e}{}
\subsection*{\small\ttfamily\raggedright EconCSLib.\allowbreak{}Probability.\allowbreak{}twoBernoulliThresholdRate}
\textbf{Source connection:} no explicit byte-pinned paper source connection is registered.
\paragraph{Lean-expanded library semantic target}
\begin{ReviewVerbatim}
(gHi gLo tHi tLo a : ℝ) → gHi * EconCSLib.Probability.bernoulliKL a tHi + gLo * EconCSLib.Probability.bernoulliKL a tLo
\end{ReviewVerbatim}

\paragraph{Exact Lean library declaration}
\begin{ReviewVerbatim}
def twoBernoulliThresholdRate
    (gHi gLo tHi tLo a : ℝ) : ℝ :=
  gHi * bernoulliKL a tHi + gLo * bernoulliKL a tLo
\end{ReviewVerbatim}

\paragraph{Recorded source-to-library screening}
\textbf{Verdict:} \texttt{not recorded}
\reviewmatch{Matches source input}
\noindent\textbf{Reviewer annotation}\par
\reviewerbox
\clearpage
\hypertarget{library-prerequisite-f185c8e6394c9317}{}
\subsection*{\small\ttfamily\raggedright EconCSLib.\allowbreak{}pmfProd}
\textbf{Source connection:} no explicit byte-pinned paper source connection is registered.
\paragraph{Lean-expanded library semantic target}
\begin{ReviewVerbatim}
{α : Type u_1} →
  {β : Type u_2} →
    [inst : Fintype α] → [inst_1 : Fintype β] → (μ : PMF α) → (ν : PMF β) → PMF.ofFintype (fun p => μ p.1 * ν p.2) ⋯
\end{ReviewVerbatim}

\paragraph{Exact Lean library declaration}
\begin{ReviewVerbatim}
noncomputable def pmfProd {α β : Type*} [Fintype α] [Fintype β]
    (μ : PMF α) (ν : PMF β) : PMF (α × β) :=
  PMF.ofFintype (fun p : α × β => μ p.1 * ν p.2) (by
    classical
    have hμ : ∑ a : α, μ a = 1 := by
      rw [← PMF.tsum_coe μ, tsum_fintype]
    have hν : ∑ b : β, ν b = 1 := by
      rw [← PMF.tsum_coe ν, tsum_fintype]
    calc
      ∑ p : α × β, μ p.1 * ν p.2
          = ∑ a : α, ∑ b : β, μ a * ν b := by
            rw [Fintype.sum_prod_type]
      _ = ∑ a : α, μ a * (∑ b : β, ν b) := by
            simp [Finset.mul_sum]
      _ = 1 := by
            simp [hμ, hν])
\end{ReviewVerbatim}

\paragraph{Recorded source-to-library screening}
\textbf{Verdict:} \texttt{not recorded}
\reviewmatch{Matches source input}
\noindent\textbf{Reviewer annotation}\par
\reviewerbox
\clearpage
\hypertarget{library-prerequisite-38eb9a31a47bfbd9}{}
\subsection*{\small\ttfamily\raggedright EconCSLib.\allowbreak{}pmfProduct}
\textbf{Source connection:} no explicit byte-pinned paper source connection is registered.
\paragraph{Lean-expanded library semantic target}
\begin{ReviewVerbatim}
(ι : Type u_1) →
  (α : Type u_2) →
    [inst : Fintype ι] →
      [inst_1 : DecidableEq ι] → [inst_2 : Fintype α] → (μ : PMF α) → PMF.ofFintype (fun f => ∏ i, μ (f i)) ⋯
\end{ReviewVerbatim}

\paragraph{Exact Lean library declaration}
\begin{ReviewVerbatim}
noncomputable def pmfProduct (ι α : Type*) [Fintype ι] [DecidableEq ι] [Fintype α]
    (μ : PMF α) : PMF (ι → α) :=
  PMF.ofFintype (fun f : ι → α => ∏ i : ι, μ (f i)) (by
    classical
    have hcoord : ∑ a : α, μ a = 1 := by
      rw [← PMF.tsum_coe μ, tsum_fintype]
    calc
      ∑ f : ι → α, ∏ i : ι, μ (f i)
          = ∏ i : ι, ∑ a : α, μ a := by
            symm
            simpa using
              (Finset.prod_univ_sum
                (t := fun _i : ι => (Finset.univ : Finset α))
                (f := fun _i a => μ a))
      _ = 1 := by simp [hcoord])
\end{ReviewVerbatim}

\paragraph{Recorded source-to-library screening}
\textbf{Verdict:} \texttt{not recorded}
\reviewmatch{Matches source input}
\noindent\textbf{Reviewer annotation}\par
\reviewerbox
\clearpage
\hypertarget{library-prerequisite-59131a14024b99ca}{}
\subsection*{\small\ttfamily\raggedright EconCSLib.\allowbreak{}Probability.\allowbreak{}twoSampleRatingLaw}
\textbf{Source connection:} no explicit byte-pinned paper source connection is registered.
\paragraph{Lean-expanded library semantic target}
\begin{ReviewVerbatim}
{Seller : Type u_1} →
  {Rating : Type u_2} →
    [inst : Fintype Rating] →
      [inst_1 : DecidableEq Rating] →
        (M : EconCSLib.Probability.FiniteRatingLDPModel Seller Rating) →
          (hi lo : Seller) →
            (nHi nLo : ℕ) →
              EconCSLib.pmfProd (EconCSLib.pmfProduct (Fin nHi) Rating (M.typeLaw hi))
                (EconCSLib.pmfProduct (Fin nLo) Rating (M.typeLaw lo))
\end{ReviewVerbatim}

\paragraph{Exact Lean library declaration}
\begin{ReviewVerbatim}
noncomputable def twoSampleRatingLaw
    {Seller Rating : Type*} [Fintype Rating] [DecidableEq Rating]
    (M : FiniteRatingLDPModel Seller Rating) (hi lo : Seller)
    (nHi nLo : ℕ) :
    PMF ((Fin nHi → Rating) × (Fin nLo → Rating)) :=
  EconCSLib.pmfProd
    (EconCSLib.pmfProduct (Fin nHi) Rating (M.typeLaw hi))
    (EconCSLib.pmfProduct (Fin nLo) Rating (M.typeLaw lo))
\end{ReviewVerbatim}

\paragraph{Recorded source-to-library screening}
\textbf{Verdict:} \texttt{not recorded}
\reviewmatch{Matches source input}
\noindent\textbf{Reviewer annotation}\par
\reviewerbox
\clearpage
\hypertarget{library-prerequisite-0777cfce9fa4aa84}{}
\subsection*{\small\ttfamily\raggedright EconCSLib.\allowbreak{}Probability.\allowbreak{}twoSampleScoreGapSum}
\textbf{Source connection:} no explicit byte-pinned paper source connection is registered.
\paragraph{Lean-expanded library semantic target}
\begin{ReviewVerbatim}
{Seller : Type u_1} →
  {Rating : Type u_2} →
    [inst : Fintype Rating] →
      [inst_1 : DecidableEq Rating] →
        (M : EconCSLib.Probability.FiniteRatingLDPModel Seller Rating) →
          {nHi nLo : ℕ} →
            (cHi cLo : ℝ) →
              (sample : (Fin nHi → Rating) × (Fin nLo → Rating)) →
                cHi * EconCSLib.Probability.finiteIidScoreSum M.score sample.1 -
                  cLo * EconCSLib.Probability.finiteIidScoreSum M.score sample.2
\end{ReviewVerbatim}

\paragraph{Exact Lean library declaration}
\begin{ReviewVerbatim}
def twoSampleScoreGapSum
    {Seller Rating : Type*} [Fintype Rating] [DecidableEq Rating]
    (M : FiniteRatingLDPModel Seller Rating) {nHi nLo : ℕ}
    (cHi cLo : ℝ)
    (sample : (Fin nHi → Rating) × (Fin nLo → Rating)) : ℝ :=
  cHi * finiteIidScoreSum M.score sample.1 -
    cLo * finiteIidScoreSum M.score sample.2
\end{ReviewVerbatim}

\paragraph{Recorded source-to-library screening}
\textbf{Verdict:} \texttt{not recorded}
\reviewmatch{Matches source input}
\noindent\textbf{Reviewer annotation}\par
\reviewerbox
\clearpage
\hypertarget{library-prerequisite-e5c7eedc40c11b7a}{}
\subsection*{\small\ttfamily\raggedright EconCSLib.\allowbreak{}pmfExp}
\textbf{Source connection:} no explicit byte-pinned paper source connection is registered.
\paragraph{Lean-expanded library semantic target}
\begin{ReviewVerbatim}
{α : Type u_1} → [inst : Fintype α] → [DecidableEq α] → (μ : PMF α) → (f : α → ℝ) → ∑ a, (μ a).toReal * f a
\end{ReviewVerbatim}

\paragraph{Exact Lean library declaration}
\begin{ReviewVerbatim}
noncomputable def pmfExp {α : Type*} [Fintype α] [DecidableEq α]
    (μ : PMF α) (f : α → ℝ) : ℝ :=
  ∑ a, (μ a).toReal * f a
\end{ReviewVerbatim}

\paragraph{Recorded source-to-library screening}
\textbf{Verdict:} \texttt{not recorded}
\reviewmatch{Matches source input}
\noindent\textbf{Reviewer annotation}\par
\reviewerbox
\clearpage
\hypertarget{library-prerequisite-7fec9dedef722fc8}{}
\subsection*{\small\ttfamily\raggedright EconCSLib.\allowbreak{}pmfProb}
\textbf{Source connection:} no explicit byte-pinned paper source connection is registered.
\paragraph{Lean-expanded library semantic target}
\begin{ReviewVerbatim}
{α : Type u_1} →
  [inst : Fintype α] →
    [inst_1 : DecidableEq α] →
      (μ : PMF α) → (p : α → Prop) → [inst_2 : DecidablePred p] → EconCSLib.pmfExp μ fun a => if p a then 1 else 0
\end{ReviewVerbatim}

\paragraph{Exact Lean library declaration}
\begin{ReviewVerbatim}
noncomputable def pmfProb {α : Type*} [Fintype α] [DecidableEq α]
    (μ : PMF α) (p : α → Prop) [DecidablePred p] : ℝ :=
  pmfExp μ (fun a => if p a then 1 else 0)
\end{ReviewVerbatim}

\paragraph{Recorded source-to-library screening}
\textbf{Verdict:} \texttt{not recorded}
\reviewmatch{Matches source input}
\noindent\textbf{Reviewer annotation}\par
\reviewerbox
\clearpage
\hypertarget{library-prerequisite-7977d3effede269a}{}
\subsection*{\small\ttfamily\raggedright EconCSLib.\allowbreak{}Probability.\allowbreak{}twoSampleScoreGapStrictLeftProb}
\textbf{Source connection:} no explicit byte-pinned paper source connection is registered.
\paragraph{Lean-expanded library semantic target}
\begin{ReviewVerbatim}
{Seller : Type u_1} →
  {Rating : Type u_2} →
    [inst : Fintype Rating] →
      [inst_1 : DecidableEq Rating] →
        (M : EconCSLib.Probability.FiniteRatingLDPModel Seller Rating) →
          (hi lo : Seller) →
            (nHi nLo : ℕ) →
              (cHi cLo : ℝ) →
                EconCSLib.pmfProb (EconCSLib.Probability.twoSampleRatingLaw M hi lo nHi nLo) fun sample =>
                  EconCSLib.Probability.twoSampleScoreGapSum M cHi cLo sample < 0
\end{ReviewVerbatim}

\paragraph{Exact Lean library declaration}
\begin{ReviewVerbatim}
def twoSampleScoreGapStrictLeftProb
    {Seller Rating : Type*} [Fintype Rating] [DecidableEq Rating]
    (M : FiniteRatingLDPModel Seller Rating) (hi lo : Seller)
    (nHi nLo : ℕ) (cHi cLo : ℝ) : ℝ :=
  EconCSLib.pmfProb (twoSampleRatingLaw M hi lo nHi nLo)
    (fun sample => twoSampleScoreGapSum M cHi cLo sample < 0)
\end{ReviewVerbatim}

\paragraph{Recorded source-to-library screening}
\textbf{Verdict:} \texttt{not recorded}
\reviewmatch{Matches source input}
\noindent\textbf{Reviewer annotation}\par
\reviewerbox
\clearpage
\hypertarget{library-prerequisite-04fdd75df8db4e26}{}
\subsection*{\small\ttfamily\raggedright EconCSLib.\allowbreak{}Probability.\allowbreak{}twoSampleScoreGapTieProb}
\textbf{Source connection:} no explicit byte-pinned paper source connection is registered.
\paragraph{Lean-expanded library semantic target}
\begin{ReviewVerbatim}
{Seller : Type u_1} →
  {Rating : Type u_2} →
    [inst : Fintype Rating] →
      [inst_1 : DecidableEq Rating] →
        (M : EconCSLib.Probability.FiniteRatingLDPModel Seller Rating) →
          (hi lo : Seller) →
            (nHi nLo : ℕ) →
              (cHi cLo : ℝ) →
                EconCSLib.pmfProb (EconCSLib.Probability.twoSampleRatingLaw M hi lo nHi nLo) fun sample =>
                  EconCSLib.Probability.twoSampleScoreGapSum M cHi cLo sample = 0
\end{ReviewVerbatim}

\paragraph{Exact Lean library declaration}
\begin{ReviewVerbatim}
def twoSampleScoreGapTieProb
    {Seller Rating : Type*} [Fintype Rating] [DecidableEq Rating]
    (M : FiniteRatingLDPModel Seller Rating) (hi lo : Seller)
    (nHi nLo : ℕ) (cHi cLo : ℝ) : ℝ :=
  EconCSLib.pmfProb (twoSampleRatingLaw M hi lo nHi nLo)
    (fun sample => twoSampleScoreGapSum M cHi cLo sample = 0)
\end{ReviewVerbatim}

\paragraph{Recorded source-to-library screening}
\textbf{Verdict:} \texttt{not recorded}
\reviewmatch{Matches source input}
\noindent\textbf{Reviewer annotation}\par
\reviewerbox
\clearpage
\hypertarget{library-prerequisite-ba4c2212912ca97d}{}
\subsection*{\small\ttfamily\raggedright EconCSLib.\allowbreak{}Probability.\allowbreak{}twoSamplePkComplementErrorProb}
\textbf{Source connection:} no explicit byte-pinned paper source connection is registered.
\paragraph{Lean-expanded library semantic target}
\begin{ReviewVerbatim}
{Seller : Type u_1} →
  {Rating : Type u_2} →
    [inst : Fintype Rating] →
      [inst_1 : DecidableEq Rating] →
        (M : EconCSLib.Probability.FiniteRatingLDPModel Seller Rating) →
          (hi lo : Seller) →
            (nHi nLo : ℕ) →
              (cHi cLo : ℝ) →
                2 * EconCSLib.Probability.twoSampleScoreGapStrictLeftProb M hi lo nHi nLo cHi cLo +
                  EconCSLib.Probability.twoSampleScoreGapTieProb M hi lo nHi nLo cHi cLo
\end{ReviewVerbatim}

\paragraph{Exact Lean library declaration}
\begin{ReviewVerbatim}
def twoSamplePkComplementErrorProb
    {Seller Rating : Type*} [Fintype Rating] [DecidableEq Rating]
    (M : FiniteRatingLDPModel Seller Rating) (hi lo : Seller)
    (nHi nLo : ℕ) (cHi cLo : ℝ) : ℝ :=
  2 * twoSampleScoreGapStrictLeftProb M hi lo nHi nLo cHi cLo +
    twoSampleScoreGapTieProb M hi lo nHi nLo cHi cLo
\end{ReviewVerbatim}

\paragraph{Recorded source-to-library screening}
\textbf{Verdict:} \texttt{not recorded}
\reviewmatch{Matches source input}
\noindent\textbf{Reviewer annotation}\par
\reviewerbox
\clearpage
\hypertarget{library-prerequisite-94bd3b0a19bf1178}{}
\subsection*{\small\ttfamily\raggedright EconCSLib.\allowbreak{}Probability.\allowbreak{}twoSampleFloorPkComplementErrorProb}
\textbf{Source connection:} no explicit byte-pinned paper source connection is registered.
\paragraph{Lean-expanded library semantic target}
\begin{ReviewVerbatim}
{Seller : Type u_1} →
  {Rating : Type u_2} →
    [inst : Fintype Rating] →
      [inst_1 : DecidableEq Rating] →
        (M : EconCSLib.Probability.FiniteRatingLDPModel Seller Rating) →
          (sampleRate : Seller → ℝ) →
            (hi lo : Seller) →
              (k : ℕ) →
                have nHi := EconCSLib.Probability.floorSampleCount sampleRate hi k;
                have nLo := EconCSLib.Probability.floorSampleCount sampleRate lo k;
                EconCSLib.Probability.twoSamplePkComplementErrorProb M hi lo nHi nLo (↑nHi)⁻¹ (↑nLo)⁻¹
\end{ReviewVerbatim}

\paragraph{Exact Lean library declaration}
\begin{ReviewVerbatim}
def twoSampleFloorPkComplementErrorProb
    {Seller Rating : Type*} [Fintype Rating] [DecidableEq Rating]
    (M : FiniteRatingLDPModel Seller Rating) (sampleRate : Seller → ℝ)
    (hi lo : Seller) (k : ℕ) : ℝ :=
  let nHi := floorSampleCount sampleRate hi k
  let nLo := floorSampleCount sampleRate lo k
  twoSamplePkComplementErrorProb M hi lo nHi nLo
    ((nHi : ℝ)⁻¹) ((nLo : ℝ)⁻¹)
\end{ReviewVerbatim}

\paragraph{Recorded source-to-library screening}
\textbf{Verdict:} \texttt{not recorded}
\reviewmatch{Matches source input}
\noindent\textbf{Reviewer annotation}\par
\reviewerbox
\clearpage
\hypertarget{library-prerequisite-11866f137653a75d}{}
\subsection*{\small\ttfamily\raggedright EconCSLib.\allowbreak{}Probability.\allowbreak{}weightedBernoulliFailureBase}
\textbf{Source connection:} no explicit byte-pinned paper source connection is registered.
\paragraph{Lean-expanded library semantic target}
\begin{ReviewVerbatim}
(gHi gLo pHi pLo : ℝ) → (1 - pLo) ^ (gLo / (gHi + gLo)) * (1 - pHi) ^ (gHi / (gHi + gLo))
\end{ReviewVerbatim}

\paragraph{Exact Lean library declaration}
\begin{ReviewVerbatim}
def weightedBernoulliFailureBase (gHi gLo pHi pLo : ℝ) : ℝ :=
  (1 - pLo) ^ (gLo / (gHi + gLo)) *
    (1 - pHi) ^ (gHi / (gHi + gLo))
\end{ReviewVerbatim}

\paragraph{Recorded source-to-library screening}
\textbf{Verdict:} \texttt{not recorded}
\reviewmatch{Matches source input}
\noindent\textbf{Reviewer annotation}\par
\reviewerbox
\clearpage
\hypertarget{library-prerequisite-ffca3d1e44838990}{}
\subsection*{\small\ttfamily\raggedright EconCSLib.\allowbreak{}Probability.\allowbreak{}weightedBernoulliSuccessBase}
\textbf{Source connection:} no explicit byte-pinned paper source connection is registered.
\paragraph{Lean-expanded library semantic target}
\begin{ReviewVerbatim}
(gHi gLo pHi pLo : ℝ) → pLo ^ (gLo / (gHi + gLo)) * pHi ^ (gHi / (gHi + gLo))
\end{ReviewVerbatim}

\paragraph{Exact Lean library declaration}
\begin{ReviewVerbatim}
def weightedBernoulliSuccessBase (gHi gLo pHi pLo : ℝ) : ℝ :=
  pLo ^ (gLo / (gHi + gLo)) *
    pHi ^ (gHi / (gHi + gLo))
\end{ReviewVerbatim}

\paragraph{Recorded source-to-library screening}
\textbf{Verdict:} \texttt{not recorded}
\reviewmatch{Matches source input}
\noindent\textbf{Reviewer annotation}\par
\reviewerbox
\clearpage
\hypertarget{library-prerequisite-2d3d970688daf0d7}{}
\subsection*{\small\ttfamily\raggedright EconCSLib.\allowbreak{}Probability.\allowbreak{}weightedBernoulliClosedRateBase}
\textbf{Source connection:} no explicit byte-pinned paper source connection is registered.
\paragraph{Lean-expanded library semantic target}
\begin{ReviewVerbatim}
(gHi gLo pHi pLo : ℝ) →
  EconCSLib.Probability.weightedBernoulliFailureBase gHi gLo pHi pLo +
    EconCSLib.Probability.weightedBernoulliSuccessBase gHi gLo pHi pLo
\end{ReviewVerbatim}

\paragraph{Exact Lean library declaration}
\begin{ReviewVerbatim}
def weightedBernoulliClosedRateBase (gHi gLo pHi pLo : ℝ) : ℝ :=
  weightedBernoulliFailureBase gHi gLo pHi pLo +
    weightedBernoulliSuccessBase gHi gLo pHi pLo
\end{ReviewVerbatim}

\paragraph{Recorded source-to-library screening}
\textbf{Verdict:} \texttt{not recorded}
\reviewmatch{Matches source input}
\noindent\textbf{Reviewer annotation}\par
\reviewerbox
\clearpage
\hypertarget{library-prerequisite-c3cc92a498459481}{}
\subsection*{\small\ttfamily\raggedright EconCSLib.\allowbreak{}Probability.\allowbreak{}weightedBernoulliClosedThresholdRate}
\textbf{Source connection:} no explicit byte-pinned paper source connection is registered.
\paragraph{Lean-expanded library semantic target}
\begin{ReviewVerbatim}
(gHi gLo pHi pLo : ℝ) → -(gHi + gLo) * Real.log (EconCSLib.Probability.weightedBernoulliClosedRateBase gHi gLo pHi pLo)
\end{ReviewVerbatim}

\paragraph{Exact Lean library declaration}
\begin{ReviewVerbatim}
def weightedBernoulliClosedThresholdRate (gHi gLo pHi pLo : ℝ) : ℝ :=
  -(gHi + gLo) *
    Real.log (weightedBernoulliClosedRateBase gHi gLo pHi pLo)
\end{ReviewVerbatim}

\paragraph{Recorded source-to-library screening}
\textbf{Verdict:} \texttt{not recorded}
\reviewmatch{Matches source input}
\noindent\textbf{Reviewer annotation}\par
\reviewerbox
\clearpage
\hypertarget{library-prerequisite-ecc26a3133fdd38b}{}
\subsection*{\small\ttfamily\raggedright EconCSLib.\allowbreak{}finiteMax}
\textbf{Source connection:} no explicit byte-pinned paper source connection is registered.
\paragraph{Lean-expanded library semantic target}
\begin{ReviewVerbatim}
{α : Type u_1} → [inst : Fintype α] → [inst_1 : Nonempty α] → (f : α → ℝ) → Finset.univ.sup' ⋯ f
\end{ReviewVerbatim}

\paragraph{Exact Lean library declaration}
\begin{ReviewVerbatim}
def finiteMax {α : Type*} [Fintype α] [Nonempty α]
    (f : α → ℝ) : ℝ :=
  (Finset.univ : Finset α).sup' Finset.univ_nonempty f
\end{ReviewVerbatim}

\paragraph{Recorded source-to-library screening}
\textbf{Verdict:} \texttt{not recorded}
\reviewmatch{Matches source input}
\noindent\textbf{Reviewer annotation}\par
\reviewerbox
\clearpage
\hypertarget{library-prerequisite-087e3287a5c76915}{}
\subsection*{\small\ttfamily\raggedright EconCSLib.\allowbreak{}finiteMin}
\textbf{Source connection:} no explicit byte-pinned paper source connection is registered.
\paragraph{Lean-expanded library semantic target}
\begin{ReviewVerbatim}
{α : Type u_1} → [inst : Fintype α] → [inst_1 : Nonempty α] → (f : α → ℝ) → Finset.univ.inf' ⋯ f
\end{ReviewVerbatim}

\paragraph{Exact Lean library declaration}
\begin{ReviewVerbatim}
def finiteMin {α : Type*} [Fintype α] [Nonempty α]
    (f : α → ℝ) : ℝ :=
  (Finset.univ : Finset α).inf' Finset.univ_nonempty f
\end{ReviewVerbatim}

\paragraph{Recorded source-to-library screening}
\textbf{Verdict:} \texttt{not recorded}
\reviewmatch{Matches source input}
\noindent\textbf{Reviewer annotation}\par
\reviewerbox
\clearpage
\hypertarget{library-prerequisite-dae3217c424b40b2}{}
\subsection*{\small\ttfamily\raggedright EconCSLib.\allowbreak{}strictUpperPairSet}
\textbf{Source connection:} no explicit byte-pinned paper source connection is registered.
\paragraph{Lean-expanded library semantic target}
\begin{ReviewVerbatim}
∀ (a : ℝ × ℝ), {q | q.2 < q.1} a
\end{ReviewVerbatim}

\paragraph{Exact Lean library declaration}
\begin{ReviewVerbatim}
def strictUpperPairSet : Set (ℝ × ℝ) :=
  {q : ℝ × ℝ | q.2 < q.1}
\end{ReviewVerbatim}

\paragraph{Recorded source-to-library screening}
\textbf{Verdict:} \texttt{not recorded}
\reviewmatch{Matches source input}
\noindent\textbf{Reviewer annotation}\par
\reviewerbox

\clearpage
\hypertarget{source-claim-d4f7ad79732f053c}{}
\hypertarget{source-section-f10b5f68e4acf3dc}{}
\section*{Source claims}
\subsection*{Review row 1}
\textbf{Source locator:} {\footnotesize\raggedright cited publication, lines 187-\allowbreak{}204\par}
\paragraph{Verbatim source input}
\begin{ReviewVerbatim}
3

Model and optimization

We now formalize our model and show how to optimize the rating function to maximize the learning
rate. We focus on finding an optimal β : [0, 1] 7→
[0, 1], a map from item quality θ to the probability
it should receive a positive rating. This section requires no data: we characterize the optimal system.
3.1

Model and problem specification

Our model is constructed to emphasize the rating system’s learning rate. Time is discrete (k =
0, 1, 2, . . .). Informally: there is a set of items. Each
time step, buyers match with the items and leave a
rating according to β(θ). The platform records the
ratings and ranks the items. Formally:
\end{ReviewVerbatim}


\paragraph{Expanded PaperInterface specification}
\begin{ReviewVerbatim}
GJ19OptimalBinaryRatingSystems.ProofBridge.source_quality_domain = GJ19OptimalBinaryRatingSystems.sourceQualityDomain
\end{ReviewVerbatim}

\paragraph{Lean proof endpoint (built separately)}\begin{ReviewVerbatim}
GJ19OptimalBinaryRatingSystems.PaperInterface.source_quality_domain
\end{ReviewVerbatim}

\paragraph{Recorded source-to-Spec screening}
\textbf{Verdict:} \texttt{not recorded}
\\
\textbf{Reason:} not recorded
\reviewmatch{Matches source input}
\noindent\textbf{Reviewer annotation}\par
\reviewerbox
\clearpage
\hypertarget{source-claim-7979c322a62ef2c2}{}
\section*{Review row 2}
\textbf{Source locator:} {\footnotesize\raggedright cited publication, lines 199-\allowbreak{}218\par}
\paragraph{Verbatim source input}
\begin{ReviewVerbatim}
Our model is constructed to emphasize the rating system’s learning rate. Time is discrete (k =
0, 1, 2, . . .). Informally: there is a set of items. Each
time step, buyers match with the items and leave a
rating according to β(θ). The platform records the
ratings and ranks the items. Formally:

existence of a nondecreasing match function g(θ),
where item θ receives nk (θ) = bkg(θ)c matches, and
thus ratings, up to time k. In other words, item
1
time steps.
θ is matched approximately every g(θ)
g(θ) ≤ 1 and bounded away from 0, i.e. ∃c > 0:
g(θ) > c. This accumulation captures the feature
that better items may be more likely to match.
Ratings. The key quantity for our subsequent
analysis is the probability of a positive rating for
each θ, β(θ) , Pr(positive rating|θ). Let y` (θ) ∼
Bernoulli(β(θ)) be the rating an item of quality θ
receives at the `th time it matches.
\end{ReviewVerbatim}


\paragraph{Expanded PaperInterface specification}
\begin{ReviewVerbatim}
GJ19OptimalBinaryRatingSystems.ProofBridge.source_matching_count = GJ19OptimalBinaryRatingSystems.sourceMatchingCount
\end{ReviewVerbatim}

\paragraph{Lean proof endpoint (built separately)}\begin{ReviewVerbatim}
GJ19OptimalBinaryRatingSystems.PaperInterface.source_matching_count
\end{ReviewVerbatim}

\paragraph{Recorded source-to-Spec screening}
\textbf{Verdict:} \texttt{not recorded}
\\
\textbf{Reason:} not recorded
\reviewmatch{Matches source input}
\noindent\textbf{Reviewer annotation}\par
\reviewerbox
\clearpage
\hypertarget{source-claim-62613dce8c1c3874}{}
\section*{Review row 3}
\textbf{Source locator:} {\footnotesize\raggedright cited publication, lines 199-\allowbreak{}218\par}
\paragraph{Verbatim source input}
\begin{ReviewVerbatim}
Our model is constructed to emphasize the rating system’s learning rate. Time is discrete (k =
0, 1, 2, . . .). Informally: there is a set of items. Each
time step, buyers match with the items and leave a
rating according to β(θ). The platform records the
ratings and ranks the items. Formally:

existence of a nondecreasing match function g(θ),
where item θ receives nk (θ) = bkg(θ)c matches, and
thus ratings, up to time k. In other words, item
1
time steps.
θ is matched approximately every g(θ)
g(θ) ≤ 1 and bounded away from 0, i.e. ∃c > 0:
g(θ) > c. This accumulation captures the feature
that better items may be more likely to match.
Ratings. The key quantity for our subsequent
analysis is the probability of a positive rating for
each θ, β(θ) , Pr(positive rating|θ). Let y` (θ) ∼
Bernoulli(β(θ)) be the rating an item of quality θ
receives at the `th time it matches.
\end{ReviewVerbatim}


\paragraph{Expanded PaperInterface specification}
\begin{ReviewVerbatim}
GJ19OptimalBinaryRatingSystems.ProofBridge.source_matching_function =
  GJ19OptimalBinaryRatingSystems.SourceMatchingFunction
\end{ReviewVerbatim}

\paragraph{Lean proof endpoint (built separately)}\begin{ReviewVerbatim}
GJ19OptimalBinaryRatingSystems.PaperInterface.source_matching_function
\end{ReviewVerbatim}

\paragraph{Recorded source-to-Spec screening}
\textbf{Verdict:} \texttt{not recorded}
\\
\textbf{Reason:} not recorded
\reviewmatch{Matches source input}
\noindent\textbf{Reviewer annotation}\par
\reviewerbox
\clearpage
\hypertarget{source-claim-ad2abdefec264024}{}
\section*{Review row 4}
\textbf{Source locator:} {\footnotesize\raggedright cited publication, lines 219-\allowbreak{}237\par}
\paragraph{Verbatim source input}
\begin{ReviewVerbatim}
Aggregating ratings and ranking sellers.
These ratings are aggregated into a reputation score,
xk (θ), at each time k. The score is the fraction
of positive ratings received up to time k: xk (θ) ,
Pnk (θ)
1
`=0 y` (θ) with x0 (θ) , 0 for all θ. Thus,
nk (θ)
xk (θ) ∼ nk1(θ) Binomial(β(θ), nk (θ)).
System state. The state of the system is given by a
joint distribution µk (Θ, X), which gives the mass of
items of quality θ ∈ Θ ⊂ [0, 1] with aggregate score
xk (θ) ∈ X ⊂ [0, 1] at time k. Because our model is
a continuum, the evolution of the system state µk
follows a deterministic dynamical system.
We have described these dynamics at the level of individual items; however, such statements should be
interpreted as describing the evolution of the joint
distribution µk . The state update for µk is determined by the mass of items that match and the distributions of their ratings. A formal description of
the state evolution is in Appendix Section B.1.
\end{ReviewVerbatim}


\paragraph{Expanded PaperInterface specification}
\begin{ReviewVerbatim}
GJ19OptimalBinaryRatingSystems.ProofBridge.source_empirical_reputation_score =
  GJ19OptimalBinaryRatingSystems.sourceEmpiricalReputationScore
\end{ReviewVerbatim}

\paragraph{Lean proof endpoint (built separately)}\begin{ReviewVerbatim}
GJ19OptimalBinaryRatingSystems.PaperInterface.source_empirical_reputation_score
\end{ReviewVerbatim}

\paragraph{Recorded source-to-Spec screening}
\textbf{Verdict:} \texttt{not recorded}
\\
\textbf{Reason:} not recorded
\reviewmatch{Matches source input}
\noindent\textbf{Reviewer annotation}\par
\reviewerbox
\clearpage
\hypertarget{source-claim-ab291bf70e6ff37c}{}
\section*{Review row 5}
\textbf{Source locator:} {\footnotesize\raggedright cited publication, lines 228-\allowbreak{}237\par}
\paragraph{Verbatim source input}
\begin{ReviewVerbatim}
System state. The state of the system is given by a
joint distribution µk (Θ, X), which gives the mass of
items of quality θ ∈ Θ ⊂ [0, 1] with aggregate score
xk (θ) ∈ X ⊂ [0, 1] at time k. Because our model is
a continuum, the evolution of the system state µk
follows a deterministic dynamical system.
We have described these dynamics at the level of individual items; however, such statements should be
interpreted as describing the evolution of the joint
distribution µk . The state update for µk is determined by the mass of items that match and the distributions of their ratings. A formal description of
the state evolution is in Appendix Section B.1.
\end{ReviewVerbatim}


\paragraph{Expanded PaperInterface specification}
\begin{ReviewVerbatim}
GJ19OptimalBinaryRatingSystems.ProofBridge.source_system_state = GJ19OptimalBinaryRatingSystems.SourceSystemState
\end{ReviewVerbatim}

\paragraph{Lean proof endpoint (built separately)}\begin{ReviewVerbatim}
GJ19OptimalBinaryRatingSystems.PaperInterface.source_system_state
\end{ReviewVerbatim}

\paragraph{Recorded source-to-Spec screening}
\textbf{Verdict:} \texttt{not recorded}
\\
\textbf{Reason:} not recorded
\reviewmatch{Matches source input}
\noindent\textbf{Reviewer annotation}\par
\reviewerbox
\clearpage
\hypertarget{source-claim-b8da8ea96b79a0b8}{}
\section*{Review row 6}
\textbf{Source locator:} {\footnotesize\raggedright cited publication, lines 238-\allowbreak{}256\par}
\paragraph{Verbatim source input}
\begin{ReviewVerbatim}
Platform objective. The platform wishes to rank
the items accurately. Given β and θ1 > θ2 , define:
[U+0001 control character]
Pk (θ1 , θ2 |β) =µk xk (θ1 ) > xk (θ2 )|θ1 , θ2
[U+0001 control character]
− µk xk (θ1 ) < xk (θ2 )|θ1 , θ2
(1)

Items. The system consists of a set [0, 1] of items,
where each item is associated with a unique (but unknown) quality θ ∈ [0, 1]; i.e., the system consists of
a continuum of a unit mass of items whose unknown
qualities are uniform2 in [0, 1]. Below, we discretize
the continuous quality space [0, 1] into M types, to
calculate a stepwise increasing β. We will make clear
why we introduce a continuum but then discretize.

This expression captures observed score ranking’s
accuracy. When θ1 > θ2 but xk (θ1 ) < xk (θ2 ), the
ranking mistakenly orders θ1 below θ2 . A good system has large Pk (θ1 , θ2 |β). Integrating across items
\end{ReviewVerbatim}


\paragraph{Expanded PaperInterface specification}
\begin{ReviewVerbatim}
∀ (probabilityCorrect probabilityIncorrect : ℝ),
  GJ19OptimalBinaryRatingSystems.ProofBridge.source_definition_pairwise_accuracy_eq1 probabilityCorrect
      probabilityIncorrect =
    probabilityCorrect - probabilityIncorrect
\end{ReviewVerbatim}

\paragraph{Lean proof endpoint (built separately)}\begin{ReviewVerbatim}
GJ19OptimalBinaryRatingSystems.PaperInterface.source_definition_pairwise_accuracy_eq1
\end{ReviewVerbatim}

\paragraph{Recorded source-to-Spec screening}
\textbf{Verdict:} \texttt{not recorded}
\\
\textbf{Reason:} not recorded
\reviewmatch{Matches source input}
\noindent\textbf{Reviewer annotation}\par
\reviewerbox
\clearpage
\hypertarget{source-claim-88670d8e30834161}{}
\section*{Review row 7}
\textbf{Source locator:} {\footnotesize\raggedright cited publication, lines 257-\allowbreak{}288\par}
\paragraph{Verbatim source input}
\begin{ReviewVerbatim}
creates the following objective for each time k:
Z
Wk =
w(θ1 , θ2 )Pk (θ1 , θ2 |β)dθ1 dθ2
(2)

Matching with buyers. Items accumulate ratings
over time by matching with buyers. We assume the

Weight function w(θ1 , θ2 ) > 0 indicates how much
the platform cares about not mistaking a quality θ1
item with Ra quality θ2 item. We consider scaled w
such that θ1 >θ2 w(θ1 , θ2 )dθ1 dθ2 = 1.

1
Note that a rating is not the same as a reward; buyers
often give positive ratings after bad experiences.
2
Any distribution can be handled by considering θ to
be the item’s quantile rather than its absolute quality.

θ1 >θ2

Our first question then becomes: What β yields the
highest value of Wk ? As discussed above, the plat-


[form-feed in source extraction]
Designing Optimal Binary Rating Systems

form influences β through the design of its rating
system. The optimal choice of β sets the benchmark.
Discussion. Objective function. The specification
(2) of the objective is quite rich. It contains scaled
\end{ReviewVerbatim}


\paragraph{Expanded PaperInterface specification}
\begin{ReviewVerbatim}
∀ (weight pairwiseAccuracy : ℝ → ℝ → ℝ) (Wk : ℝ),
  (GJ19OptimalBinaryRatingSystems.ProofBridge.source_definition_weighted_objective_eq2 weight pairwiseAccuracy Wk =
      ∀ θ1 ∈ Set.Icc 0 1, ∀ θ2 ∈ Set.Ico 0 θ1, 0 < weight θ1 θ2) ∧
    ∫ (θ1 : ℝ) in Set.Icc 0 1, ∫ (θ2 : ℝ) in Set.Ico 0 θ1, weight θ1 θ2 = 1 ∧
      Wk = ∫ (θ1 : ℝ) in Set.Icc 0 1, ∫ (θ2 : ℝ) in Set.Ico 0 θ1, weight θ1 θ2 * pairwiseAccuracy θ1 θ2
\end{ReviewVerbatim}

\paragraph{Lean proof endpoint (built separately)}\begin{ReviewVerbatim}
GJ19OptimalBinaryRatingSystems.PaperInterface.source_definition_weighted_objective_eq2
\end{ReviewVerbatim}

\paragraph{Recorded source-to-Spec screening}
\textbf{Verdict:} \texttt{not recorded}
\\
\textbf{Reason:} not recorded
\reviewmatch{Matches source input}
\noindent\textbf{Reviewer annotation}\par
\reviewerbox
\clearpage
\hypertarget{source-claim-55d3b1e0bb545f86}{}
\section*{Review row 8}
\textbf{Source locator:} {\footnotesize\raggedright cited publication, lines 321-\allowbreak{}360\par}
\paragraph{Verbatim source input}
\begin{ReviewVerbatim}
Large deviations & discretization

Recall the question: What β yields the highest value
of Wk ?. We now refine objective Wk and constrain β
to form a non-degenerate, feasible optimization task.
Large deviation rate function. Wk is not one
objective: it has a different value per time k, and
no single β simultaneously optimizes Wk for all
3
We use θ1 θ2 (θ1 − θ2 ), (1 − θ1 )(1 − θ2 )(θ1 − θ2 ), and
1
( 2 − θ1 )2 ( 12 − θ2 )2 (θ1 − θ2 ) as examples.

k.4 Considering asymptotic performance is also
insufficient: when β is strictly increasing in θ,
limk→∞ Pk (θ1 , θ2 |β) = 1 ∀θ1 , θ2 by the law of large
numbers. Thus, W , limk→∞ Wk = 1, and any such
β is asymptotically optimal.
For this reason, we consider maximization of the
rate at which Wk converges, i.e., how fast the estimated item ranking converges to the true item ranking. We use a large deviations approach (Dembo
and Zeitouni, 2010) to quantify this convergence
rate. Formally, given sequence Yk ≤ limk→∞ Yk =
Y , the large deviations rate of convergence is
− limk→∞ k1 log(Y − Yk ) = c. If c exists then Yk approaches Y exponentially fast: Y − Yk = e−kc+o(k) .
Then, we wish to choose β to maximize Wk ’s large
deviations rate, r = − limk→∞ k1 log(W − Wk ).
Discretizing β. Unfortunately, even this problem
is degenerate if we consider continuous β: for any β
that is not piecewise constant, the large deviations
rate of convergence is zero, i.e., convergence of Wk
to its limit is only polynomially fast, and characterizing the dependence of this convergence rate on β
is intractable. Thus, the rate of convergence for Wk
is not a satisfactory objective with continuous β.
We make progress by discretizing β; in particular,
we restrict attention to optimization over stepwise
increasing β functions.5 Among stepwise increasing
β, the large deviations rate of Wk to its limiting
value W can be shown to be nondegenerate, i.e. ∃c >
0 s.t. W − Wk = e−kc+o(k) . (See Lemma C.4 in the
Appendix for further discussion.)
\end{ReviewVerbatim}


\paragraph{Expanded PaperInterface specification}
\begin{ReviewVerbatim}
∀ (W : ℝ) (Wk : ℕ → ℝ) (rate : ℝ),
  GJ19OptimalBinaryRatingSystems.ProofBridge.source_definition_large_deviation_rate W Wk rate =
    EconCSLib.Probability.HasExponentialRate (fun k => W - Wk k) rate
\end{ReviewVerbatim}

\paragraph{Lean proof endpoint (built separately)}\begin{ReviewVerbatim}
GJ19OptimalBinaryRatingSystems.PaperInterface.source_definition_large_deviation_rate
\end{ReviewVerbatim}

\paragraph{Recorded source-to-Spec screening}
\textbf{Verdict:} \texttt{not recorded}
\\
\textbf{Reason:} not recorded
\reviewmatch{Matches source input}
\noindent\textbf{Reviewer annotation}\par
\reviewerbox
\clearpage
\hypertarget{source-claim-9c0b6bf527ad55ca}{}
\section*{Review row 9}
\textbf{Source locator:} {\footnotesize\raggedright cited publication, lines 354-\allowbreak{}368\par}
\paragraph{Verbatim source input}
\begin{ReviewVerbatim}
We make progress by discretizing β; in particular,
we restrict attention to optimization over stepwise
increasing β functions.5 Among stepwise increasing
β, the large deviations rate of Wk to its limiting
value W can be shown to be nondegenerate, i.e. ∃c >
0 s.t. W − Wk = e−kc+o(k) . (See Lemma C.4 in the
Appendix for further discussion.)
Notationally, we will calculate an optimal stepwise
increasing β with M levels, i.e. there are M intervals
Si ⊂ [0, 1] and levels ti such that when θ1 , θ2 ∈ Si ,
then β(θ1 ) = β(θ2 ) , ti . The challenge is calculating an optimal S ∗ = {Si } and t∗ = {ti }.
The physical interpretation is that we group the
items into M subsets (types) Si ⊂ [0, 1]. When
items θ1 , θ2 are in the same subset, then their asymptotic reputation scores are the same, limk xk (θ1 ) =
limk xk (θ2 ) , ti . These items cannot be distinguished from one another even asymptotically.
\end{ReviewVerbatim}


\paragraph{Expanded PaperInterface specification}
\begin{ReviewVerbatim}
∀ (m : ℕ) (cutpoints : ℕ → ℝ) (levels : Fin (m + 2) → ℝ),
  GJ19OptimalBinaryRatingSystems.ProofBridge.source_definition_step_rule_partition_levels m cutpoints levels =
      GJ19OptimalBinaryRatingSystems.intervalCutpointsEndpointFeasible m cutpoints ∧
    GJ19OptimalBinaryRatingSystems.BinaryEndpointLevelVector levels
\end{ReviewVerbatim}

\paragraph{Lean proof endpoint (built separately)}\begin{ReviewVerbatim}
GJ19OptimalBinaryRatingSystems.PaperInterface.source_definition_step_rule_partition_levels
\end{ReviewVerbatim}

\paragraph{Recorded source-to-Spec screening}
\textbf{Verdict:} \texttt{not recorded}
\\
\textbf{Reason:} not recorded
\reviewmatch{Matches source input}
\noindent\textbf{Reviewer annotation}\par
\reviewerbox
\clearpage
\hypertarget{source-claim-b2f4ad11d62b1527}{}
\section*{Review row 10}
\textbf{Source locator:} {\footnotesize\raggedright cited publication, lines 459-\allowbreak{}476\par}
\paragraph{Verbatim source input}
\begin{ReviewVerbatim}
Our optimization problem: Within the class of
stepwise increasing functions with M levels, find the
β that is optimal, i.e. is

For Kendall’s τ and Spearman’s ρ, the optimal intervals are simply equispaced in [0, 1], i.e. Si∗ =
[ Mi , i+1
M ), because the entire item quality distribution is equally important. For other objective weight
functions w, the difficulty of finding the optimal subsets S ∗ depends on the properties of w. Since S ∗ is
trivial for Kendall’s τ and Spearman’s ρ – and w
is just an analytic tool that formalizes a platform’s
goals – we focus on finding the optimal levels t∗ .

(1) Asymptotically optimal. It yields the highest limiting value of Wk . AND
(2) Rate optimal. It yields the fastest large deviations rate r among asymptotically optimal β.
[U+0001 control character]
A remarkable result of our paper is a O M log2 M
[U+000F control character]
procedure to find an optimal β with M levels.
\end{ReviewVerbatim}


\paragraph{Expanded PaperInterface specification}
\begin{ReviewVerbatim}
∀ {α : Type u_1} (feasible : α → Prop) (primary secondary : α → ℝ) (x : α),
  feasible x ∧ ∀ (y : α), feasible y → primary y < primary x ∨ primary y = primary x ∧ secondary y ≤ secondary x
\end{ReviewVerbatim}

\paragraph{Lean proof endpoint (built separately)}\begin{ReviewVerbatim}
GJ19OptimalBinaryRatingSystems.PaperInterface.source_definition_lexicographic_optimality
\end{ReviewVerbatim}

\paragraph{Recorded source-to-Spec screening}
\textbf{Verdict:} \texttt{not recorded}
\\
\textbf{Reason:} not recorded
\reviewmatch{Matches source input}
\noindent\textbf{Reviewer annotation}\par
\reviewerbox
\clearpage
\hypertarget{source-claim-84eec684fc1865ab}{}
\section*{Review row 11}
\textbf{Source locator:} {\footnotesize\raggedright cited publication, lines 488-\allowbreak{}507\par}
\paragraph{Verbatim source input}
\begin{ReviewVerbatim}
Theorem 3.1. The β defined by the following
Suppose we started with a set of M items. Then the
choices of S ∗ and t∗ is optimal:
only remaining challenge is to equalize the rates at
R
P
which each item is separated from others: the large
∗
S = arg max 0≤i<j<M θ2 ∈Si ,θ1 ∈Sj w(θ1 , θ2 )d(θ1 , θ2 )
deviations rate is unaffected by the weight function
∗
6
w (it does not appear in the simplification of r(t)).
t = arg max r(t), where gi , inf θ∈Si g(θ) and
In other words, given a discrete set of M items
1
(equiv, given S ∗ ), calculating the optimal t is equivr(t) , − lim log(W − Wk )
(3)
alent to solving a maximin problem for the rates at
k→∞ k
\end{ReviewVerbatim}


\paragraph{Expanded PaperInterface specification}
\begin{ReviewVerbatim}
∀ (μ : MeasureTheory.Measure ℝ) [MeasureTheory.IsFiniteMeasure (μ.prod μ)] [(μ.prod μ).IsOpenPosMeasure]
  (S : GJ19OptimalBinaryRatingSystems.Theorem31SourceFiniteDiscretizationWeightedModel μ) (limitingValue : (ℕ → ℝ) → ℝ)
  (rate : (ℕ → ℝ) → (Fin (S.m + 2) → ℝ) → ℝ),
  EconCSLib.Optimization.IsMaximizerOn
      (GJ19OptimalBinaryRatingSystems.monotoneIntervalCutpointsEndpointFeasible (S.m + 2)) limitingValue S.cut →
    (∀ (cut : ℕ → ℝ),
        GJ19OptimalBinaryRatingSystems.monotoneIntervalCutpointsEndpointFeasible (S.m + 2) cut →
          limitingValue cut = limitingValue S.cut → cut = S.cut) →
      (∀ (levels : Fin (S.m + 2) → ℝ),
          GJ19OptimalBinaryRatingSystems.BinaryEndpointLevelVector levels →
            rate S.cut levels =
              GJ19OptimalBinaryRatingSystems.binaryEndpointAwareAdjacentRateObjective levels S.sampleRate) →
        ∃ levels,
          ∃ (hlevels : GJ19OptimalBinaryRatingSystems.BinaryEndpointLevelVector levels),
            GJ19OptimalBinaryRatingSystems.BinaryEndpointAwareAdjacentRatesEqualize levels S.sampleRate ∧
              EconCSLib.Probability.ExponentialRateCertificate
                  (GJ19OptimalBinaryRatingSystems.theorem31SourceWbar μ S.cut ⋯ S.sampleRate levels hlevels S.weight)
                  (GJ19OptimalBinaryRatingSystems.binaryEndpointAwareAdjacentRateObjective levels S.sampleRate) ∧
                EconCSLib.Optimization.IsLexicographicMaximizerOn
                  (fun design =>
                    GJ19OptimalBinaryRatingSystems.monotoneIntervalCutpointsEndpointFeasible (S.m + 2) design.1 ∧
                      GJ19OptimalBinaryRatingSystems.BinaryEndpointLevelVector design.2)
                  (fun design => limitingValue design.1) (fun design => rate design.1 design.2) (S.cut, levels)
\end{ReviewVerbatim}

\paragraph{Lean proof endpoint (built separately)}\begin{ReviewVerbatim}
GJ19OptimalBinaryRatingSystems.PaperInterface.theorem31_source_matching_function_unique_value_argmax_lexicographic
\end{ReviewVerbatim}

\paragraph{Recorded source-to-Spec screening}
\textbf{Verdict:} \texttt{not recorded}
\\
\textbf{Reason:} not recorded
\reviewmatch{Matches source input}
\noindent\textbf{Reviewer annotation}\par
\reviewerbox
\clearpage
\hypertarget{source-claim-8a7eb113492e3a82}{}
\section*{Review row 12}
\textbf{Source locator:} {\footnotesize\raggedright cited publication, lines 501-\allowbreak{}518\par}
\paragraph{Verbatim source input}
\begin{ReviewVerbatim}
t = arg max r(t), where gi , inf θ∈Si g(θ) and
In other words, given a discrete set of M items
1
(equiv, given S ∗ ), calculating the optimal t is equivr(t) , − lim log(W − Wk )
(3)
alent to solving a maximin problem for the rates at
k→∞ k
which
each type is distinguished from each of the
= min inf {gi+1 KL(a||ti+1 ) + gi KL(a||ti )}
0≤i≤M −2 a∈R
others. Thus, the algorithm below also solves the
inverse bandits problem in which we wish to rank
The proof is in the Appendix. The main hurdle is
the M arms, and we can choose the structure of the
showing that the continuum of rates for Pk (θ1 , θ2 )
(binary) observations at each arm.
for each pair θ1 , θ2 translates into a rate for Wk .
\end{ReviewVerbatim}


\paragraph{Expanded PaperInterface specification}
\begin{ReviewVerbatim}
∀ {m : ℕ} (cutpoints : ℕ → ℝ) (g : ℝ → ℝ) (j : Fin (m + 1)),
  Monotone cutpoints →
    Monotone g →
      ∀ (tHi tLo : ℝ),
        GJ19OptimalBinaryRatingSystems.sourceCellMatchingRate cutpoints g
              (GJ19OptimalBinaryRatingSystems.adjacentLowIndex j) =
            g (cutpoints ↑(GJ19OptimalBinaryRatingSystems.adjacentLowIndex j)) ∧
          GJ19OptimalBinaryRatingSystems.sourceCellMatchingRate cutpoints g
                (GJ19OptimalBinaryRatingSystems.adjacentHighIndex j) =
              g (cutpoints ↑(GJ19OptimalBinaryRatingSystems.adjacentHighIndex j)) ∧
            GJ19OptimalBinaryRatingSystems.paperAdjacentBinaryRatingRate
                (GJ19OptimalBinaryRatingSystems.sourceCellMatchingRate cutpoints g
                  (GJ19OptimalBinaryRatingSystems.adjacentHighIndex j))
                (GJ19OptimalBinaryRatingSystems.sourceCellMatchingRate cutpoints g
                  (GJ19OptimalBinaryRatingSystems.adjacentLowIndex j))
                tHi tLo =
              sInf
                (Set.range fun a =>
                  GJ19OptimalBinaryRatingSystems.sourceCellMatchingRate cutpoints g
                        (GJ19OptimalBinaryRatingSystems.adjacentHighIndex j) *
                      GJ19OptimalBinaryRatingSystems.paperBernoulliKL a tHi +
                    GJ19OptimalBinaryRatingSystems.sourceCellMatchingRate cutpoints g
                        (GJ19OptimalBinaryRatingSystems.adjacentLowIndex j) *
                      GJ19OptimalBinaryRatingSystems.paperBernoulliKL a tLo)
\end{ReviewVerbatim}

\paragraph{Lean proof endpoint (built separately)}\begin{ReviewVerbatim}
GJ19OptimalBinaryRatingSystems.PaperInterface.source_theorem31_adjacent_rate_eq3
\end{ReviewVerbatim}

\paragraph{Recorded source-to-Spec screening}
\textbf{Verdict:} \texttt{not recorded}
\\
\textbf{Reason:} not recorded
\reviewmatch{Matches source input}
\noindent\textbf{Reviewer annotation}\par
\reviewerbox
\clearpage
\hypertarget{source-claim-50a2d9627ea4a9d0}{}
\section*{Review row 13}
\textbf{Source locator:} {\footnotesize\raggedright cited publication, lines 501-\allowbreak{}518\par}
\paragraph{Verbatim source input}
\begin{ReviewVerbatim}
t = arg max r(t), where gi , inf θ∈Si g(θ) and
In other words, given a discrete set of M items
1
(equiv, given S ∗ ), calculating the optimal t is equivr(t) , − lim log(W − Wk )
(3)
alent to solving a maximin problem for the rates at
k→∞ k
which
each type is distinguished from each of the
= min inf {gi+1 KL(a||ti+1 ) + gi KL(a||ti )}
0≤i≤M −2 a∈R
others. Thus, the algorithm below also solves the
inverse bandits problem in which we wish to rank
The proof is in the Appendix. The main hurdle is
the M arms, and we can choose the structure of the
showing that the continuum of rates for Pk (θ1 , θ2 )
(binary) observations at each arm.
for each pair θ1 , θ2 translates into a rate for Wk .
\end{ReviewVerbatim}


\paragraph{Expanded PaperInterface specification}
\begin{ReviewVerbatim}
∀ {m : ℕ} (cut : ℕ → ℝ) (g : ℝ → ℝ) (i : Fin (m + 2)),
  cut ↑i ≤ cut (↑i + 1) →
    MonotoneOn g (Set.Icc (cut ↑i) (cut (↑i + 1))) →
      GJ19OptimalBinaryRatingSystems.sourceCellMatchingRate cut g i = g (cut ↑i)
\end{ReviewVerbatim}

\paragraph{Lean proof endpoint (built separately)}\begin{ReviewVerbatim}
GJ19OptimalBinaryRatingSystems.PaperInterface.theorem31_source_cell_matching_rate_eq_lower_cutpoint
\end{ReviewVerbatim}

\paragraph{Recorded source-to-Spec screening}
\textbf{Verdict:} \texttt{not recorded}
\\
\textbf{Reason:} not recorded
\reviewmatch{Matches source input}
\noindent\textbf{Reviewer annotation}\par
\reviewerbox
\clearpage
\hypertarget{source-claim-74286919e0e8cecc}{}
\section*{Review row 14}
\textbf{Source locator:} {\footnotesize\raggedright cited publication, lines 519-\allowbreak{}526\par}
\paragraph{Verbatim source input}
\begin{ReviewVerbatim}
We further note that the choice of M is not conse6
b
1−b
quential; in the Appendix Section B.4 we show that
KL(a||b) = a log a + (1 − a) log 1−a is the Kullbackin an appropriate sense, a sequence of optimal βM
Leibler (KL) divergence between Bernoulli random varifor
each M converges as M gets large.
ables with success probabilities a and b respectively.
\end{ReviewVerbatim}


\paragraph{Expanded PaperInterface specification}
\begin{ReviewVerbatim}
∀ (a b : ℝ),
  GJ19OptimalBinaryRatingSystems.paperBernoulliKL a b = a * Real.log (a / b) + (1 - a) * Real.log ((1 - a) / (1 - b))
\end{ReviewVerbatim}

\paragraph{Lean proof endpoint (built separately)}\begin{ReviewVerbatim}
GJ19OptimalBinaryRatingSystems.PaperInterface.definition_bernoulli_kl_formula
\end{ReviewVerbatim}

\paragraph{Recorded source-to-Spec screening}
\textbf{Verdict:} \texttt{not recorded}
\\
\textbf{Reason:} not recorded
\reviewmatch{Matches source input}
\noindent\textbf{Reviewer annotation}\par
\reviewerbox
\clearpage
\hypertarget{source-claim-d2fd3442bfa83879}{}
\section*{Review row 15}
\textbf{Source locator:} {\footnotesize\raggedright cited publication, lines 530-\allowbreak{}549\par}
\paragraph{Verbatim source input}
\begin{ReviewVerbatim}
Algorithm to find the optimal levels We now
describe how to find t∗ , the maximizer of r(t).
The following lemma describes a system of equations
to find the t∗ that maximizes r. It states that t∗
equalizes the rates at which each interval i is separated from its neighbors. The proof involves manipulation of r and convexity, and is in the appendix.
Lemma 3.1. The unique solution t∗ to the following
system of equations maximizes r(t):
r(t0 , t1 ) = r(t1 , t2 ) = · · · = r(tM −2 , tM −1 )

(4)

t0 = 0, tM −1 = 1
[U+0014 control character][U+0012 control character]
gi−1
gi
r(ti−1 , ti ) , − log (1 − ti−1 ) gi−1 +gi (1 − ti ) gi−1 +gi
gi−1

gi
\end{ReviewVerbatim}


\paragraph{Expanded PaperInterface specification}
\begin{ReviewVerbatim}
∀ {m : ℕ},
  1 < m →
    ∀ (sampleRate : Fin (m + 2) → ℝ),
      (∀ k < m + 2, 0 < GJ19OptimalBinaryRatingSystems.binaryEndpointSampleRateNat sampleRate k) →
        ∃! levels,
          GJ19OptimalBinaryRatingSystems.BinaryEndpointLevelVector levels ∧
            GJ19OptimalBinaryRatingSystems.BinaryEndpointAwareAdjacentRatesEqualize levels sampleRate ∧
              EconCSLib.Optimization.IsMaximizerOn GJ19OptimalBinaryRatingSystems.BinaryEndpointLevelVector
                (fun xs => GJ19OptimalBinaryRatingSystems.binaryEndpointAwareAdjacentRateObjective xs sampleRate) levels
\end{ReviewVerbatim}

\paragraph{Lean proof endpoint (built separately)}\begin{ReviewVerbatim}
GJ19OptimalBinaryRatingSystems.PaperInterface.source_lemma31_equalization_eq4
\end{ReviewVerbatim}

\paragraph{Recorded source-to-Spec screening}
\textbf{Verdict:} \texttt{not recorded}
\\
\textbf{Reason:} not recorded
\reviewmatch{Matches source input}
\noindent\textbf{Reviewer annotation}\par
\reviewerbox
\clearpage
\hypertarget{source-claim-e679d8fd634fddb9}{}
\section*{Review row 16}
\textbf{Source locator:} {\footnotesize\raggedright cited publication, lines 542-\allowbreak{}571\par}
\paragraph{Verbatim source input}
\begin{ReviewVerbatim}
[U+0014 control character][U+0012 control character]
gi−1
gi
r(ti−1 , ti ) , − log (1 − ti−1 ) gi−1 +gi (1 − ti ) gi−1 +gi
gi−1

gi

ALGORITHM 1: Nested Bisection
Data: Number of intervals M , match function g
Result: β levels, i.e. {t0 . . . tM −1 }
Function main (M , δ, g)
while Uncertainty region for tM −2 is bigger
than error tolerance do
Calculate r(tM −2 , 1), the rate between
current guess for tM −2 , and tM −1 = 1.
Fixing tM −2 , find t1 . . . tM −3 such that
r(t1 , t2 ) = r(t2 , t3 ) = · · · =
r(tM −3 , tM −2 ) = r(tM −2 , 1), which can be
done through a sequence of bisection
routines.
Calculate r(0, t1 ), the rate between current
guess for t1 , and t0 = 0.
Compare r(tM −2 , 1) and r(0, t1 ), adjust
uncertainty region for tM −2 accordingly.
return {ti }

[U+0013 control character]gi−1 +gi [U+0015 control character]

+ti−1 gi−1 +gi ti gi−1 +gi
\end{ReviewVerbatim}


\paragraph{Expanded PaperInterface specification}
\begin{ReviewVerbatim}
∀ (gLo gHi tLo tHi : ℝ),
  GJ19OptimalBinaryRatingSystems.paperAdjacentBinaryClosedRate gLo gHi tLo tHi =
    -(gLo + gHi) *
      Real.log
        ((1 - tLo) ^ (gLo / (gLo + gHi)) * (1 - tHi) ^ (gHi / (gLo + gHi)) +
          tLo ^ (gLo / (gLo + gHi)) * tHi ^ (gHi / (gLo + gHi)))
\end{ReviewVerbatim}

\paragraph{Lean proof endpoint (built separately)}\begin{ReviewVerbatim}
GJ19OptimalBinaryRatingSystems.PaperInterface.lemma31_closed_adjacent_rate_formula
\end{ReviewVerbatim}

\paragraph{Recorded source-to-Spec screening}
\textbf{Verdict:} \texttt{not recorded}
\\
\textbf{Reason:} not recorded
\reviewmatch{Matches source input}
\noindent\textbf{Reviewer annotation}\par
\reviewerbox
\clearpage
\hypertarget{source-claim-71d943c2f56bac62}{}
\section*{Review row 17}
\textbf{Source locator:} {\footnotesize\raggedright cited publication, lines 550-\allowbreak{}587\par}
\paragraph{Verbatim source input}
\begin{ReviewVerbatim}
ALGORITHM 1: Nested Bisection
Data: Number of intervals M , match function g
Result: β levels, i.e. {t0 . . . tM −1 }
Function main (M , δ, g)
while Uncertainty region for tM −2 is bigger
than error tolerance do
Calculate r(tM −2 , 1), the rate between
current guess for tM −2 , and tM −1 = 1.
Fixing tM −2 , find t1 . . . tM −3 such that
r(t1 , t2 ) = r(t2 , t3 ) = · · · =
r(tM −3 , tM −2 ) = r(tM −2 , 1), which can be
done through a sequence of bisection
routines.
Calculate r(0, t1 ), the rate between current
guess for t1 , and t0 = 0.
Compare r(tM −2 , 1) and r(0, t1 ), adjust
uncertainty region for tM −2 accordingly.
return {ti }

[U+0013 control character]gi−1 +gi [U+0015 control character]

+ti−1 gi−1 +gi ti gi−1 +gi
We do not know of any algorithm that efficiently
and provably solves such convex equality systems
in general. However, we leverage some structure in
our setting to develop an algorithm, NestedBisection, with run-time and optimality guarantees. The
efficiency of our algorithm results from the property
that, given a rate, ti is uniquely determined by the
value of either of the adjacent levels ti−1 , ti+1 , reducing an exponentially large search space to an almost
linear one. Physically, i.e., we only need to separate
each type of item from its neighboring types.
Below we include pseudo-code. Akin to branch and
bound, the algorithm proceeds via bisection on the
optimal value of tM −2 . For each candidate value
of tM −2 , the other values can be found using a sequence of bisection subroutines. These values approximately obey all the equalities in the system (4)
except the first. The direction of the first equality’s
violation reveals how to change the interval for the
next outer bisection iteration.
\end{ReviewVerbatim}


\paragraph{Expanded PaperInterface specification}
\begin{ReviewVerbatim}
∀ (m outerSteps innerSteps : ℕ) (sampleRate : Fin (m + 2) → ℝ) (grid : ℝ),
  GJ19OptimalBinaryRatingSystems.theorem32WeightedNestedBisectionOutput m outerSteps innerSteps sampleRate grid =
    have gLast :=
      sampleRate (GJ19OptimalBinaryRatingSystems.adjacentLowIndex GJ19OptimalBinaryRatingSystems.lastAdjacentIndex);
    have candidate := fun lastLow =>
      GJ19OptimalBinaryRatingSystems.theorem32WeightedBackwardGridRateBisectionLevels m innerSteps sampleRate grid
        (gLast * -Real.log lastLow) lastLow;
    have above := fun lastLow =>
      GJ19OptimalBinaryRatingSystems.theorem32OuterSourceWeightedRateAbove (candidate lastLow) sampleRate;
    have lastLow := (EconCSLib.Optimization.realBisectionRun above outerSteps (1 - 1 / ↑(m + 1)) (1 - grid)).2;
    candidate lastLow
\end{ReviewVerbatim}

\paragraph{Lean proof endpoint (built separately)}\begin{ReviewVerbatim}
GJ19OptimalBinaryRatingSystems.PaperInterface.theorem32_weighted_nested_bisection_output
\end{ReviewVerbatim}

\paragraph{Recorded source-to-Spec screening}
\textbf{Verdict:} \texttt{not recorded}
\\
\textbf{Reason:} not recorded
\reviewmatch{Matches source input}
\noindent\textbf{Reviewer annotation}\par
\reviewerbox
\clearpage
\hypertarget{source-claim-9c02a19f4b41b0c2}{}
\section*{Review row 18}
\textbf{Source locator:} {\footnotesize\raggedright cited publication, lines 588-\allowbreak{}598\par}
\paragraph{Verbatim source input}
\begin{ReviewVerbatim}
Theorem 3.2. [U+0001 control character]NestedBisection finds an [U+000F control character]-optimal t
in O M log2 M
operations, where [U+000F control character]-optimal means
[U+000F control character]
that r(t) is within additive constant [U+000F control character] of optimal.
The proof is in the appendix. The main difficulty
is finding a Lipschitz constant [U+000F control character](δ) for how much
the rate changes with a shift δ in a level ti . This
requires lower bounding t1 as a function of M . In
practice, the algorithm runs instantaneously on a
modern machine (e.g. for M = 200).
\end{ReviewVerbatim}


\paragraph{Expanded PaperInterface specification}
\begin{ReviewVerbatim}
∀ {m : ℕ},
  1 < m →
    ∀ (optimal sampleRate : Fin (m + 2) → ℝ) (eps : ℝ),
      GJ19OptimalBinaryRatingSystems.BinaryEndpointLevelVector optimal →
        GJ19OptimalBinaryRatingSystems.BinaryEndpointAwareAdjacentRatesEqualize optimal sampleRate →
          (∀ (idx : Fin (m + 2)), 0 < sampleRate idx) →
            (∀ {a b : Fin (m + 2)}, ↑a ≤ ↑b → sampleRate a ≤ sampleRate b) →
              sampleRate
                    (GJ19OptimalBinaryRatingSystems.adjacentHighIndex
                      GJ19OptimalBinaryRatingSystems.firstAdjacentIndex) =
                  1 →
                0 < eps →
                  have grid := GJ19OptimalBinaryRatingSystems.theorem32WeightedSourceGrid sampleRate eps;
                  have lastLower := 1 - 1 / ↑(m + 1);
                  have sourceDepthBudget := max ((1 - grid - lastLower) / 2) 1;
                  have runtimeLog := Real.logb 2 (max 1 (sourceDepthBudget / grid)) + 2;
                  0 < grid ∧
                    ∃ L,
                      ↑(L + 1) ≤ runtimeLog ∧
                        GJ19OptimalBinaryRatingSystems.binaryEndpointAwareAdjacentRateObjective optimal sampleRate -
                              GJ19OptimalBinaryRatingSystems.binaryEndpointAwareAdjacentRateObjective
                                (GJ19OptimalBinaryRatingSystems.theorem32WeightedNestedBisectionOutput m (L + 1) L
                                  sampleRate grid)
                                sampleRate ≤
                            eps ∧
                          GJ19OptimalBinaryRatingSystems.nestedBisectionOperationCount (m + 2) (L + 1) L ≤
                              EconCSLib.Optimization.nestedBisectionStepBound (m + 2) L ∧
                            ↑(GJ19OptimalBinaryRatingSystems.nestedBisectionOperationCount (m + 2) (L + 1) L) ≤
                              ↑(m + 2) * runtimeLog ^ 2
\end{ReviewVerbatim}

\paragraph{Lean proof endpoint (built separately)}\begin{ReviewVerbatim}
GJ19OptimalBinaryRatingSystems.PaperInterface.theorem32_weighted_nested_bisection_loss_and_runtime
\end{ReviewVerbatim}

\paragraph{Recorded source-to-Spec screening}
\textbf{Verdict:} \texttt{not recorded}
\\
\textbf{Reason:} not recorded
\reviewmatch{Matches source input}
\noindent\textbf{Reviewer annotation}\par
\reviewerbox
\clearpage
\hypertarget{source-claim-7159be549a8d27ef}{}
\section*{Review row 19}
\textbf{Source locator:} {\footnotesize\raggedright cited publication, lines 642-\allowbreak{}668\par}
\paragraph{Verbatim source input}
\begin{ReviewVerbatim}
Resources available to platform. We suppose
the platform has a set Y of possible binary questions that it could ask a buyer, e.g., “Are you satisfied with your experience” or “Is this experience
your best on our platform?”. Informally, at each
rating opportunity (i.e., match made), the rater can
be shown a single question y ∈ Y. Let ψ(θ, y) be the
probability an item quality θ would receive a positive
rating when the rater is asked question y ∈ Y.
We further suppose the platform has a set Θ of representative items for which it has access to item quality. Θ, then, is the granularity at which the platform
can collect data about historical performance, or
otherwise get expert labels. (We assume M [U+001D control character] |Θ|).
Using known set Θ, the platform can run an experiment to create an estimate ψ̂(θ, y), ∀y ∈ Y, θ ∈ Θ.
Design heuristic. How can the platform build an
effective rating system using these primitives and
β? We consider the following heuristic design: the
platform randomly shows a question y ∈ Y to each
buyer. The choice of the platform is a distribution
H(y) over y ∈ Y; in other words, for the platform
the design of the system amounts to choosing the frequency with which each question is shown. At each
rating opportunity, y is chosen from Y according to
H, independently across opportunities.
Clearly, the probability that
P an item θ receives a
positive rating is: β̃(θ) , y∈Y ψ(θ, y)H(y).
We want a distribution H such that β̃(θ) = β(θ) for
all θ ∈ [0, 1], i.e., that the positive rating probability
for each item is exactly the optimal value. However,
there may not exist any set of questions Y with associated ψ and choice of H such that β̃ = β.7
\end{ReviewVerbatim}


\paragraph{Expanded PaperInterface specification}
\begin{ReviewVerbatim}
∀ {Y : Type u_1} [inst : Fintype Y] (H : Y → ℝ), (∀ (y : Y), 0 ≤ H y) ∧ ∑ y, H y = 1
\end{ReviewVerbatim}

\paragraph{Lean proof endpoint (built separately)}\begin{ReviewVerbatim}
GJ19OptimalBinaryRatingSystems.PaperInterface.section4_question_distribution
\end{ReviewVerbatim}

\paragraph{Recorded source-to-Spec screening}
\textbf{Verdict:} \texttt{not recorded}
\\
\textbf{Reason:} not recorded
\reviewmatch{Matches source input}
\noindent\textbf{Reviewer annotation}\par
\reviewerbox
\clearpage
\hypertarget{source-claim-a5f9ed46c3edb5c1}{}
\section*{Review row 20}
\textbf{Source locator:} {\footnotesize\raggedright cited publication, lines 652-\allowbreak{}668\par}
\paragraph{Verbatim source input}
\begin{ReviewVerbatim}
Using known set Θ, the platform can run an experiment to create an estimate ψ̂(θ, y), ∀y ∈ Y, θ ∈ Θ.
Design heuristic. How can the platform build an
effective rating system using these primitives and
β? We consider the following heuristic design: the
platform randomly shows a question y ∈ Y to each
buyer. The choice of the platform is a distribution
H(y) over y ∈ Y; in other words, for the platform
the design of the system amounts to choosing the frequency with which each question is shown. At each
rating opportunity, y is chosen from Y according to
H, independently across opportunities.
Clearly, the probability that
P an item θ receives a
positive rating is: β̃(θ) , y∈Y ψ(θ, y)H(y).
We want a distribution H such that β̃(θ) = β(θ) for
all θ ∈ [0, 1], i.e., that the positive rating probability
for each item is exactly the optimal value. However,
there may not exist any set of questions Y with associated ψ and choice of H such that β̃ = β.7
\end{ReviewVerbatim}


\paragraph{Expanded PaperInterface specification}
\begin{ReviewVerbatim}
∀ {Y : Type u_1} [inst : Fintype Y] (ψ : ℝ → Y → ℝ) (H : Y → ℝ) (θ : ℝ),
  GJ19OptimalBinaryRatingSystems.sourceInducedBinaryResponse ψ H θ = ∑ y, ψ θ y * H y
\end{ReviewVerbatim}

\paragraph{Lean proof endpoint (built separately)}\begin{ReviewVerbatim}
GJ19OptimalBinaryRatingSystems.PaperInterface.section4_induced_binary_response
\end{ReviewVerbatim}

\paragraph{Recorded source-to-Spec screening}
\textbf{Verdict:} \texttt{not recorded}
\\
\textbf{Reason:} not recorded
\reviewmatch{Matches source input}
\noindent\textbf{Reviewer annotation}\par
\reviewerbox
\clearpage
\hypertarget{source-claim-3fed172202fb05fc}{}
\section*{Review row 21}
\textbf{Source locator:} {\footnotesize\raggedright cited publication, lines 669-\allowbreak{}696\par}
\paragraph{Verbatim source input}
\begin{ReviewVerbatim}
We propose the following heuristic to address this
difficulty. Choose a probability distribution H to
minimize the following L1 distance:
X
X
min
|β(θ) −
ψ̂(θ, y)H(y)|
(5)
H:kHk1 =1

θ∈Θ

y∈Y

This heuristic uses the data available to the platform, ψ̂(θ, y) for a set of items θ ∈ Θ, and designs H
so at least these items receive ratings close to their
optimal ratings β(θi ). Then, as long as ψ is wellbehaved, and Θ is representative of the full set, one
can hope that β̃(θ) u β(θ), for all θ ∈ [0, 1].
Discussion. Real-world analogue & constraints. A
special case of this system is already in place on
7

There are special cases where an exact solution exists. For example, let Y = [0, 1], and ψ(θ, y) = I[θ ≥ y].

many platforms, where the same question is always
shown. Static systems can be designed by restricting
H to only have mass at one question y. More generally, constraints can be used in optimization (5).
\end{ReviewVerbatim}


\paragraph{Expanded PaperInterface specification}
\begin{ReviewVerbatim}
∀ {Representative : Type u_1} {Y : Type u_2} [inst : Fintype Representative] [inst_1 : Fintype Y]
  (quality : Representative → ℝ) (β : ℝ → ℝ) (ψHat : Representative → Y → ℝ) (H : Y → ℝ),
  GJ19OptimalBinaryRatingSystems.sourceQuestionDesignL1Objective quality β ψHat H =
    ∑ i, |β (quality i) - ∑ y, ψHat i y * H y|
\end{ReviewVerbatim}

\paragraph{Lean proof endpoint (built separately)}\begin{ReviewVerbatim}
GJ19OptimalBinaryRatingSystems.PaperInterface.section4_l1_design_objective
\end{ReviewVerbatim}

\paragraph{Recorded source-to-Spec screening}
\textbf{Verdict:} \texttt{not recorded}
\\
\textbf{Reason:} not recorded
\reviewmatch{Matches source input}
\noindent\textbf{Reviewer annotation}\par
\reviewerbox
\clearpage
\hypertarget{source-claim-642a63ff09a11d8b}{}
\section*{Review row 22}
\textbf{Source locator:} {\footnotesize\raggedright cited publication, lines 669-\allowbreak{}696\par}
\paragraph{Verbatim source input}
\begin{ReviewVerbatim}
We propose the following heuristic to address this
difficulty. Choose a probability distribution H to
minimize the following L1 distance:
X
X
min
|β(θ) −
ψ̂(θ, y)H(y)|
(5)
H:kHk1 =1

θ∈Θ

y∈Y

This heuristic uses the data available to the platform, ψ̂(θ, y) for a set of items θ ∈ Θ, and designs H
so at least these items receive ratings close to their
optimal ratings β(θi ). Then, as long as ψ is wellbehaved, and Θ is representative of the full set, one
can hope that β̃(θ) u β(θ), for all θ ∈ [0, 1].
Discussion. Real-world analogue & constraints. A
special case of this system is already in place on
7

There are special cases where an exact solution exists. For example, let Y = [0, 1], and ψ(θ, y) = I[θ ≥ y].

many platforms, where the same question is always
shown. Static systems can be designed by restricting
H to only have mass at one question y. More generally, constraints can be used in optimization (5).
\end{ReviewVerbatim}


\paragraph{Expanded PaperInterface specification}
\begin{ReviewVerbatim}
∀ {Representative : Type u_1} {Y : Type u_2} [inst : Fintype Representative] [inst_1 : Fintype Y]
  (quality : Representative → ℝ) (β : ℝ → ℝ) (ψHat : Representative → Y → ℝ) (H : Y → ℝ),
  EconCSLib.Optimization.IsMinimizerOn GJ19OptimalBinaryRatingSystems.SourceQuestionDistribution
    (GJ19OptimalBinaryRatingSystems.sourceQuestionDesignL1Objective quality β ψHat) H
\end{ReviewVerbatim}

\paragraph{Lean proof endpoint (built separately)}\begin{ReviewVerbatim}
GJ19OptimalBinaryRatingSystems.PaperInterface.section4_question_design_solution
\end{ReviewVerbatim}

\paragraph{Recorded source-to-Spec screening}
\textbf{Verdict:} \texttt{not recorded}
\\
\textbf{Reason:} not recorded
\reviewmatch{Matches source input}
\noindent\textbf{Reviewer annotation}\par
\reviewerbox
\clearpage
\hypertarget{source-claim-ffa7fbb9f1c16e4b}{}
\section*{Review row 23}
\textbf{Source locator:} {\footnotesize\raggedright cited publication, lines 1820-\allowbreak{}1828\par}
\paragraph{Verbatim source input}
\begin{ReviewVerbatim}
Formal specification of system state update

Recall that µk (Θ, X) is the mass of items with true quality θ ∈ Θ ⊆ [0, 1] and a reputation score
x ∈ X ⊆ [0, 1] at time k. Let Ek = {θ : nk (θ) = nk−1 (θ) + 1}. These are the items who receive an
additional rating at time k; for all θ ∈ Ekc , nk (θ) = nk−1 (θ). Our system is completely deterministic,
and evolves according to the distributions of the individual seller dynamics.
For each θ ∈ Ek , x, x0 , define ω(θ, x, x0 ) as follows:
ω(θ, x, x0 ) = β(θ)I{nk (θ)x − nk−1 (θ)x0 = 1} + (1 − β(θ))I{nk (θ)x − nk−1 (θ)x0 = 0}.
\end{ReviewVerbatim}


\paragraph{Expanded PaperInterface specification}
\begin{ReviewVerbatim}
GJ19OptimalBinaryRatingSystems.ProofBridge.appendixB1_active_set =
  GJ19OptimalBinaryRatingSystems.sourceAppendixBActiveSet
\end{ReviewVerbatim}

\paragraph{Lean proof endpoint (built separately)}\begin{ReviewVerbatim}
GJ19OptimalBinaryRatingSystems.PaperInterface.appendixB1_active_set
\end{ReviewVerbatim}

\paragraph{Recorded source-to-Spec screening}
\textbf{Verdict:} \texttt{not recorded}
\\
\textbf{Reason:} not recorded
\reviewmatch{Matches source input}
\noindent\textbf{Reviewer annotation}\par
\reviewerbox
\clearpage
\hypertarget{source-claim-9053ffafcfb8db42}{}
\section*{Review row 24}
\textbf{Source locator:} {\footnotesize\raggedright cited publication, lines 1828-\allowbreak{}1831\par}
\paragraph{Verbatim source input}
\begin{ReviewVerbatim}
ω(θ, x, x0 ) = β(θ)I{nk (θ)x − nk−1 (θ)x0 = 1} + (1 − β(θ))I{nk (θ)x − nk−1 (θ)x0 = 0}.
Then ω gives the probability of transition from x0 to x when an item receives a rating. We then
have:
Z Z
\end{ReviewVerbatim}


\paragraph{Expanded PaperInterface specification}
\begin{ReviewVerbatim}
GJ19OptimalBinaryRatingSystems.ProofBridge.appendixB1_transition_kernel =
  GJ19OptimalBinaryRatingSystems.sourceAppendixBTransitionKernel
\end{ReviewVerbatim}

\paragraph{Lean proof endpoint (built separately)}\begin{ReviewVerbatim}
GJ19OptimalBinaryRatingSystems.PaperInterface.appendixB1_transition_kernel
\end{ReviewVerbatim}

\paragraph{Recorded source-to-Spec screening}
\textbf{Verdict:} \texttt{not recorded}
\\
\textbf{Reason:} not recorded
\reviewmatch{Matches source input}
\noindent\textbf{Reviewer annotation}\par
\reviewerbox
\clearpage
\hypertarget{source-claim-6724e8b8f279b968}{}
\section*{Review row 25}
\textbf{Source locator:} {\footnotesize\raggedright cited publication, lines 1832-\allowbreak{}1852\par}
\paragraph{Verbatim source input}
\begin{ReviewVerbatim}
Z
Z Z
1

µk+1 (Θ, X) =
Ek

x0 =0

ω(θ, x, x0 )dxµk (dx0 , dθ) +

x∈X

Ekc

µk (dx, dθ).
x∈X

It is straightforward but tedious to check that the preceding dynamics are well defined, given our
primitives.
\end{ReviewVerbatim}


\paragraph{Expanded PaperInterface specification}
\begin{ReviewVerbatim}
GJ19OptimalBinaryRatingSystems.ProofBridge.appendixB1_state_update =
  GJ19OptimalBinaryRatingSystems.sourceAppendixBStateUpdateMass
\end{ReviewVerbatim}

\paragraph{Lean proof endpoint (built separately)}\begin{ReviewVerbatim}
GJ19OptimalBinaryRatingSystems.PaperInterface.appendixB1_state_update
\end{ReviewVerbatim}

\paragraph{Recorded source-to-Spec screening}
\textbf{Verdict:} \texttt{not recorded}
\\
\textbf{Reason:} not recorded
\reviewmatch{Matches source input}
\noindent\textbf{Reviewer annotation}\par
\reviewerbox
\clearpage
\hypertarget{source-claim-52ebae870dd67c82}{}
\section*{Review row 26}
\textbf{Source locator:} {\footnotesize\raggedright cited publication, lines 1927-\allowbreak{}1949\par}
\paragraph{Verbatim source input}
\begin{ReviewVerbatim}
Formalization of effect of matching rates shifting

Matching concentrating at the top items moves mass of β(θ) away from high θ, and subsequently
mass of H(y) away from the questions that help distinguish the top items, as observed in Figures 1b and 3b above. Informally, this occurs because when matching concentrates, top items are
accumulating many ratings more ratings comparatively, and so the amount of information needed
per rating is comparatively less. We formalize this intuition in Lemma B.1 below.
The lemma states that if matching rates shift such that there is an index k above which matching
rates increase and below which they decrease, then correspondingly the levels of β, (i.e. ti ) become
closer together above k.
Lemma B.1. Suppose k, g, g̃ such that ∀j ∈ {k+1 . . . M −1}, gj ≥ g̃j , and ∀j ∈ {0 . . . k−1}, gj ≤ g̃j ,
and gk = g̃k . Then, t∗k ≥ t̃∗k .
Proof. This proof is similar to that of Lemma 3.1, except that with the matching function changing
the overall rate function can either increase, decrease, or stay the same. Suppose the overall rate
function decreased or stayed the same when the matching function changed from g̃ to g. Then
gM −2 > g̃M −2 and the target rate is no larger, and so t∗M −2 > t̃∗M −2 . Then, t∗M −3 > t̃∗M −3 (a smaller
width is needed because the matching rates are higher and the rate is no larger, and the next value
also increased). This shifting continues until t∗k+1 > t̃∗k+1 . Then, t∗k > t̃∗k .
Suppose instead that the overall rate function increased when the matching function changed
from g̃ to g. Then g1 < g̃1 and the target rate is larger, and so t∗1 > t̃∗1 . Then, t∗2 > t̃∗2 (a larger
width is needed and the previous value also increased). This shifting continues until t∗k−1 > t̃∗k−1 .
Then, t∗k > t̃∗k .
\end{ReviewVerbatim}


\paragraph{Expanded PaperInterface specification}
\begin{ReviewVerbatim}
∀ {m k : ℕ},
  0 < k →
    ∀ (hkm : k < m + 1) (sampleRate shiftedRate levels shiftedLevels : Fin (m + 2) → ℝ),
      GJ19OptimalBinaryRatingSystems.BinaryEndpointLevelVector levels →
        GJ19OptimalBinaryRatingSystems.BinaryEndpointLevelVector shiftedLevels →
          GJ19OptimalBinaryRatingSystems.BinaryEndpointAwareAdjacentRatesEqualize levels sampleRate →
            GJ19OptimalBinaryRatingSystems.BinaryEndpointAwareAdjacentRatesEqualize shiftedLevels shiftedRate →
              (∀ (i : Fin (m + 2)), 0 < sampleRate i) →
                (∀ (i : Fin (m + 2)), 0 < shiftedRate i) →
                  (∀ (i : Fin (m + 2)), k < ↑i → shiftedRate i ≤ sampleRate i) →
                    (∀ (i : Fin (m + 2)), ↑i < k → sampleRate i ≤ shiftedRate i) →
                      sampleRate ⟨k, ⋯⟩ = shiftedRate ⟨k, ⋯⟩ → shiftedLevels ⟨k, ⋯⟩ ≤ levels ⟨k, ⋯⟩
\end{ReviewVerbatim}

\paragraph{Lean proof endpoint (built separately)}\begin{ReviewVerbatim}
GJ19OptimalBinaryRatingSystems.PaperInterface.lemmaB1_matching_rate_shift
\end{ReviewVerbatim}

\paragraph{Recorded source-to-Spec screening}
\textbf{Verdict:} \texttt{not recorded}
\\
\textbf{Reason:} not recorded
\reviewmatch{Matches source input}
\noindent\textbf{Reviewer annotation}\par
\reviewerbox
\clearpage
\hypertarget{source-claim-f84fba16d2553222}{}
\section*{Review row 27}
\textbf{Source locator:} {\footnotesize\raggedright cited publication, lines 1951-\allowbreak{}1979\par}
\paragraph{Verbatim source input}
\begin{ReviewVerbatim}
Limit of β as M → ∞

w denote the optimal β with M intervals for weight function w, with intervals {S wM } =
Let βM
i
wM
wM . Let q
wM , swM ), i.e. the quantile of
{[si , swM
)}
and
levels
t
(θ)
=
i/M
when
θ
∈
[s
wM
i+1
i
i+1
interval item of type θ is in.
Then, we have the following convergence result for βM .

Theorem B.1. Let g be uniform. Suppose w such that qwM converges uniformly. Then, ∀C ∈
\end{ReviewVerbatim}


\paragraph{Expanded PaperInterface specification}
\begin{ReviewVerbatim}
∀ (M : ℕ) (cellIndex : ℝ → Fin M) (levels : Fin M → ℝ) (quantile beta : ℝ → ℝ),
  (GJ19OptimalBinaryRatingSystems.ProofBridge.source_theoremB1_quantile_representation M cellIndex levels quantile
        beta =
      ∀ (θ : ℝ), quantile θ = ↑↑(cellIndex θ) / ↑M) ∧
    ∀ (θ : ℝ), beta θ = levels (cellIndex θ)
\end{ReviewVerbatim}

\paragraph{Lean proof endpoint (built separately)}\begin{ReviewVerbatim}
GJ19OptimalBinaryRatingSystems.PaperInterface.source_theoremB1_quantile_representation
\end{ReviewVerbatim}

\paragraph{Recorded source-to-Spec screening}
\textbf{Verdict:} \texttt{not recorded}
\\
\textbf{Reason:} not recorded
\reviewmatch{Matches source input}
\noindent\textbf{Reviewer annotation}\par
\reviewerbox
\clearpage
\hypertarget{source-claim-5249ffde05e3b08f}{}
\section*{Review row 28}
\textbf{Source locator:} {\footnotesize\raggedright cited publication, lines 1980-\allowbreak{}2003\par}
\paragraph{Verbatim source input}
\begin{ReviewVerbatim}
w
w uniformly as N → ∞.
N, ∃β w s.t. βC2
N +1 → β
The proof is technical and is below. We leverage the fact that, for g uniform, the levels of β2M
can be analytically written as a function of the levels of βM . We believe (numerically observe)
that this theorem holds for the entire sequence as opposed to the each such subsequence, and for
general matching functions g. However, our proof technique does not carry over, and the proof
would leverage more global properties of the optimal βM .
Furthermore, the condition on w is light. For example, it holds for Kendall’s τ , Spearman’s ρ,
and all other examples mentioned in this work.
This convergence result suggests that the choice of M when calculating a asymptotic and rate
optimal β is not consequential. As M increases, the limiting value of Wk increases to 1 (i.e. the
asymptotic value increases), but the optimal rate decreases to 0. As discussed above, with strictly
increasing and continuous β, the asymptotic value is 1 but the large deviations rate does not exist,
i.e. convergence is polynomial.
This result could potentially be strengthened as follows: first, show convergence on the entire
sequence as opposed to these exponential subsequences, as conjectured; second, show desirable
properties of the limiting function itself. It is conceivable but not necessarily true that the limiting
9


[form-feed in source extraction]
function is “better” than other strictly continuous increasing functions in some rate sense, even
though the comparison through large deviations rate is degenerate.
\end{ReviewVerbatim}


\paragraph{Expanded PaperInterface specification}
\begin{ReviewVerbatim}
∀ (betaSeq quantileSeq : ℕ → ℝ → ℝ) (quantileLimit : ℝ → ℝ) (levels : (m : ℕ) → Fin (m + 2) → ℝ)
  (levelIndex : (m : ℕ) → ℝ → Fin (m + 2)),
  (∀ (m : ℕ) (θ : ℝ), betaSeq (m + 2) θ = levels m (levelIndex m θ)) →
    (∀ (m : ℕ),
        EconCSLib.Optimization.IsMaximizerOn GJ19OptimalBinaryRatingSystems.BinaryEndpointLevelVector
          (fun xs => GJ19OptimalBinaryRatingSystems.binaryEndpointAwareAdjacentRateObjective xs fun x => 1)
          (levels m)) →
      (∀ (m : ℕ), ∀ θ ∈ Set.Icc 0 1, ↑(levelIndex m θ) = min ⌊↑(m + 2) * quantileSeq (m + 2) θ⌋₊ (m + 1)) →
        (∀ (m : ℕ), ∀ θ ∈ Set.Icc 0 1, quantileSeq (m + 2) θ ∈ Set.Icc 0 1) →
          TendstoUniformlyOn quantileSeq quantileLimit Filter.atTop (Set.Icc 0 1) →
            GJ19OptimalBinaryRatingSystems.theoremB1UniformOptimalSubsequencePrincipleTo betaSeq quantileSeq
              quantileLimit
\end{ReviewVerbatim}

\paragraph{Lean proof endpoint (built separately)}\begin{ReviewVerbatim}
GJ19OptimalBinaryRatingSystems.PaperInterface.paper_theoremB1_uniform_subsequence_principle_to_of_quantile_floor_tendstoUniformlyOn_geometric_mesh
\end{ReviewVerbatim}

\paragraph{Recorded source-to-Spec screening}
\textbf{Verdict:} \texttt{not recorded}
\\
\textbf{Reason:} not recorded
\reviewmatch{Matches source input}
\noindent\textbf{Reviewer annotation}\par
\reviewerbox
\clearpage
\hypertarget{source-claim-c34337b3884bb5d6}{}
\section*{Review row 29}
\textbf{Source locator:} {\footnotesize\raggedright cited publication, lines 2005-\allowbreak{}2027\par}
\paragraph{Verbatim source input}
\begin{ReviewVerbatim}
Learning ψ(θ, y) through experiments

Now, we show how a platform would run an experiment to decide to learn ψ(θ, y). In particular,
one potential issue is that the platform does not have any items with know quality that it can use
as representative items in its optimization. In this case, we show that it can use ratings within the
experiment itself to identify these representative items. The results essentially follow from the law
of large numbers.
We assume that |Θ| = L representative items i ∈ {1 . . . L} are in the experiment, and each are
matched N times. The experiment proceeds as follows: every time an item is matched, show the
buyer a random question from Y. For each word y ∈ Y, track the empirical ψ̂(i, y), the proportion
of times a positive response was given to question y. Alternatively, if Y is totally ordered (i.e. a
positive rating for a given y also implies positive ratings would be given to all “easier” y 0 ), and
can be phrased as a multiple choice question, data collection can be faster: e.g., as we do in our
experiments: Y consists of a set of totally ordered adjectives that can describe the item; the rater is
asked to pick an adjective out of the set; this is interpreted as the item receiving a positive answer
to the questions induced by the chosen answer and all worse adjectives, and a negative answer to
all better adjectives.
First, suppose the platform approximately knows the quality θi of each item i, and θi are evenly
distributed in [0, 1]. Suppose the items are ordered by index, i.e. θ1 < θ2 < · · · < θL . Then let
ψ̂(θ, y) = ψ̂(i, y) when θ ∈ [θi−1 , θi ]. Call this procedure KnownTypeExperiment.
Lemma B.2. Suppose ψ(θ, y) is Lipschitz continuous in θ. With KnownTypeExperiment, ψ̂(i, y) →
ψ(θi , y)∀y uniformly as N → ∞. As L → ∞, ψ̂(θ, y) → ψ(θ, y)∀θ uniformly.
\end{ReviewVerbatim}


\paragraph{Expanded PaperInterface specification}
\begin{ReviewVerbatim}
∀ (Representative : Type u_1) (Y : Type u_2) [inst : Fintype Representative] [inst_1 : Fintype Y]
  (value : GJ19OptimalBinaryRatingSystems.SourceKnownTypeExperiment Representative Y),
  ∃ quality empiricalResponse, value.quality = quality ∧ value.empiricalResponse = empiricalResponse
\end{ReviewVerbatim}

\paragraph{Lean proof endpoint (built separately)}\begin{ReviewVerbatim}
GJ19OptimalBinaryRatingSystems.PaperInterface.appendixB5_known_type_experiment
\end{ReviewVerbatim}

\paragraph{Recorded source-to-Spec screening}
\textbf{Verdict:} \texttt{not recorded}
\\
\textbf{Reason:} not recorded
\reviewmatch{Matches source input}
\noindent\textbf{Reviewer annotation}\par
\reviewerbox
\clearpage
\hypertarget{source-claim-b26a46506065f742}{}
\section*{Review row 30}
\textbf{Source locator:} {\footnotesize\raggedright cited publication, lines 2005-\allowbreak{}2027\par}
\paragraph{Verbatim source input}
\begin{ReviewVerbatim}
Learning ψ(θ, y) through experiments

Now, we show how a platform would run an experiment to decide to learn ψ(θ, y). In particular,
one potential issue is that the platform does not have any items with know quality that it can use
as representative items in its optimization. In this case, we show that it can use ratings within the
experiment itself to identify these representative items. The results essentially follow from the law
of large numbers.
We assume that |Θ| = L representative items i ∈ {1 . . . L} are in the experiment, and each are
matched N times. The experiment proceeds as follows: every time an item is matched, show the
buyer a random question from Y. For each word y ∈ Y, track the empirical ψ̂(i, y), the proportion
of times a positive response was given to question y. Alternatively, if Y is totally ordered (i.e. a
positive rating for a given y also implies positive ratings would be given to all “easier” y 0 ), and
can be phrased as a multiple choice question, data collection can be faster: e.g., as we do in our
experiments: Y consists of a set of totally ordered adjectives that can describe the item; the rater is
asked to pick an adjective out of the set; this is interpreted as the item receiving a positive answer
to the questions induced by the chosen answer and all worse adjectives, and a negative answer to
all better adjectives.
First, suppose the platform approximately knows the quality θi of each item i, and θi are evenly
distributed in [0, 1]. Suppose the items are ordered by index, i.e. θ1 < θ2 < · · · < θL . Then let
ψ̂(θ, y) = ψ̂(i, y) when θ ∈ [θi−1 , θi ]. Call this procedure KnownTypeExperiment.
Lemma B.2. Suppose ψ(θ, y) is Lipschitz continuous in θ. With KnownTypeExperiment, ψ̂(i, y) →
ψ(θi , y)∀y uniformly as N → ∞. As L → ∞, ψ̂(θ, y) → ψ(θ, y)∀θ uniformly.
\end{ReviewVerbatim}


\paragraph{Expanded PaperInterface specification}
\begin{ReviewVerbatim}
∀ (Ω : Type u_1) (Item : Type u_2) (Y : Type u_3)
  (value : GJ19OptimalBinaryRatingSystems.SourceRandomQuestionExperiment Ω Item Y),
  ∃ question positiveResponse, value.question = question ∧ value.positiveResponse = positiveResponse
\end{ReviewVerbatim}

\paragraph{Lean proof endpoint (built separately)}\begin{ReviewVerbatim}
GJ19OptimalBinaryRatingSystems.PaperInterface.appendixB5_random_question_experiment
\end{ReviewVerbatim}

\paragraph{Recorded source-to-Spec screening}
\textbf{Verdict:} \texttt{not recorded}
\\
\textbf{Reason:} not recorded
\reviewmatch{Matches source input}
\noindent\textbf{Reviewer annotation}\par
\reviewerbox
\clearpage
\hypertarget{source-claim-8ccc8a8f737a5e7c}{}
\section*{Review row 31}
\textbf{Source locator:} {\footnotesize\raggedright cited publication, lines 2005-\allowbreak{}2027\par}
\paragraph{Verbatim source input}
\begin{ReviewVerbatim}
Learning ψ(θ, y) through experiments

Now, we show how a platform would run an experiment to decide to learn ψ(θ, y). In particular,
one potential issue is that the platform does not have any items with know quality that it can use
as representative items in its optimization. In this case, we show that it can use ratings within the
experiment itself to identify these representative items. The results essentially follow from the law
of large numbers.
We assume that |Θ| = L representative items i ∈ {1 . . . L} are in the experiment, and each are
matched N times. The experiment proceeds as follows: every time an item is matched, show the
buyer a random question from Y. For each word y ∈ Y, track the empirical ψ̂(i, y), the proportion
of times a positive response was given to question y. Alternatively, if Y is totally ordered (i.e. a
positive rating for a given y also implies positive ratings would be given to all “easier” y 0 ), and
can be phrased as a multiple choice question, data collection can be faster: e.g., as we do in our
experiments: Y consists of a set of totally ordered adjectives that can describe the item; the rater is
asked to pick an adjective out of the set; this is interpreted as the item receiving a positive answer
to the questions induced by the chosen answer and all worse adjectives, and a negative answer to
all better adjectives.
First, suppose the platform approximately knows the quality θi of each item i, and θi are evenly
distributed in [0, 1]. Suppose the items are ordered by index, i.e. θ1 < θ2 < · · · < θL . Then let
ψ̂(θ, y) = ψ̂(i, y) when θ ∈ [θi−1 , θi ]. Call this procedure KnownTypeExperiment.
Lemma B.2. Suppose ψ(θ, y) is Lipschitz continuous in θ. With KnownTypeExperiment, ψ̂(i, y) →
ψ(θi , y)∀y uniformly as N → ∞. As L → ∞, ψ̂(θ, y) → ψ(θ, y)∀θ uniformly.
\end{ReviewVerbatim}


\paragraph{Expanded PaperInterface specification}
\begin{ReviewVerbatim}
∀ {Ω : Type u_1} {Item : Type u_2} {Y : Type u_3} [inst : DecidableEq Y]
  (E : GJ19OptimalBinaryRatingSystems.SourceRandomQuestionExperiment Ω Item Y) (i : Item) (y : Y) (N : ℕ) (ω : Ω),
  GJ19OptimalBinaryRatingSystems.sourceExperimentEmpiricalQuestionResponse E i y N ω =
    GJ19OptimalBinaryRatingSystems.sourceExperimentPositiveQuestionCount E i y N ω /
      GJ19OptimalBinaryRatingSystems.sourceExperimentQuestionCount E i y N ω
\end{ReviewVerbatim}

\paragraph{Lean proof endpoint (built separately)}\begin{ReviewVerbatim}
GJ19OptimalBinaryRatingSystems.PaperInterface.appendixB5_empirical_question_response
\end{ReviewVerbatim}

\paragraph{Recorded source-to-Spec screening}
\textbf{Verdict:} \texttt{not recorded}
\\
\textbf{Reason:} not recorded
\reviewmatch{Matches source input}
\noindent\textbf{Reviewer annotation}\par
\reviewerbox
\clearpage
\hypertarget{source-claim-276fd5feab006dec}{}
\section*{Review row 32}
\textbf{Source locator:} {\footnotesize\raggedright cited publication, lines 2027-\allowbreak{}2031\par}
\paragraph{Verbatim source input}
\begin{ReviewVerbatim}
ψ(θi , y)∀y uniformly as N → ∞. As L → ∞, ψ̂(θ, y) → ψ(θ, y)∀θ uniformly.
Proof. The proof follows directly from the Strong Law of Large Numbers. As N → ∞, ∀i, ψ̂(i, x) →
ψ(θi , x) uniformly. Now, let L → ∞. ∀[U+000F control character], ∃L0 s.t. ∀L > L0 , ∀θ, ∃i s.t. |θ −θi | < [U+000F control character]. ψ(θ, x) is Lipschitz
in θ by assumption, and so ψ̂(θ, x) → ψ(θ, x) uniformly.
We now relax the assumption that the platform has an existing set of items with known qualities.
\end{ReviewVerbatim}


\paragraph{Expanded PaperInterface specification}
\begin{ReviewVerbatim}
∀ {Ω : Type u_1} {Y : Type u_2} [inst : MeasurableSpace Ω] [Fintype Y] [inst_2 : DecidableEq Y]
  (P : MeasureTheory.Measure Ω)
  (E : (L : ℕ) → GJ19OptimalBinaryRatingSystems.SourceRandomQuestionExperiment Ω (Fin (L + 1)) Y)
  (quality : (L : ℕ) → Fin (L + 1) → ℝ) (representative : (L : ℕ) → ℝ → Fin (L + 1)) (psi : ℝ → Y → ℝ) (H : Y → ℝ),
  (∀ (y : Y), 0 < H y) →
    (∀ (L : ℕ) (i : Fin (L + 1)) (y : Y),
        MeasureTheory.Integrable (GJ19OptimalBinaryRatingSystems.sourceQuestionAskedIndicator (E L) i y 0) P) →
      (∀ (L : ℕ) (i : Fin (L + 1)) (y : Y),
          Pairwise
            (Function.onFun (fun x1 x2 => ProbabilityTheory.IndepFun x1 x2 P)
              (GJ19OptimalBinaryRatingSystems.sourceQuestionAskedIndicator (E L) i y))) →
        (∀ (L : ℕ) (i : Fin (L + 1)) (y : Y) (n : ℕ),
            ProbabilityTheory.IdentDistrib (GJ19OptimalBinaryRatingSystems.sourceQuestionAskedIndicator (E L) i y n)
              (GJ19OptimalBinaryRatingSystems.sourceQuestionAskedIndicator (E L) i y 0) P P) →
          (∀ (L : ℕ) (i : Fin (L + 1)) (y : Y),
              ∫ (ω : Ω), GJ19OptimalBinaryRatingSystems.sourceQuestionAskedIndicator (E L) i y 0 ω ∂P = H y) →
            (∀ (L : ℕ) (i : Fin (L + 1)) (y : Y),
                MeasureTheory.Integrable (GJ19OptimalBinaryRatingSystems.sourceQuestionPositiveIndicator (E L) i y 0)
                  P) →
              (∀ (L : ℕ) (i : Fin (L + 1)) (y : Y),
                  Pairwise
                    (Function.onFun (fun x1 x2 => ProbabilityTheory.IndepFun x1 x2 P)
                      (GJ19OptimalBinaryRatingSystems.sourceQuestionPositiveIndicator (E L) i y))) →
                (∀ (L : ℕ) (i : Fin (L + 1)) (y : Y) (n : ℕ),
                    ProbabilityTheory.IdentDistrib
                      (GJ19OptimalBinaryRatingSystems.sourceQuestionPositiveIndicator (E L) i y n)
                      (GJ19OptimalBinaryRatingSystems.sourceQuestionPositiveIndicator (E L) i y 0) P P) →
                  (∀ (L : ℕ) (i : Fin (L + 1)) (y : Y),
                      ∫ (ω : Ω), GJ19OptimalBinaryRatingSystems.sourceQuestionPositiveIndicator (E L) i y 0 ω ∂P =
                        H y * psi (quality L i) y) →
                    (∀ δ > 0,
                        ∀ᶠ (L : ℕ) in Filter.atTop, ∀ θ ∈ Set.Icc 0 1, dist (quality L (representative L θ)) θ < δ) →
                      ∀ {K : ℝ},
                        0 ≤ K →
                          (∀ (θ θ' : ℝ) (y : Y), dist (psi θ y) (psi θ' y) ≤ K * dist θ θ') →
                            ∀ᵐ (ω : Ω) ∂P,
                              ∀ ε > 0,
                                ∀ᶠ (L : ℕ) (N : ℕ) in Filter.atTop,
                                  ∀ θ ∈ Set.Icc 0 1,
                                    ∀ (y : Y),
                                      dist
                                          (GJ19OptimalBinaryRatingSystems.sourceExperimentEmpiricalQuestionResponse
                                            (E L) (representative L θ) y N ω)
                                          (psi θ y) <
                                        ε
\end{ReviewVerbatim}

\paragraph{Lean proof endpoint (built separately)}\begin{ReviewVerbatim}
GJ19OptimalBinaryRatingSystems.PaperInterface.paper_lemmaB2_knownTypeExperiment_uniform_convergence_of_lipschitz_tracking
\end{ReviewVerbatim}

\paragraph{Recorded source-to-Spec screening}
\textbf{Verdict:} \texttt{not recorded}
\\
\textbf{Reason:} not recorded
\reviewmatch{Matches source input}
\noindent\textbf{Reviewer annotation}\par
\reviewerbox
\clearpage
\hypertarget{source-claim-76b42fc11ca5198e}{}
\section*{Review row 33}
\textbf{Source locator:} {\footnotesize\raggedright cited publication, lines 2005-\allowbreak{}2031\par}
\paragraph{Verbatim source input}
\begin{ReviewVerbatim}
Learning ψ(θ, y) through experiments

Now, we show how a platform would run an experiment to decide to learn ψ(θ, y). In particular,
one potential issue is that the platform does not have any items with know quality that it can use
as representative items in its optimization. In this case, we show that it can use ratings within the
experiment itself to identify these representative items. The results essentially follow from the law
of large numbers.
We assume that |Θ| = L representative items i ∈ {1 . . . L} are in the experiment, and each are
matched N times. The experiment proceeds as follows: every time an item is matched, show the
buyer a random question from Y. For each word y ∈ Y, track the empirical ψ̂(i, y), the proportion
of times a positive response was given to question y. Alternatively, if Y is totally ordered (i.e. a
positive rating for a given y also implies positive ratings would be given to all “easier” y 0 ), and
can be phrased as a multiple choice question, data collection can be faster: e.g., as we do in our
experiments: Y consists of a set of totally ordered adjectives that can describe the item; the rater is
asked to pick an adjective out of the set; this is interpreted as the item receiving a positive answer
to the questions induced by the chosen answer and all worse adjectives, and a negative answer to
all better adjectives.
First, suppose the platform approximately knows the quality θi of each item i, and θi are evenly
distributed in [0, 1]. Suppose the items are ordered by index, i.e. θ1 < θ2 < · · · < θL . Then let
ψ̂(θ, y) = ψ̂(i, y) when θ ∈ [θi−1 , θi ]. Call this procedure KnownTypeExperiment.
Lemma B.2. Suppose ψ(θ, y) is Lipschitz continuous in θ. With KnownTypeExperiment, ψ̂(i, y) →
ψ(θi , y)∀y uniformly as N → ∞. As L → ∞, ψ̂(θ, y) → ψ(θ, y)∀θ uniformly.
Proof. The proof follows directly from the Strong Law of Large Numbers. As N → ∞, ∀i, ψ̂(i, x) →
ψ(θi , x) uniformly. Now, let L → ∞. ∀[U+000F control character], ∃L0 s.t. ∀L > L0 , ∀θ, ∃i s.t. |θ −θi | < [U+000F control character]. ψ(θ, x) is Lipschitz
in θ by assumption, and so ψ̂(θ, x) → ψ(θ, x) uniformly.
We now relax the assumption that the platform has an existing set of items with known qualities.
\end{ReviewVerbatim}


\paragraph{Expanded PaperInterface specification}
\begin{ReviewVerbatim}
∀ {Ω : Type u_1} {Representative : Type u_2} {Y : Type u_3} [inst : MeasurableSpace Ω] [Fintype Representative]
  [Fintype Y] [inst_3 : DecidableEq Y] (P : MeasureTheory.Measure Ω)
  (E : GJ19OptimalBinaryRatingSystems.SourceRandomQuestionExperiment Ω Representative Y) (quality : Representative → ℝ)
  (ψ : ℝ → Y → ℝ) (H : Y → ℝ),
  (∀ (y : Y), 0 < H y) →
    (∀ (i : Representative) (y : Y),
        MeasureTheory.Integrable (GJ19OptimalBinaryRatingSystems.sourceQuestionAskedIndicator E i y 0) P) →
      (∀ (i : Representative) (y : Y),
          Pairwise
            (Function.onFun (fun x1 x2 => ProbabilityTheory.IndepFun x1 x2 P)
              (GJ19OptimalBinaryRatingSystems.sourceQuestionAskedIndicator E i y))) →
        (∀ (i : Representative) (y : Y) (n : ℕ),
            ProbabilityTheory.IdentDistrib (GJ19OptimalBinaryRatingSystems.sourceQuestionAskedIndicator E i y n)
              (GJ19OptimalBinaryRatingSystems.sourceQuestionAskedIndicator E i y 0) P P) →
          (∀ (i : Representative) (y : Y),
              ∫ (ω : Ω), GJ19OptimalBinaryRatingSystems.sourceQuestionAskedIndicator E i y 0 ω ∂P = H y) →
            (∀ (i : Representative) (y : Y),
                MeasureTheory.Integrable (GJ19OptimalBinaryRatingSystems.sourceQuestionPositiveIndicator E i y 0) P) →
              (∀ (i : Representative) (y : Y),
                  Pairwise
                    (Function.onFun (fun x1 x2 => ProbabilityTheory.IndepFun x1 x2 P)
                      (GJ19OptimalBinaryRatingSystems.sourceQuestionPositiveIndicator E i y))) →
                (∀ (i : Representative) (y : Y) (n : ℕ),
                    ProbabilityTheory.IdentDistrib
                      (GJ19OptimalBinaryRatingSystems.sourceQuestionPositiveIndicator E i y n)
                      (GJ19OptimalBinaryRatingSystems.sourceQuestionPositiveIndicator E i y 0) P P) →
                  (∀ (i : Representative) (y : Y),
                      ∫ (ω : Ω), GJ19OptimalBinaryRatingSystems.sourceQuestionPositiveIndicator E i y 0 ω ∂P =
                        H y * ψ (quality i) y) →
                    ∀ᵐ (ω : Ω) ∂P,
                      ∀ (i : Representative) (y : Y),
                        Filter.Tendsto
                          (fun N => GJ19OptimalBinaryRatingSystems.sourceExperimentEmpiricalQuestionResponse E i y N ω)
                          Filter.atTop (nhds (ψ (quality i) y))
\end{ReviewVerbatim}

\paragraph{Lean proof endpoint (built separately)}\begin{ReviewVerbatim}
GJ19OptimalBinaryRatingSystems.PaperInterface.lemmaB2_knownTypeExperiment_random_question_slln
\end{ReviewVerbatim}

\paragraph{Recorded source-to-Spec screening}
\textbf{Verdict:} \texttt{not recorded}
\\
\textbf{Reason:} not recorded
\reviewmatch{Matches source input}
\noindent\textbf{Reviewer annotation}\par
\reviewerbox
\clearpage
\hypertarget{source-claim-e38b113829a05c8d}{}
\section*{Review row 34}
\textbf{Source locator:} {\footnotesize\raggedright cited publication, lines 2032-\allowbreak{}2044\par}
\paragraph{Verbatim source input}
\begin{ReviewVerbatim}
Suppose instead the platform has many items L of unknown quality who are expected to match N
times each over the experiment time period. For each item, the platform would again ask questions
from Y, drawn according to any distribution (with positive mass on each question). Then generate
ψ̂(θ, y) as follows: first, rank the items according to their ratings during the experiment itself.
Then, for each y, ψ̂(θ, y) is
the empirical
performance of the θth percentile item in the ranking, i.e.
[U+0002 control character] i−1
[U+0003 control character]
i
ψ̂(θ, y) = ψ̂(θi , y) for θ ∈ L , L . Call this procedure UnknownTypeExperiment.
Lemma B.3. Suppose ψ(θ, y) is Lipschitz continuous in θ. With UnknownTypeExperiment, ψ̂(θ, y) →
ψ(θ, y)∀y, θ uniformly as L, N → ∞.
\end{ReviewVerbatim}


\paragraph{Expanded PaperInterface specification}
\begin{ReviewVerbatim}
∀ (Item : Type u_1) (Y : Type u_2) [inst : Fintype Item] [inst_1 : Fintype Y]
  (value : GJ19OptimalBinaryRatingSystems.SourceUnknownTypeExperiment Item Y),
  ∃ trueQuality empiricalResponse empiricalAverageScore rankedItem,
    value.trueQuality = trueQuality ∧
      value.empiricalResponse = empiricalResponse ∧
        value.empiricalAverageScore = empiricalAverageScore ∧ value.rankedItem = rankedItem
\end{ReviewVerbatim}

\paragraph{Lean proof endpoint (built separately)}\begin{ReviewVerbatim}
GJ19OptimalBinaryRatingSystems.PaperInterface.appendixB5_unknown_type_experiment
\end{ReviewVerbatim}

\paragraph{Recorded source-to-Spec screening}
\textbf{Verdict:} \texttt{not recorded}
\\
\textbf{Reason:} not recorded
\reviewmatch{Matches source input}
\noindent\textbf{Reviewer annotation}\par
\reviewerbox
\clearpage
\hypertarget{source-claim-045f11cd1b489b98}{}
\section*{Review row 35}
\textbf{Source locator:} {\footnotesize\raggedright cited publication, lines 2044-\allowbreak{}2060\par}
\paragraph{Verbatim source input}
\begin{ReviewVerbatim}
ψ(θ, y)∀y, θ uniformly as L, N → ∞.
Proof. Fix L. Denote each item in the experiment as i ∈ {1 . . . L} (with true quality θi 6= θj ),
and each item has N samples. Without loss of generality, assume the items are indexed according
to their rank on the average of their scores on the samples, defined as the percentage of positive
ratings received. i = 1 is then the worst item, and i = L is the best item according to scores in the
experiment.
For ψ(θ, x) increasing in θ, as N → ∞, P r(θi > θj |i < j) → 0 almost surely by SLLN, and for
a fixed L, {θi } this convergence[U+0002 control character]is uniform.
Furthermore, by SLLN, ψ̂(i, x) → ψ(θi , x) as N → ∞.
[U+0003 control character]
i−1 i
Recall ψ̂(θ, x) = ψ̂(i, x) for θ ∈ L , L .
10


[form-feed in source extraction]
Now, let L → ∞. ∀[U+000F control character], ∃L0 s.t. ∀L > L0 , ∀θ, ∃i s.t. |θ − θi | < [U+000F control character]. ψ(θ, x) is Lipschitz in θ by
assumption, and so ψ̂(θ, x) → ψ(θ, x) uniformly.
\end{ReviewVerbatim}


\paragraph{Expanded PaperInterface specification}
\begin{ReviewVerbatim}
∀ {Ω : Type u_1} {Y : Type u_2} [inst : MeasurableSpace Ω] [inst_1 : Fintype Y] [inst_2 : DecidableEq Y]
  (P : MeasureTheory.Measure Ω)
  (E : (L : ℕ) → GJ19OptimalBinaryRatingSystems.SourceRandomQuestionExperiment Ω (Fin (L + 1)) Y)
  (quality : (L : ℕ) → Fin (L + 1) → ℝ) (rankRepresentative : (L : ℕ) → ℝ → Fin (L + 1)) (psi : ℝ → Y → ℝ) (H : Y → ℝ),
  (∀ (y : Y), 0 < H y) →
    (∀ (L : ℕ) (i : Fin (L + 1)) (y : Y),
        MeasureTheory.Integrable (GJ19OptimalBinaryRatingSystems.sourceQuestionAskedIndicator (E L) i y 0) P) →
      (∀ (L : ℕ) (i : Fin (L + 1)) (y : Y),
          Pairwise
            (Function.onFun (fun x1 x2 => ProbabilityTheory.IndepFun x1 x2 P)
              (GJ19OptimalBinaryRatingSystems.sourceQuestionAskedIndicator (E L) i y))) →
        (∀ (L : ℕ) (i : Fin (L + 1)) (y : Y) (n : ℕ),
            ProbabilityTheory.IdentDistrib (GJ19OptimalBinaryRatingSystems.sourceQuestionAskedIndicator (E L) i y n)
              (GJ19OptimalBinaryRatingSystems.sourceQuestionAskedIndicator (E L) i y 0) P P) →
          (∀ (L : ℕ) (i : Fin (L + 1)) (y : Y),
              ∫ (ω : Ω), GJ19OptimalBinaryRatingSystems.sourceQuestionAskedIndicator (E L) i y 0 ω ∂P = H y) →
            (∀ (L : ℕ) (i : Fin (L + 1)) (y : Y),
                MeasureTheory.Integrable (GJ19OptimalBinaryRatingSystems.sourceQuestionPositiveIndicator (E L) i y 0)
                  P) →
              (∀ (L : ℕ) (i : Fin (L + 1)) (y : Y),
                  Pairwise
                    (Function.onFun (fun x1 x2 => ProbabilityTheory.IndepFun x1 x2 P)
                      (GJ19OptimalBinaryRatingSystems.sourceQuestionPositiveIndicator (E L) i y))) →
                (∀ (L : ℕ) (i : Fin (L + 1)) (y : Y) (n : ℕ),
                    ProbabilityTheory.IdentDistrib
                      (GJ19OptimalBinaryRatingSystems.sourceQuestionPositiveIndicator (E L) i y n)
                      (GJ19OptimalBinaryRatingSystems.sourceQuestionPositiveIndicator (E L) i y 0) P P) →
                  (∀ (L : ℕ) (i : Fin (L + 1)) (y : Y),
                      ∫ (ω : Ω), GJ19OptimalBinaryRatingSystems.sourceQuestionPositiveIndicator (E L) i y 0 ω ∂P =
                        H y * psi (quality L i) y) →
                    (∀ (L : ℕ) (i : Fin (L + 1)),
                        MeasureTheory.Integrable
                          (GJ19OptimalBinaryRatingSystems.sourceAggregatePositiveIndicator (E L) i 0) P) →
                      (∀ (L : ℕ) (i : Fin (L + 1)),
                          Pairwise
                            (Function.onFun (fun x1 x2 => ProbabilityTheory.IndepFun x1 x2 P)
                              (GJ19OptimalBinaryRatingSystems.sourceAggregatePositiveIndicator (E L) i))) →
                        (∀ (L : ℕ) (i : Fin (L + 1)) (n : ℕ),
                            ProbabilityTheory.IdentDistrib
                              (GJ19OptimalBinaryRatingSystems.sourceAggregatePositiveIndicator (E L) i n)
                              (GJ19OptimalBinaryRatingSystems.sourceAggregatePositiveIndicator (E L) i 0) P P) →
                          (∀ (L : ℕ) (i : Fin (L + 1)),
                              ∫ (ω : Ω),
                                  GJ19OptimalBinaryRatingSystems.sourceAggregatePositiveIndicator (E L) i 0 ω ∂P =
                                GJ19OptimalBinaryRatingSystems.sourceExpectedAggregateScore psi H (quality L i)) →
                            (∀ (L : ℕ),
                                StrictMono fun i =>
                                  GJ19OptimalBinaryRatingSystems.sourceExpectedAggregateScore psi H (quality L i)) →
                              (∀ δ > 0,
                                  ∀ᶠ (L : ℕ) in Filter.atTop,
                                    ∀ θ ∈ Set.Icc 0 1, dist (quality L (rankRepresentative L θ)) θ < δ) →
                                ∀ {K : ℝ},
                                  0 ≤ K →
                                    (∀ (θ θ' : ℝ) (y : Y), dist (psi θ y) (psi θ' y) ≤ K * dist θ θ') →
                                      ∀ᵐ (ω : Ω) ∂P,
                                        ∀ ε > 0,
                                          ∀ᶠ (L : ℕ) (N : ℕ) in Filter.atTop,
                                            (StrictMono fun i =>
                                                GJ19OptimalBinaryRatingSystems.sourceExperimentEmpiricalMean
                                                  (GJ19OptimalBinaryRatingSystems.sourceAggregatePositiveIndicator (E L)
                                                    i)
                                                  N ω) ∧
                                              ∀ θ ∈ Set.Icc 0 1,
                                                ∀ (y : Y),
                                                  dist
                                                      (GJ19OptimalBinaryRatingSystems.sourceExperimentEmpiricalQuestionResponse
                                                        (E L) (rankRepresentative L θ) y N ω)
                                                      (psi θ y) <
                                                    ε
\end{ReviewVerbatim}

\paragraph{Lean proof endpoint (built separately)}\begin{ReviewVerbatim}
GJ19OptimalBinaryRatingSystems.PaperInterface.paper_lemmaB3_unknownTypeExperiment_uniform_convergence_of_ranking_tracking
\end{ReviewVerbatim}

\paragraph{Recorded source-to-Spec screening}
\textbf{Verdict:} \texttt{not recorded}
\\
\textbf{Reason:} not recorded
\reviewmatch{Matches source input}
\noindent\textbf{Reviewer annotation}\par
\reviewerbox
\clearpage
\hypertarget{source-claim-99b64aa08841b772}{}
\section*{Review row 36}
\textbf{Source locator:} {\footnotesize\raggedright cited publication, lines 2032-\allowbreak{}2060\par}
\paragraph{Verbatim source input}
\begin{ReviewVerbatim}
Suppose instead the platform has many items L of unknown quality who are expected to match N
times each over the experiment time period. For each item, the platform would again ask questions
from Y, drawn according to any distribution (with positive mass on each question). Then generate
ψ̂(θ, y) as follows: first, rank the items according to their ratings during the experiment itself.
Then, for each y, ψ̂(θ, y) is
the empirical
performance of the θth percentile item in the ranking, i.e.
[U+0002 control character] i−1
[U+0003 control character]
i
ψ̂(θ, y) = ψ̂(θi , y) for θ ∈ L , L . Call this procedure UnknownTypeExperiment.
Lemma B.3. Suppose ψ(θ, y) is Lipschitz continuous in θ. With UnknownTypeExperiment, ψ̂(θ, y) →
ψ(θ, y)∀y, θ uniformly as L, N → ∞.
Proof. Fix L. Denote each item in the experiment as i ∈ {1 . . . L} (with true quality θi 6= θj ),
and each item has N samples. Without loss of generality, assume the items are indexed according
to their rank on the average of their scores on the samples, defined as the percentage of positive
ratings received. i = 1 is then the worst item, and i = L is the best item according to scores in the
experiment.
For ψ(θ, x) increasing in θ, as N → ∞, P r(θi > θj |i < j) → 0 almost surely by SLLN, and for
a fixed L, {θi } this convergence[U+0002 control character]is uniform.
Furthermore, by SLLN, ψ̂(i, x) → ψ(θi , x) as N → ∞.
[U+0003 control character]
i−1 i
Recall ψ̂(θ, x) = ψ̂(i, x) for θ ∈ L , L .
10


[form-feed in source extraction]
Now, let L → ∞. ∀[U+000F control character], ∃L0 s.t. ∀L > L0 , ∀θ, ∃i s.t. |θ − θi | < [U+000F control character]. ψ(θ, x) is Lipschitz in θ by
assumption, and so ψ̂(θ, x) → ψ(θ, x) uniformly.
\end{ReviewVerbatim}


\paragraph{Expanded PaperInterface specification}
\begin{ReviewVerbatim}
∀ {Ω : Type u_1} {Y : Type u_2} {L : ℕ} [inst : MeasurableSpace Ω] [inst_1 : Fintype Y] [inst_2 : DecidableEq Y]
  (P : MeasureTheory.Measure Ω) (E : GJ19OptimalBinaryRatingSystems.SourceRandomQuestionExperiment Ω (Fin L) Y)
  (quality : Fin L → ℝ) (ψ : ℝ → Y → ℝ) (H : Y → ℝ),
  (∀ (y : Y), 0 < H y) →
    (∀ (i : Fin L) (y : Y),
        MeasureTheory.Integrable (GJ19OptimalBinaryRatingSystems.sourceQuestionAskedIndicator E i y 0) P) →
      (∀ (i : Fin L) (y : Y),
          Pairwise
            (Function.onFun (fun x1 x2 => ProbabilityTheory.IndepFun x1 x2 P)
              (GJ19OptimalBinaryRatingSystems.sourceQuestionAskedIndicator E i y))) →
        (∀ (i : Fin L) (y : Y) (n : ℕ),
            ProbabilityTheory.IdentDistrib (GJ19OptimalBinaryRatingSystems.sourceQuestionAskedIndicator E i y n)
              (GJ19OptimalBinaryRatingSystems.sourceQuestionAskedIndicator E i y 0) P P) →
          (∀ (i : Fin L) (y : Y),
              ∫ (ω : Ω), GJ19OptimalBinaryRatingSystems.sourceQuestionAskedIndicator E i y 0 ω ∂P = H y) →
            (∀ (i : Fin L) (y : Y),
                MeasureTheory.Integrable (GJ19OptimalBinaryRatingSystems.sourceQuestionPositiveIndicator E i y 0) P) →
              (∀ (i : Fin L) (y : Y),
                  Pairwise
                    (Function.onFun (fun x1 x2 => ProbabilityTheory.IndepFun x1 x2 P)
                      (GJ19OptimalBinaryRatingSystems.sourceQuestionPositiveIndicator E i y))) →
                (∀ (i : Fin L) (y : Y) (n : ℕ),
                    ProbabilityTheory.IdentDistrib
                      (GJ19OptimalBinaryRatingSystems.sourceQuestionPositiveIndicator E i y n)
                      (GJ19OptimalBinaryRatingSystems.sourceQuestionPositiveIndicator E i y 0) P P) →
                  (∀ (i : Fin L) (y : Y),
                      ∫ (ω : Ω), GJ19OptimalBinaryRatingSystems.sourceQuestionPositiveIndicator E i y 0 ω ∂P =
                        H y * ψ (quality i) y) →
                    (∀ (i : Fin L),
                        MeasureTheory.Integrable (GJ19OptimalBinaryRatingSystems.sourceAggregatePositiveIndicator E i 0)
                          P) →
                      (∀ (i : Fin L),
                          Pairwise
                            (Function.onFun (fun x1 x2 => ProbabilityTheory.IndepFun x1 x2 P)
                              (GJ19OptimalBinaryRatingSystems.sourceAggregatePositiveIndicator E i))) →
                        (∀ (i : Fin L) (n : ℕ),
                            ProbabilityTheory.IdentDistrib
                              (GJ19OptimalBinaryRatingSystems.sourceAggregatePositiveIndicator E i n)
                              (GJ19OptimalBinaryRatingSystems.sourceAggregatePositiveIndicator E i 0) P P) →
                          (∀ (i : Fin L),
                              ∫ (ω : Ω), GJ19OptimalBinaryRatingSystems.sourceAggregatePositiveIndicator E i 0 ω ∂P =
                                GJ19OptimalBinaryRatingSystems.sourceExpectedAggregateScore ψ H (quality i)) →
                            (StrictMono fun i =>
                                GJ19OptimalBinaryRatingSystems.sourceExpectedAggregateScore ψ H (quality i)) →
                              ∀ᵐ (ω : Ω) ∂P,
                                TendstoUniformlyOn
                                    (fun N p =>
                                      GJ19OptimalBinaryRatingSystems.sourceExperimentEmpiricalQuestionResponse E p.1 p.2
                                        N ω)
                                    (fun p => ψ (quality p.1) p.2) Filter.atTop Set.univ ∧
                                  ∀ᶠ (N : ℕ) in Filter.atTop,
                                    StrictMono fun i =>
                                      GJ19OptimalBinaryRatingSystems.sourceExperimentEmpiricalMean
                                        (GJ19OptimalBinaryRatingSystems.sourceAggregatePositiveIndicator E i) N ω
\end{ReviewVerbatim}

\paragraph{Lean proof endpoint (built separately)}\begin{ReviewVerbatim}
GJ19OptimalBinaryRatingSystems.PaperInterface.lemmaB3_unknownTypeExperiment_random_question_response_and_rank
\end{ReviewVerbatim}

\paragraph{Recorded source-to-Spec screening}
\textbf{Verdict:} \texttt{not recorded}
\\
\textbf{Reason:} not recorded
\reviewmatch{Matches source input}
\noindent\textbf{Reviewer annotation}\par
\reviewerbox
\clearpage
\hypertarget{source-claim-7ea3c8ab498f92dc}{}
\section*{Review row 37}
\textbf{Source locator:} {\footnotesize\raggedright cited publication, lines 2083-\allowbreak{}2131\par}
\paragraph{Verbatim source input}
\begin{ReviewVerbatim}
Rate functions for Pk (θ1 , θ2 )

Lemma C.1.
1
lim − log [µ((xk (θ1 ) − xk (θ2 )) ≤ 0|θ1 , θ2 )] = inf {g(θ1 )I(a|θ1 ) + g(θ2 )I(a|θ2 )}
a∈R
k→∞ k
where I(a|`) = supz {za − Λ(z|θ)}, and Λ(z|θ) is the log moment generating function of a single
sample from x(θ1 ) and g(θ) is the sampling rate.
Proof. limk→∞ − k1 log [µ((xk (θ1 ) − xk (θ2 )) ≤ 0|θ1 , θ2 )]
[U+0015 control character]
[U+0014 control character]Z
1
µ((xk (θ1 ) = a|θ1 )µ(xk (θ2 ) ≥ a|θ2 )da
= lim − log
k→∞ k
[U+0014 control character]Za∈R
[U+0015 control character]
1
−kg(θ1 )I(a|θ1 ) −kg(θ2 )I(a|θ2 )
= lim − log
e
e
da
k→∞ k
a∈R
= inf {g(θ1 )I(a|θ1 ) + g(θ2 )I(a|θ2 )}

(1)
(2)
Laplace principle

a∈R

(3)

Where (2) is a basic result from large deviations, where kg(θi ) is the number of samples item
of quality θi has received.
Note that this lemma also appears in Glynn and Juneja (2004), which uses the Gartner-Ellis
Theorem in the proof. Our proof is conceptually similar but instead uses Laplace’s principle.
1−b
Recall that KL(a||b) = a log ab + (1 − a) log 1−a
is the Kullback-Leibler (KL) divergence between
Bernoulli random variables with success probabilities a and b respectively. It is well known that
for a Bernoulli random variable with success probability t,
I(a|t) = KL(a||t)
11
\end{ReviewVerbatim}


\paragraph{Expanded PaperInterface specification}
\begin{ReviewVerbatim}
∀ {Seller : Type u_1} {Pair : Type u_2} (successProb : Seller → ℝ) (hprob0 : ∀ (θ : Seller), 0 ≤ successProb θ)
  (hprob1 : ∀ (θ : Seller), successProb θ ≤ 1),
  (∀ (θ : Seller), 0 < successProb θ) →
    (∀ (θ : Seller), successProb θ < 1) →
      ∀ (sampleRate : Seller → ℝ) (pairHi pairLo : Pair → Seller),
        (∀ (p : Pair), 0 < sampleRate (pairHi p)) →
          (∀ (p : Pair), 0 < sampleRate (pairLo p)) →
            ∀ (a z : Pair → ℝ),
              (∀ (p : Pair), z p ≤ 0) →
                (∀ (p : Pair),
                    HasDerivAt
                      (fun t =>
                        (GJ19OptimalBinaryRatingSystems.binaryRatingModel successProb hprob0 hprob1).logMGF (pairHi p)
                          t)
                      (a p) (z p * (sampleRate (pairHi p))⁻¹)) →
                  (∀ (p : Pair),
                      HasDerivAt
                        (fun t =>
                          (GJ19OptimalBinaryRatingSystems.binaryRatingModel successProb hprob0 hprob1).logMGF (pairLo p)
                            t)
                        (a p) (-(z p * (sampleRate (pairLo p))⁻¹))) →
                    Nonempty
                      (EconCSLib.Probability.PairwiseThresholdRateTopLdpCertificate
                        (GJ19OptimalBinaryRatingSystems.binaryRatingModel successProb hprob0 hprob1) sampleRate pairHi
                        pairLo)
\end{ReviewVerbatim}

\paragraph{Lean proof endpoint (built separately)}\begin{ReviewVerbatim}
GJ19OptimalBinaryRatingSystems.PaperInterface.source_lemmaC1_pairwise_error_rate_from_derivatives
\end{ReviewVerbatim}

\paragraph{Recorded source-to-Spec screening}
\textbf{Verdict:} \texttt{not recorded}
\\
\textbf{Reason:} not recorded
\reviewmatch{Matches source input}
\noindent\textbf{Reviewer annotation}\par
\reviewerbox
\clearpage
\hypertarget{source-claim-5122f84bd17b3a2c}{}
\section*{Review row 38}
\textbf{Source locator:} {\footnotesize\raggedright cited publication, lines 2133-\allowbreak{}2161\par}
\paragraph{Verbatim source input}
\begin{ReviewVerbatim}
Lemma C.2. Let θ1 > θ2 and I(a|θ) = KL(a||β(θ)). Further, Let P k (θ1 , θ2 ) = 1 − Pk (θ1 , θ2 ).
Then,
1
− lim log P k (θ1 , θ2 ) = inf {g(θ1 )I(a|θ1 ) + g(θ2 )I(a|θ2 )} ,
(4)
a∈R
k→∞ k
Proof. Follows directly from Lemma C.1.
− lim

1

k→∞ k

log P k (θ1 , θ2 |β)

1
= lim − log [1 + µk (xk (θ1 ) − xk (θ2 ) < 0|θ1 , θ2 ) − µk (xk (θ1 ) − xk (θ2 ) > 0|θ1 , θ2 )]
k→∞ k
1
= lim − log [2µk (xk (θ1 ) − xk (θ2 ) < 0|θ1 , θ2 ) + µk (xk (θ1 ) − xk (θ2 ) = 0|θ1 , θ2 )]
k→∞ k
1
= lim − log [µk (xk (θ1 ) − xk (θ2 ) ≤ 0|θ1 , θ2 )]
k→∞ k
= inf {g(θ1 )I(a|θ1 ) + g(θ2 )I(a|θ2 )}
Lemma C.1
a∈R
\end{ReviewVerbatim}


\paragraph{Expanded PaperInterface specification}
\begin{ReviewVerbatim}
∀ {Seller : Type u_1} (successProb : Seller → ℝ) (hprob0 : ∀ (θ : Seller), 0 ≤ successProb θ)
  (hprob1 : ∀ (θ : Seller), successProb θ ≤ 1) (sampleRate : Seller → ℝ) (hi lo : Seller),
  0 < sampleRate hi →
    0 < sampleRate lo →
      0 < successProb hi →
        successProb hi < 1 →
          0 < successProb lo →
            successProb lo < 1 →
              successProb lo ≤ successProb hi →
                EconCSLib.Probability.ExponentialRateCertificate
                  (EconCSLib.Probability.twoSampleFloorPkComplementErrorProb
                    (GJ19OptimalBinaryRatingSystems.binaryRatingModel successProb hprob0 hprob1) sampleRate hi lo)
                  (EconCSLib.Probability.weightedBernoulliClosedThresholdRate (sampleRate hi) (sampleRate lo)
                    (successProb hi) (successProb lo))
\end{ReviewVerbatim}

\paragraph{Lean proof endpoint (built separately)}\begin{ReviewVerbatim}
GJ19OptimalBinaryRatingSystems.PaperInterface.source_lemmaC2_binary_complement_rate
\end{ReviewVerbatim}

\paragraph{Recorded source-to-Spec screening}
\textbf{Verdict:} \texttt{not recorded}
\\
\textbf{Reason:} not recorded
\reviewmatch{Matches source input}
\noindent\textbf{Reviewer annotation}\par
\reviewerbox
\clearpage
\hypertarget{source-claim-d95b8ab740a8ceda}{}
\section*{Review row 39}
\textbf{Source locator:} {\footnotesize\raggedright cited publication, lines 2163-\allowbreak{}2223\par}
\paragraph{Verbatim source input}
\begin{ReviewVerbatim}
Laplace’s principle with sequence of rate functions

In order to derive a rate function for W k = (limk Wk ) − Wk , we need to be able to relate its rate to
that of P k (θ1 , θ2 ). The following theorem, related to Laplace’s principle for large deviations allows
us to do so.
Theorem C.1. Suppose that X is compact with finite Lebesgue measure µ(X) < ∞. Suppose that
ϕ(x) has an essential infimum ϕ on X, that ϕn (x) has an essential infimum ϕn , that both ϕ and
all ϕn are nonnegative, and that ϕn → ϕ uniformly:
lim sup |ϕn (x) − ϕ(x)| = 0.

n→∞ x∈X

Then:

1
log
n→∞ n

Z

lim

e−nϕn (x) dx = −ϕ.

(5)

X

Proof. First, we note that for all x and n, e−nϕn (x) ≤ e−nϕn . Therefore, letting (∗) denote the
LHS of (5), we have:
Z
1
(∗) ≤ lim log
e−nϕn dx = −ϕ,
n→∞ n
X
where the last limit follows from the fact that ϕn converges uniformly to ϕ, so that ϕn → ϕ.
Next, for [U+000F control character] > 0 let An ([U+000F control character]) = {x : ϕn (x) ≤ ϕn + [U+000F control character]} and let A([U+000F control character]) = {x : ϕ(x) ≤ ϕ + [U+000F control character]}. It
follows (again by uniform convergence) that for all sufficiently large n, A([U+000F control character]/2) ⊆ An ([U+000F control character]), so that
µ(A([U+000F control character]/2)) ≤ µ(An ([U+000F control character])) for all sufficiently large n. Further, µ(A([U+000F control character]/2)) > 0, since ϕ is the essential
infimum.
Since:
Z
e−nϕn (x) dx ≥ µ(An ([U+000F control character]))e−n(ϕn +[U+000F control character]) ,
X

12


[form-feed in source extraction]
it follows that:

1
log µ(An ([U+000F control character])).
n→∞ n
To complete the proof, observe that since µ(An ([U+000F control character])) is bounded below by a positive constant for all
sufficiently large n, the last limit is zero. Therefore:
(∗) ≥ −ϕ − [U+000F control character] + lim

(∗) ≥ −ϕ − [U+000F control character].
Since [U+000F control character] was arbitrary, this completes the proof.
Remark C.1. Let X = [0, 1] × [0, 1], ϕn (θ1 , θ2 ) = − n1 log P n (θ1 , θ2 ). Then, all the conditions for
\end{ReviewVerbatim}


\paragraph{Expanded PaperInterface specification}
\begin{ReviewVerbatim}
∀ {Ω : Type u_1} [inst : MeasurableSpace Ω] (μ : MeasureTheory.Measure Ω) [MeasureTheory.IsFiniteMeasure μ]
  (phiSeq : ℕ → Ω → ℝ) (phi : Ω → ℝ) {rate : ℝ},
  (∀ (k : ℕ), MeasureTheory.Integrable (fun x => Real.exp (-↑k * phiSeq k x)) μ) →
    EconCSLib.Probability.HasAEEssentialInfimum μ phi rate →
      (∀ ε > 0, ∀ᶠ (k : ℕ) in Filter.atTop, ∀ (x : Ω), |phiSeq k x - phi x| ≤ ε) →
        EconCSLib.Probability.HasExponentialRate (fun k => ∫ (x : Ω), Real.exp (-↑k * phiSeq k x) ∂μ) rate
\end{ReviewVerbatim}

\paragraph{Lean proof endpoint (built separately)}\begin{ReviewVerbatim}
GJ19OptimalBinaryRatingSystems.PaperInterface.source_theoremC1_laplace_principle
\end{ReviewVerbatim}

\paragraph{Recorded source-to-Spec screening}
\textbf{Verdict:} \texttt{not recorded}
\\
\textbf{Reason:} not recorded
\reviewmatch{Matches source input}
\noindent\textbf{Reviewer annotation}\par
\reviewerbox
\clearpage
\hypertarget{source-claim-16ee341855124303}{}
\section*{Review row 40}
\textbf{Source locator:} {\footnotesize\raggedright cited publication, lines 2224-\allowbreak{}2225\par}
\paragraph{Verbatim source input}
\begin{ReviewVerbatim}
Theorem C.1 are met.
\end{ReviewVerbatim}


\paragraph{Expanded PaperInterface specification}
\begin{ReviewVerbatim}
∀ (weight : ℝ × ℝ → ℝ) (phiSeq : ℕ → ℝ × ℝ → ℝ) (phi : ℝ × ℝ → ℝ) {rate W : ℝ},
  (∀ (k : ℕ),
      MeasureTheory.Integrable (fun x => weight x * Real.exp (-↑k * phiSeq k x))
        ((MeasureTheory.volume.restrict (Set.Icc 0 1)).prod (MeasureTheory.volume.restrict (Set.Icc 0 1)))) →
    0 < W →
      (∀ᵐ (x : ℝ × ℝ) ∂(MeasureTheory.volume.restrict (Set.Icc 0 1)).prod (MeasureTheory.volume.restrict (Set.Icc 0 1)),
          0 ≤ weight x) →
        (∀ᵐ (x :
            ℝ × ℝ) ∂(MeasureTheory.volume.restrict (Set.Icc 0 1)).prod (MeasureTheory.volume.restrict (Set.Icc 0 1)),
            weight x ≤ W) →
          EconCSLib.Probability.HasAEEssentialInfimum
              ((MeasureTheory.volume.restrict (Set.Icc 0 1)).prod (MeasureTheory.volume.restrict (Set.Icc 0 1))) phi
              rate →
            EconCSLib.Probability.HasPositiveWeightNearAEEssentialInfimum
                ((MeasureTheory.volume.restrict (Set.Icc 0 1)).prod (MeasureTheory.volume.restrict (Set.Icc 0 1)))
                weight phi rate →
              (∀ ε > 0, ∀ᶠ (k : ℕ) in Filter.atTop, ∀ (x : ℝ × ℝ), |phiSeq k x - phi x| ≤ ε) →
                EconCSLib.Probability.HasExponentialRate
                  (fun k =>
                    ∫ (x : ℝ × ℝ),
                      weight x *
                        Real.exp
                          (-↑k *
                            phiSeq k
                              x) ∂(MeasureTheory.volume.restrict (Set.Icc 0 1)).prod
                        (MeasureTheory.volume.restrict (Set.Icc 0 1)))
                  rate
\end{ReviewVerbatim}

\paragraph{Lean proof endpoint (built separately)}\begin{ReviewVerbatim}
GJ19OptimalBinaryRatingSystems.PaperInterface.source_remarkC1_full_square_application_under_explicit_conditions
\end{ReviewVerbatim}

\paragraph{Recorded source-to-Spec screening}
\textbf{Verdict:} \texttt{not recorded}
\\
\textbf{Reason:} not recorded
\reviewmatch{Matches source input}
\noindent\textbf{Reviewer annotation}\par
\reviewerbox
\clearpage
\hypertarget{source-claim-7529d1dedcffdbe3}{}
\section*{Review row 41}
\textbf{Source locator:} {\footnotesize\raggedright cited publication, lines 2227-\allowbreak{}2364\par}
\paragraph{Verbatim source input}
\begin{ReviewVerbatim}
Rate function for Wk

Our next lemma shows that we can obtain a nontrivial large deviations rate for Wk when β is a
step-wise increasing
R function.
Recall Wk = θ1 >θ2 w(θ1 , θ2 )Pk (θ1 , θ2 |β)d(θ1 , θ2 ).
Let P k (θ1 , θ2 ) = 1 − Pk (θ1 , θ2 ).
R
Further, let W k = (limk Wk ) − Wk = θ1 >θ2 w(θ1 , θ2 )P k (θ1 , θ2 |β)d(θ1 , θ2 ). (recall we assumed
w integrates to 1 without loss of generality).
Lemma C.3. Suppose β is piecewise constant with M levels {ti }. Let gi , inf θ∈Si g(θ) = g(si )
Then,
1
− lim log W k = min inf {gi+1 I(a|ti+1 ) + gi I(a|ti )} , r,
(6)
0≤i≤M −2 a∈R
k→∞ k
where I(a|t) = KL(a||t) as defined in Lemma C.2.
Proof. When β is step-wise increasing with M levels {ti }, then
X Z
Wk =
w(θ1 , θ2 )P k (θ1 , θ2 |β)d(θ1 , θ2 )
0≤i<j<M

θ2 ∈Si ,θ1 ∈Sj

as P k (θ1 , θ2 ) = 0 when β(θ1 ) = β(θ2 ).

13


[form-feed in source extraction]
− limk→∞ k1 log W k
1
= − lim log
k→∞ k
1
= − lim log
k→∞ k

Z
w(θ1 , θ2 )P k (θ1 , θ2 |β)d(θ1 , θ2 )
X Z
w(θ1 , θ2 )P k (θ1 , θ2 |β)d(θ1 , θ2 )

0≤i<j<M

=
=

min

inf

(8)

θ2 ∈Si ,θ1 ∈Sj

" Z
#!
1
= − max
lim
log
w(θ1 , θ2 )P k (θj , θi |β)d(θ1 , θ2 )
0≤i<j<M k→∞ k
θ2 ∈Si ,θ1 ∈Sj
[U+0013 control character]
[U+0012 control character]
[U+0003 control character]
1[U+0002 control character]
= − max
sup
log w(θ1 , θ2 )P k (θj , θi |β)
lim
0≤i<j<M θ1 ∈Sj ,θ2 ∈Si k→∞ k
[U+0012 control character]
[U+0013 control character]
1
= − max
sup
lim log P k (θj , θi |β)
0≤i<j<M θ1 ∈Sj ,θ2 ∈Si k→∞ k
[U+0012 control character]
[U+0013 control character]
1
− lim log P k (θj , θi |β)
= min
inf
0≤i<j<M θ1 ∈Sj ,θ2 ∈Si
k→∞ k
=

(7)

θ1 >θ2

inf {g(θ1 )I(a|tj ) + g(θ2 )I(a|ti )}

0≤i<j<M θ1 ∈Sj ,θ2 ∈Si a∈R

(9)
(10)
(11)
(12)
(13)

min

inf {gj I(a|tj ) + gi I(a|ti )}

(14)

min

inf {gi+1 I(a|ti+1 ) + gi I(a|ti )}

(15)

0≤i<j<M a∈R

0≤i<M −1 a∈R

The last line follows from adjacent ti , ti+1 dominating the rate due to monotonicity properties.
Line (10) follows from Theorem C.1.
[U+0011 control character]i
h
[U+0010 control character]P
N [U+000F control character]
[U+000F control character]
= maxN
Line (9) follows from: ∀a[U+000F control character]i ≥ 0, lim sup[U+000F control character]→0 [U+000F control character] log
a
i lim sup[U+000F control character]→0 [U+000F control character] log(ai ), which
i
i
is a finite case version (with fewer assumptions) of Theorem C.1. See, e.g., Lemma 1.2.15 in (Dembo
and Zeitouni, 2010) for a proof of this property.
\end{ReviewVerbatim}


\paragraph{Expanded PaperInterface specification}
\begin{ReviewVerbatim}
∀ {Ω : Type u_1} {Component : Type u_2} [inst : MeasurableSpace Ω] [inst_1 : Fintype Component] [DecidableEq Component]
  [Nonempty Component] (μ : MeasureTheory.Measure Ω) [MeasureTheory.IsFiniteMeasure μ]
  (P : EconCSLib.Probability.FiniteMeasurableSetPartition μ Component) (weight : Ω → ℝ) (kernel : ℕ → Ω → ℝ)
  (phi : Component → Ω → ℝ) (rate W : Component → ℝ),
  (∀ (k : ℕ) (component : Component),
      MeasureTheory.IntegrableOn (fun x => weight x * kernel k x) (P.pieceSet component) μ) →
    (∀ (component : Component), MeasureTheory.IntegrableOn weight (P.pieceSet component) μ) →
      (∀ (component : Component), 0 < W component) →
        (∀ (component : Component), ∀ᵐ (x : Ω) ∂μ.restrict (P.pieceSet component), 0 ≤ weight x) →
          (∀ (component : Component), ∀ᵐ (x : Ω) ∂μ.restrict (P.pieceSet component), weight x ≤ W component) →
            (∀ (component : Component),
                EconCSLib.Probability.HasAEEssentialInfimum (μ.restrict (P.pieceSet component)) (phi component)
                  (rate component)) →
              (∀ (component : Component),
                  EconCSLib.Probability.HasPositiveWeightNearAEEssentialInfimum (μ.restrict (P.pieceSet component))
                    weight (phi component) (rate component)) →
                (∀ (k : ℕ) (x : Ω), 0 < kernel k x) →
                  (∀ (component : Component),
                      ∀ ε > 0,
                        ∀ᶠ (k : ℕ) in Filter.atTop, ∀ (x : Ω), |-Real.log (kernel k x) / ↑k - phi component x| ≤ ε) →
                    ∃ minComponent,
                      (∀ (component : Component), rate minComponent ≤ rate component) ∧
                        EconCSLib.Probability.HasExponentialRate
                          (fun k => ∫ (x : Ω) in P.support, weight x * kernel k x ∂μ) (rate minComponent)
\end{ReviewVerbatim}

\paragraph{Lean proof endpoint (built separately)}\begin{ReviewVerbatim}
GJ19OptimalBinaryRatingSystems.PaperInterface.paper_lemmaC3_partition_integral_hasExponentialRate_of_min_component_uniform_logRate
\end{ReviewVerbatim}

\paragraph{Recorded source-to-Spec screening}
\textbf{Verdict:} \texttt{not recorded}
\\
\textbf{Reason:} not recorded
\reviewmatch{Matches source input}
\noindent\textbf{Reviewer annotation}\par
\reviewerbox
\clearpage
\hypertarget{source-claim-043c064a62132f60}{}
\section*{Review row 42}
\textbf{Source locator:} {\footnotesize\raggedright cited publication, lines 2366-\allowbreak{}2399\par}
\paragraph{Verbatim source input}
\begin{ReviewVerbatim}
Proof. =⇒ follows directly from Lemma C.3: inf a∈R {gi+1 I(a|ti+1 ) + gi I(a|ti )} > 0 when ti 6= ti+1 ,
which holds when β is piece-wise constant with the appropriate number of levels.
⇐= Consider β that is not piece-wise constant. Recall that we further assume that β is nondecreasing, and discontinuous only on a measure 0 set. Following algebra steps similar to those in
Lemma C.3, but for general β:
Z
1
1
log W k = − lim log
w(θ1 , θ2 )P k (θ1 , θ2 |β)d(θ1 , θ2 )
k→∞ k
k→∞ k
θ1 >θ2
[U+0012 control character]
[U+0013 control character]
1
= inf − lim log P k (θj , θi |β)
θ1 >θ2
k→∞ k
=0

− lim

(16)
(17)
(18)

Where the last line follows from β continuous at some θ1 , and so limθ2 →θ1 P k (θ1 , θ2 |β) = 1.
Intuitively, what goes wrong with continuous β is that P k (θ1 , θ2 |β) does not converge uniformly:
∀[U+000F control character], k, ∃θ2 6= θ1 P k (θ1 , θ2 ) > [U+000F control character]
14


[form-feed in source extraction]
i.e. close by items are very hard to distinguish from one another. Then, because the large deviations
rate of W k is dominated by the worst rates under the integral, we don’t get a positive rate.
\end{ReviewVerbatim}


\paragraph{Expanded PaperInterface specification}
\begin{ReviewVerbatim}
∀ (μ : MeasureTheory.Measure ℝ) [MeasureTheory.SFinite μ] [MeasureTheory.IsFiniteMeasure (μ.prod μ)]
  [(μ.prod μ).IsOpenPosMeasure] {m : ℕ},
  0 < m →
    ∀ (cut : ℕ → ℝ) (hmono : Monotone cut),
      (∀ i < m + 2, cut i < cut (i + 1)) →
        ∀ (levels sampleRate : Fin (m + 2) → ℝ)
          (hlevels : GJ19OptimalBinaryRatingSystems.BinaryEndpointLevelVector levels),
          (∀ (idx : Fin (m + 2)), 0 < sampleRate idx) →
            (∀ {a b : Fin (m + 2)}, ↑a ≤ ↑b → sampleRate a ≤ sampleRate b) →
              ∀ (weight : ℝ × ℝ → ℝ),
                (∀
                    (component :
                      { piece // GJ19OptimalBinaryRatingSystems.theorem31OrderedNontrivialPairSelected piece }),
                    MeasureTheory.IntegrableOn weight
                      ((GJ19OptimalBinaryRatingSystems.theorem31_ordered_quality_pair_partition μ (m + 2) cut hmono
                            GJ19OptimalBinaryRatingSystems.theorem31OrderedNontrivialPairSelected).pieceSet
                        component)
                      (μ.prod μ)) →
                  (∀
                      (component :
                        { piece // GJ19OptimalBinaryRatingSystems.theorem31OrderedNontrivialPairSelected piece }),
                      ∀ᵐ (x :
                        ℝ ×
                          ℝ) ∂(μ.prod μ).restrict
                          ((GJ19OptimalBinaryRatingSystems.theorem31_ordered_quality_pair_partition μ (m + 2) cut hmono
                                GJ19OptimalBinaryRatingSystems.theorem31OrderedNontrivialPairSelected).pieceSet
                            component),
                        0 ≤ weight x) →
                    Continuous weight →
                      (∀ (component : GJ19OptimalBinaryRatingSystems.theorem31OrderedNontrivialPairComponent m),
                          0 <
                            weight
                              (GJ19OptimalBinaryRatingSystems.theorem31_ordered_quality_pair_component_midpoint cut
                                component)) →
                        ∀ (lo hi : ℝ)
                          (R : GJ19OptimalBinaryRatingSystems.assumption_lemmaC4_positive_support_interval_model μ),
                          ¬GJ19OptimalBinaryRatingSystems.lemmaC4FiniteStepOnIoo R.successProb R.lo R.hi →
                            (GJ19OptimalBinaryRatingSystems.lemmaC4FiniteRangeOnIoo
                                  (GJ19OptimalBinaryRatingSystems.cutpointStepSuccessProb cut levels) lo hi ∧
                                ∃ rate,
                                  0 < rate ∧
                                    EconCSLib.Probability.ExponentialRateCertificate
                                      (GJ19OptimalBinaryRatingSystems.lemmaC4TieErasedSourceWbar μ
                                        (GJ19OptimalBinaryRatingSystems.theorem31SelectedPullbackSourceWeight weight)
                                        (GJ19OptimalBinaryRatingSystems.theorem31SelectedPullbackSourceKernel μ cut
                                          hmono sampleRate levels hlevels))
                                      rate) ∧
                              EconCSLib.Probability.HasExponentialRate
                                  (GJ19OptimalBinaryRatingSystems.lemmaC4TieErasedSourceWbar μ R.weight
                                    (GJ19OptimalBinaryRatingSystems.lemmaC4RawSourcePbarKernel R))
                                  0 ∧
                                ∀ (rate : ℝ),
                                  0 < rate →
                                    ¬EconCSLib.Probability.ExponentialRateCertificate
                                        (GJ19OptimalBinaryRatingSystems.lemmaC4TieErasedSourceWbar μ R.weight
                                          (GJ19OptimalBinaryRatingSystems.lemmaC4RawSourcePbarKernel R))
                                        rate
\end{ReviewVerbatim}

\paragraph{Lean proof endpoint (built separately)}\begin{ReviewVerbatim}
GJ19OptimalBinaryRatingSystems.PaperInterface.paper_lemmaC4_finite_selected_pullback_positive_rate_and_nonfiniteStep_source_zero_rate
\end{ReviewVerbatim}

\paragraph{Recorded source-to-Spec screening}
\textbf{Verdict:} \texttt{not recorded}
\\
\textbf{Reason:} not recorded
\reviewmatch{Matches source input}
\noindent\textbf{Reviewer annotation}\par
\reviewerbox
\clearpage
\hypertarget{source-claim-15e577a1127ae1d6}{}
\section*{Review row 43}
\textbf{Source locator:} {\footnotesize\raggedright cited publication, lines 2405-\allowbreak{}2409\par}
\paragraph{Verbatim source input}
\begin{ReviewVerbatim}
convex, with minima at a = b, when a, b ∈ (0, 1). Note that inf a {gi KL(a||ti ) + g(i + 1)KL(a||ti+1 )},
for all feasible g, is also continuous and strictly convex in ti , ti+1 , with minima at ti = ti+1 .
One consequence of the above fact is that fixing either ti or ti+1 and moving the other farther
away monotonically increases KL, while moving it closer decreases KL.
C.4.1
\end{ReviewVerbatim}


\paragraph{Expanded PaperInterface specification}
\begin{ReviewVerbatim}
∀ {gHi gLo pLo pHi : ℝ},
  0 < gHi →
    0 < gLo →
      0 < pLo →
        pLo < pHi →
          pHi < 1 →
            ContinuousAt (fun q => EconCSLib.Probability.weightedBernoulliClosedThresholdRate gHi gLo q.1 q.2)
                (pHi, pLo) ∧
              (∀ p ∈ Set.Ioo 0 1, EconCSLib.Probability.weightedBernoulliClosedThresholdRate gHi gLo p p = 0) ∧
                0 < EconCSLib.Probability.weightedBernoulliClosedThresholdRate gHi gLo pHi pLo ∧
                  StrictMonoOn (fun x => EconCSLib.Probability.weightedBernoulliClosedThresholdRate gHi gLo x pLo)
                      (Set.Icc pLo pHi) ∧
                    StrictAntiOn (fun x => EconCSLib.Probability.weightedBernoulliClosedThresholdRate gHi gLo pHi x)
                      (Set.Icc pLo pHi)
\end{ReviewVerbatim}

\paragraph{Lean proof endpoint (built separately)}\begin{ReviewVerbatim}
GJ19OptimalBinaryRatingSystems.PaperInterface.source_remarkC2_weighted_kl_separation_monotonicity
\end{ReviewVerbatim}

\paragraph{Recorded source-to-Spec screening}
\textbf{Verdict:} \texttt{not recorded}
\\
\textbf{Reason:} not recorded
\reviewmatch{Matches source input}
\noindent\textbf{Reviewer annotation}\par
\reviewerbox
\clearpage
\hypertarget{source-claim-f9871b4133903309}{}
\section*{Review row 44}
\textbf{Source locator:} {\footnotesize\raggedright cited publication, lines 2670-\allowbreak{}2714\par}
\paragraph{Verbatim source input}
\begin{ReviewVerbatim}
Additional necessary lemmas

Now, we begin the set-up that will lead to a proof for Theorem 3.2. It turns out that proving the
theorem requires, in the process, essentially proving our convergence result with M → ∞, Theorem B.1. For Theorem 3.2, we need a lower bound for t1 as a function of M . This seems hard
to do in general. Luckily, in our case, there is a property for how t∗ changes when M is doubled.
Using this property, we can derive that t∗1 ≥ O(M −3 ).
−1
Recall that step-wise increasing β with M intervals Si = [si , si+1 ) has levels {ti }M
i=0 , where
t0 = 1, tM −1 = 1, and s0 , 0, sM , 1.
Furthermore, we use the following notation for the large deviation rate
#
"
g
g

ri = −(gi−1 + gi ) log (1 − ti−1 )

gi−1
gi−1 +gi

(1 − ti )

gi
gi−1 +gi

i−1
gi−1 +gi

+ ti−1

g

i
+gi

ti i−1

(23)

for i ∈ {1 . . . M − 1}, which implies r1 = −g1 log(1 − t1 ) and rM −1 = −gM −2 log(tM −2 ).
We further use rM −1 to be the rate achieved by the optimal βM with M intervals.
−1
Lemma C.5. Suppose g uniform, i.e. gi = 1, ∀i and that βM has values√{ti }M
\end{ReviewVerbatim}


\paragraph{Expanded PaperInterface specification}
\begin{ReviewVerbatim}
∀ {m : ℕ} (successProb sampleRate : Fin (m + 2) → ℝ) (i : Fin (m + 1)),
  GJ19OptimalBinaryRatingSystems.binaryEndpointAwareAdjacentRate successProb sampleRate i =
    if ↑i = 0 then
      sampleRate (GJ19OptimalBinaryRatingSystems.adjacentHighIndex i) *
        -Real.log (1 - successProb (GJ19OptimalBinaryRatingSystems.adjacentHighIndex i))
    else
      if ↑i = m then
        sampleRate (GJ19OptimalBinaryRatingSystems.adjacentLowIndex i) *
          -Real.log (successProb (GJ19OptimalBinaryRatingSystems.adjacentLowIndex i))
      else
        EconCSLib.Probability.weightedBernoulliClosedThresholdRate
          (sampleRate (GJ19OptimalBinaryRatingSystems.adjacentHighIndex i))
          (sampleRate (GJ19OptimalBinaryRatingSystems.adjacentLowIndex i))
          (successProb (GJ19OptimalBinaryRatingSystems.adjacentHighIndex i))
          (successProb (GJ19OptimalBinaryRatingSystems.adjacentLowIndex i))
\end{ReviewVerbatim}

\paragraph{Lean proof endpoint (built separately)}\begin{ReviewVerbatim}
GJ19OptimalBinaryRatingSystems.PaperInterface.source_appendixC5_rate_notation_eq23
\end{ReviewVerbatim}

\paragraph{Recorded source-to-Spec screening}
\textbf{Verdict:} \texttt{not recorded}
\\
\textbf{Reason:} not recorded
\reviewmatch{Matches source input}
\noindent\textbf{Reviewer annotation}\par
\reviewerbox
\clearpage
\hypertarget{source-claim-f8aedb48228e5085}{}
\section*{Review row 45}
\textbf{Source locator:} {\footnotesize\raggedright cited publication, lines 2715-\allowbreak{}2895\par}
\paragraph{Verbatim source input}
\begin{ReviewVerbatim}
i=0[U+0001 control character] . Then β2M −1
2M −2
1
0
0
0
has values {ti }i=0 , where t2i = ti , ∀i ∈ {0 . . . M − 1}, t1 = 2 1 − 1 − t1 and t02M −3 =
√
1
2 (1 + tM −2 ).

Proof. We first set the values t02i = ti and then optimally choose the remaining values t0k , k odd.
Then, we show that the resulting large deviation rates between all adjacent pairs are equal. Then,
by the proof of Lemma 3.1, which showed that equalizing the rates between adjacent intervals is a
−2
sufficient condition for optimality, β2M −1 has the levels {t0i }2M
i=0 .
Let r0 denote rates between adjacent t0 as r does for t. Supposing t02 = t1 , we find t01 such that
0
r1 = r20 and t01 < t02 .
[U+0014 control character]q
[U+0015 control character]
q
0
0
0
0
0
− log(1 − t1 ) = −2 log
(1 − t1 )(1 − t2 ) + t1 t2
q
=⇒ 1 − t01 = (1 − t01 )(1 − t02 ) + t01 t02 + 2 (1 − t01 )(1 − t02 )t01 t02
[U+0012 control character]
[U+0013 control character]
q
[U+0001 control character]
√
1
1
0
0
=⇒ t1 =
1 − 1 − t2 =
1 − 1 − t1
2
2
1
0
0
0
Similarly, r2M
−3 = r2M −2 when t2M −3 = 2 (1 +
by choosing such t01 , t02M −3 .

√

0
0
tM −2 ). It follows that r10 = r20 = r2M
−3 = r2M −2

17


[form-feed in source extraction]
0
Next, we find t0k ∈ (t0k−1 , t0k+1 ) for k ∈ {3, 5, . . . 2M − 5} such that the rates rk0 = rk+1
.

q
q
hq
i
hq
i
0
0
0
0
0
0
(1 − tk )(1 − tk−1 ) + tk tk−1 = −2 log
(1 − tk )(1 − tk+1 ) + t0k t0k+1
−2 log
q
[U+F8F9 private-use glyph]2
[U+F8EE private-use glyph]q
0
0
1
−
t
−
1
−
t
k+1
k−1
c
[U+F8FB private-use glyph]
q
q
=⇒ t0k =
, where c = [U+F8F0 private-use glyph]
1+c
0
0
− t
t
k−1

k+1

Now, we show that rk0 = rj0 , ∀j, k by showing that the difference between each rate ri and its
0 is constant. r = r , ∀j, k by assumption and so r 0 = r 0 , ∀j, k follows.
analogous rate r2i
j
k
j
k
√
1
0
(1
+
t
).
Thus
if
r
=
− log x for some x, then
rM −1 = − log tM −2 and r2M
=
−
log
i
M
−2
−2
2
√
r2i = − log 21 (1 + x) would imply that all the rates are equal. Thus, it is sufficient to show that
q
i2 1 h
i
hq
p
p
(1 − t02i−1 )(1 − t02i ) + t02i−1 t02i =
1 + (1 − ti−1 )(1 − ti ) + ti−1 ti
2
"s[U+0012 control character]
#2
r
[U+0013 control character]
i
p
p
c
c
1h
≡
1−
1 + (1 − ti−1 )(1 − ti ) + ti−1 ti
(1 − ti ) +
ti =
1+c
1+c
2
[U+0015 control character]2
[U+0014 control character]√
√
1 − ti − 1 − ti−1
√
where c =
√
ti−1 − ti

(24)
(25)

The proof for (25) is algebraically tedious and is shown in Remark C.3 below.
\end{ReviewVerbatim}


\paragraph{Expanded PaperInterface specification}
\begin{ReviewVerbatim}
∀ {m : ℕ},
  0 < m →
    ∀ {oldLevels : Fin (m + 2) → ℝ},
      GJ19OptimalBinaryRatingSystems.BinaryEndpointLevelVector oldLevels →
        (GJ19OptimalBinaryRatingSystems.BinaryEndpointAwareAdjacentRatesEqualize oldLevels fun x => 1) →
          GJ19OptimalBinaryRatingSystems.BinaryEndpointLevelVector
              (GJ19OptimalBinaryRatingSystems.uniformDoubledEndpointLevels oldLevels) ∧
            (GJ19OptimalBinaryRatingSystems.BinaryEndpointAwareAdjacentRatesEqualize
                (GJ19OptimalBinaryRatingSystems.uniformDoubledEndpointLevels oldLevels) fun x => 1) ∧
              EconCSLib.Optimization.IsMaximizerOn GJ19OptimalBinaryRatingSystems.BinaryEndpointLevelVector
                (fun xs => GJ19OptimalBinaryRatingSystems.binaryEndpointAwareAdjacentRateObjective xs fun x => 1)
                (GJ19OptimalBinaryRatingSystems.uniformDoubledEndpointLevels oldLevels)
\end{ReviewVerbatim}

\paragraph{Lean proof endpoint (built separately)}\begin{ReviewVerbatim}
GJ19OptimalBinaryRatingSystems.PaperInterface.source_lemmaC5_uniform_doubled_chain
\end{ReviewVerbatim}

\paragraph{Recorded source-to-Spec screening}
\textbf{Verdict:} \texttt{not recorded}
\\
\textbf{Reason:} not recorded
\reviewmatch{Matches source input}
\noindent\textbf{Reviewer annotation}\par
\reviewerbox
\clearpage
\hypertarget{source-claim-428f19cd38ba7b0a}{}
\section*{Review row 46}
\textbf{Source locator:} {\footnotesize\raggedright cited publication, lines 2884-\allowbreak{}3038\par}
\paragraph{Verbatim source input}
\begin{ReviewVerbatim}
[U+0014 control character]√
√
1 − ti − 1 − ti−1
√
where c =
√
ti−1 − ti

(24)
(25)

The proof for (25) is algebraically tedious and is shown in Remark C.3 below.
Then, by the proof of Lemma 3.1, which shows that equalizing the rates inside the minimization
−2
terms implies an optimal {ti }, β2M −1 has the levels {t0i }2M
i=0 .
Remark C.3.
"s[U+0012 control character]

c
1−
1+c

r

[U+0013 control character]
(1 − ti ) +

#2

i
p
p
1h
1 + (1 − ti−1 )(1 − ti ) + ti−1 ti
2
[U+0015 control character]2
[U+0014 control character]√
√
1 − ti − 1 − ti−1
√
where c =
√
ti−1 − ti

c
ti
1+c

=

√
√
√
√
Proof. Let x = ti , y = 1 − ti , z = ti−1 , and w = 1 − ti−1 . Note that x > z, w > y, y =
1 − x2 , w = 1 − z 2 . Then,
c
(y − w)2
1
(x − z)2
=
, and
=
c+1
2 − 2xz − 2yw
c+1
2 − 2xz − 2yw
1
(To show the above two equalities, factor out (x−z)
2 from numerator and denominator, and substi2
2
tute y = 1 − x , w = 1 − z ).

18


[form-feed in source extraction]
Now, the left hand side:
"s[U+0012 control character]

c
1−
1+c

#2
c
(1 − ti ) +
ti
1+c
i2
hp
p
(x − z)2 y 2 + (y − w)2 x2
r

[U+0013 control character]

1
2 − 2xz − 2yw
(x − z)2 y 2 + (y − w)2 x2 + 2xy(x − z)(w − y)
=
2 − 2xz − 2yw
2
2
2
2
z y + w x − 2wxyz
=
2 − 2xz − 2yw
=

p
p
(y − w)2 = w − y, (x − z)2 = x − z

The right hand side:
i
p
p
1h
1 + (1 − ti−1 )(1 − ti ) + ti−1 ti
2
1
= [1 + (wy + xz)]
2
Multiplying both sides by 2 − 2xz − 2yw, we have:
"s[U+0012 control character]

c
1−
1+c

r

[U+0013 control character]
(1 − ti ) +

c
ti
1+c

#2
=

i
p
p
1h
1 + (1 − ti−1 )(1 − ti ) + ti−1 ti
2

≡ z 2 y 2 + w2 x2 − 2wxyz = 1 − (wy + xz)2
≡ z 2 (1 − x2 ) + (1 − z 2 )x2 − 2wxyz = 1 − w2 y 2 − x2 z 2 − 2wxyz
≡ z 2 − 2x2 z 2 + x2 = 1 − (1 − z 2 )(1 − x2 ) − x2 z 2
≡0=0
\end{ReviewVerbatim}


\paragraph{Expanded PaperInterface specification}
\begin{ReviewVerbatim}
(∀ {m k : ℕ},
    0 < k →
      ∀ (hkm : k < m) (oldLevels : Fin (m + 2) → ℝ),
        GJ19OptimalBinaryRatingSystems.uniformDoubledEndpointLevels oldLevels ⟨2 * k + 1, ⋯⟩ =
          EconCSLib.Probability.bernoulliInteriorEqualSplit (oldLevels ⟨k, ⋯⟩) (oldLevels ⟨k + 1, ⋯⟩)) ∧
  ∀ {pLo pHi : ℝ},
    0 ≤ pLo →
      pHi ≤ 1 →
        pLo < pHi →
          EconCSLib.Probability.weightedBernoulliClosedThresholdRate 1 1
              (EconCSLib.Probability.bernoulliInteriorEqualSplit pLo pHi) pLo =
            EconCSLib.Probability.weightedBernoulliClosedThresholdRate 1 1 pHi
              (EconCSLib.Probability.bernoulliInteriorEqualSplit pLo pHi)
\end{ReviewVerbatim}

\paragraph{Lean proof endpoint (built separately)}\begin{ReviewVerbatim}
GJ19OptimalBinaryRatingSystems.PaperInterface.source_lemmaC5_refinement_equations24_25
\end{ReviewVerbatim}

\paragraph{Recorded source-to-Spec screening}
\textbf{Verdict:} \texttt{not recorded}
\\
\textbf{Reason:} not recorded
\reviewmatch{Matches source input}
\noindent\textbf{Reviewer annotation}\par
\reviewerbox
\clearpage
\hypertarget{source-claim-4cea7b91f971b40f}{}
\section*{Review row 47}
\textbf{Source locator:} {\footnotesize\raggedright cited publication, lines 3042-\allowbreak{}3092\par}
\paragraph{Verbatim source input}
\begin{ReviewVerbatim}
0

Proof. Let M = 2N , M 0 = 2N +1 − 1, for some N . We show that rM ≤ 12 rM . The corollary follows
0
by noting that rK < rK ∀K 0 > K and that rK < ∞, ∀K.
0

0

M
rM − rM = − log tM
M −2 + log tM 0 −2
[U+0014 control character]
[U+0015 control character]
q
1 1 M
M
= − log tM −2 + log
+
t
2 2 M −2
[U+F8EE private-use glyph]
[U+F8F9 private-use glyph]
1 1
1
1 [U+F8FB private-use glyph]
= log [U+F8F0 private-use glyph] M + q
2 tM −2 2 tM

Lemma C.5

M −2

1
≥ − log tM
M −2
2
1
0
=⇒ rM ≤ rM
2

q

19

M
tM
M −2 ≥ tM −2


[form-feed in source extraction]
M
\end{ReviewVerbatim}


\paragraph{Expanded PaperInterface specification}
\begin{ReviewVerbatim}
∀ (levels : (N : ℕ) → Fin (N + 1 + 2) → ℝ),
  (∀ (N : ℕ), GJ19OptimalBinaryRatingSystems.BinaryEndpointLevelVector (levels N)) →
    (∀ (N : ℕ), GJ19OptimalBinaryRatingSystems.BinaryEndpointAwareAdjacentRatesEqualize (levels N) fun x => 1) →
      Filter.Tendsto
        (fun N =>
          GJ19OptimalBinaryRatingSystems.binaryEndpointAwareAdjacentRate (levels N) (fun x => 1)
            GJ19OptimalBinaryRatingSystems.lastAdjacentIndex)
        Filter.atTop (nhds 0)
\end{ReviewVerbatim}

\paragraph{Lean proof endpoint (built separately)}\begin{ReviewVerbatim}
GJ19OptimalBinaryRatingSystems.PaperInterface.paper_corollaryC2_uniform_equalized_last_rate_tendsto_zero
\end{ReviewVerbatim}

\paragraph{Recorded source-to-Spec screening}
\textbf{Verdict:} \texttt{not recorded}
\\
\textbf{Reason:} not recorded
\reviewmatch{Matches source input}
\noindent\textbf{Reviewer annotation}\par
\reviewerbox
\clearpage
\hypertarget{source-claim-510e870d600e0e8e}{}
\section*{Review row 48}
\textbf{Source locator:} {\footnotesize\raggedright cited publication, lines 3094-\allowbreak{}3125\par}
\paragraph{Verbatim source input}
\begin{ReviewVerbatim}
k − tk−1 < δ.

Proof. This corollary follows directly from Corollary C.1. If the rates are upper bounded, then so
are the level differences.
We first find where the rate is minimized given a width between levels of δ
hp
i
p
xm = arg min −2 log
(1 − x − δ)(1 − x) + x(x + δ)
x

1 1
= − δ
2 2
Then given an upper bound of [U+000F control character] on the rate, there is a bound on δ determined by the largest
possible difference at levels symmetric around 21 .
" r

1
1
r = −2 log 2 ( − δ)( + δ)
2
2
[U+0002 control character]
[U+0003 control character]
= − log 1 − 4δ 2
1p
≥ [U+000F control character] when δ >
1 − e−[U+000F control character]
2
\end{ReviewVerbatim}


\paragraph{Expanded PaperInterface specification}
\begin{ReviewVerbatim}
∀ (levels : (N : ℕ) → Fin (N + 1 + 2) → ℝ),
  (∀ (N : ℕ), GJ19OptimalBinaryRatingSystems.BinaryEndpointLevelVector (levels N)) →
    (∀ (N : ℕ), GJ19OptimalBinaryRatingSystems.BinaryEndpointAwareAdjacentRatesEqualize (levels N) fun x => 1) →
      Filter.Tendsto (fun N => GJ19OptimalBinaryRatingSystems.binaryEndpointAdjacentMaxWidth (levels N)) Filter.atTop
        (nhds 0)
\end{ReviewVerbatim}

\paragraph{Lean proof endpoint (built separately)}\begin{ReviewVerbatim}
GJ19OptimalBinaryRatingSystems.PaperInterface.paper_corollaryC2_uniform_equalized_adjacent_mesh_tendsto_zero
\end{ReviewVerbatim}

\paragraph{Recorded source-to-Spec screening}
\textbf{Verdict:} \texttt{not recorded}
\\
\textbf{Reason:} not recorded
\reviewmatch{Matches source input}
\noindent\textbf{Reviewer annotation}\par
\reviewerbox
\clearpage
\hypertarget{source-claim-6e85310ac4deebe8}{}
\section*{Review row 49}
\textbf{Source locator:} {\footnotesize\raggedright cited publication, lines 3131-\allowbreak{}3159\par}
\paragraph{Verbatim source input}
\begin{ReviewVerbatim}
Proof. Note that, with uniform matching, ∀x ∈ (0, 1], y ∈ [0, 1 − x] the rate with values ti−1 =
y, ti = y + x is no more than the last with tM −2 = 1 − x. With width x, in other words, the extreme
points have a larger rates than the middle points. For i ∈
/ {1, M − 1}:
ri = inf {gi−1 KL(a||ti−1 ) + gi KL(a||ti )}
a

= inf {KL(a||y) + KL(a||y + x)}
uniform matching
a
h
i
1
1
1
1
= −2 log (1 − y) 2 (1 − y − x) 2 + y 2 (y + x) 2
(26)
h
i
= − log (1 − y)(1 − y − x) + y(y + x) + 2 [(1 − y)(1 − y − x)y(y + x)]1/2
≤ − log(1 − x)
Where line (26) follows from line (22).
By the proof of Lemma 3.1, the optimal levels equalize the rates between each level. Then,
when g is non-decreasing, gM −2 ≥ g` , ∀` ∈ {1 . . . M − 3}. Then, at the same level differences, the
rate corresponding to the last level is no smaller. Thus, to equalize the rates, the last width must
be no larger than any other width. Thus, tM −2 ≥ 1 − L1 .
N +1 −1
\end{ReviewVerbatim}


\paragraph{Expanded PaperInterface specification}
\begin{ReviewVerbatim}
∀ {m : ℕ},
  0 < m →
    ∀ {levels sampleRate : Fin (m + 2) → ℝ},
      GJ19OptimalBinaryRatingSystems.BinaryEndpointLevelVector levels →
        GJ19OptimalBinaryRatingSystems.BinaryEndpointAwareAdjacentRatesEqualize levels sampleRate →
          (∀ (idx : Fin (m + 2)), 0 < sampleRate idx) →
            (∀ {a b : Fin (m + 2)}, ↑a ≤ ↑b → sampleRate a ≤ sampleRate b) →
              1 - 1 / ↑(m + 1) ≤
                levels
                  (GJ19OptimalBinaryRatingSystems.adjacentLowIndex GJ19OptimalBinaryRatingSystems.lastAdjacentIndex)
\end{ReviewVerbatim}

\paragraph{Lean proof endpoint (built separately)}\begin{ReviewVerbatim}
GJ19OptimalBinaryRatingSystems.PaperInterface.lemmaC6_monotone_matching_penultimate_level_bound
\end{ReviewVerbatim}

\paragraph{Recorded source-to-Spec screening}
\textbf{Verdict:} \texttt{not recorded}
\\
\textbf{Reason:} not recorded
\reviewmatch{Matches source input}
\noindent\textbf{Reviewer annotation}\par
\reviewerbox
\clearpage
\hypertarget{source-claim-2d8e50cd70ac2429}{}
\section*{Review row 50}
\textbf{Source locator:} {\footnotesize\raggedright cited publication, lines 3161-\allowbreak{}3233\par}
\paragraph{Verbatim source input}
\begin{ReviewVerbatim}
20

N

≥ 15 r2 .


[form-feed in source extraction]
1
Proof. Let K = 2N , K 0 = 2N +1 − 1. Note that tK
K−1 ≥ 2 by Lemma C.6.
0

0

K
rK − rK = − log tK
K−2 + log tK 0 −2
[U+0014 control character]
[U+0015 control character]
q
1 1 K
K
= − log tK−2 + log
+
t
2 2 K−2
[U+F8EE private-use glyph]
[U+F8F9 private-use glyph]
1
1
1
1
[U+F8FB private-use glyph]
= log [U+F8F0 private-use glyph] K + q
2 tK−2 2 tK

Lemma C.5

K−2

≤ log

tK
K−2

[U+0001 control character]

[U+0014 control character]
[U+0015 control character]
[U+0001 control character]− 45
1 K [U+0001 control character]−1 1 K [U+0001 control character]− 21
1
K
K
≤ tK−2
when tK−2 ∈
t
+
t
,1
2 K−2
2 K−2
2

− 54

1
0
=⇒ rK ≥ rK
5

−3 .
Lemma C.8. With uniform matching (gi = 1), ∃C > 0 s.t. ∀M, tM
\end{ReviewVerbatim}


\paragraph{Expanded PaperInterface specification}
\begin{ReviewVerbatim}
∀ {m : ℕ},
  0 < m →
    ∀ {oldLevels : Fin (m + 2) → ℝ},
      GJ19OptimalBinaryRatingSystems.BinaryEndpointLevelVector oldLevels →
        (GJ19OptimalBinaryRatingSystems.BinaryEndpointAwareAdjacentRatesEqualize oldLevels fun x => 1) →
          (1 / 5 * GJ19OptimalBinaryRatingSystems.binaryEndpointAwareAdjacentRateObjective oldLevels fun x => 1) ≤
            GJ19OptimalBinaryRatingSystems.binaryEndpointAwareAdjacentRateObjective
              (GJ19OptimalBinaryRatingSystems.lemmaC5_uniform_doubled_endpoint_levels oldLevels) fun x => 1
\end{ReviewVerbatim}

\paragraph{Lean proof endpoint (built separately)}\begin{ReviewVerbatim}
GJ19OptimalBinaryRatingSystems.PaperInterface.lemmaC7_uniform_doubled_objective_rate_ge_one_fifth_old_objective
\end{ReviewVerbatim}

\paragraph{Recorded source-to-Spec screening}
\textbf{Verdict:} \texttt{not recorded}
\\
\textbf{Reason:} not recorded
\reviewmatch{Matches source input}
\noindent\textbf{Reviewer annotation}\par
\reviewerbox
\clearpage
\hypertarget{source-claim-905297a638ebd8e0}{}
\section*{Review row 51}
\textbf{Source locator:} {\footnotesize\raggedright cited publication, lines 3234-\allowbreak{}3267\par}
\paragraph{Verbatim source input}
\begin{ReviewVerbatim}
1 ≥ CM

Proof. By Lemma C.7, ∃C2 > 0 s.t. rM ≥ C2 5−dlog2 M e . Then
M
− log(1 − tM
1 )=r

≥ C2 5−dlog2 M e
h
i
−dlog2 M e
=⇒ tM
1 ≥ 1 − exp −C2 5
h
i
− 1
≥ 1 − exp −C3 M log5 2
e−1
− 1
C3 M log5 2
e
M
=⇒ ∃C > 0 s.t. t1 ≥ CM −3

e−x ≤ 1 −

≥

e−1
x for x ∈ [0, 1]
e

−3 .
Corollary C.3. With monotonically non-decreasing g, ∃C > 0 s.t. ∀M, tM
\end{ReviewVerbatim}


\paragraph{Expanded PaperInterface specification}
\begin{ReviewVerbatim}
∀ {m : ℕ},
  0 < m →
    ∀ {levels : Fin (m + 2) → ℝ},
      GJ19OptimalBinaryRatingSystems.BinaryEndpointLevelVector levels →
        (GJ19OptimalBinaryRatingSystems.BinaryEndpointAwareAdjacentRatesEqualize levels fun x => 1) →
          (1 / ↑(m + 1)) ^ 2 / 2 ≤
            levels (GJ19OptimalBinaryRatingSystems.adjacentHighIndex GJ19OptimalBinaryRatingSystems.firstAdjacentIndex)
\end{ReviewVerbatim}

\paragraph{Lean proof endpoint (built separately)}\begin{ReviewVerbatim}
GJ19OptimalBinaryRatingSystems.PaperInterface.source_lemmaC8_uniform_first_level_polynomial_lower_bound
\end{ReviewVerbatim}

\paragraph{Recorded source-to-Spec screening}
\textbf{Verdict:} \texttt{not recorded}
\\
\textbf{Reason:} not recorded
\reviewmatch{Matches source input}
\noindent\textbf{Reviewer annotation}\par
\reviewerbox
\clearpage
\hypertarget{source-claim-d7f37d83d608e4fa}{}
\section*{Review row 52}
\textbf{Source locator:} {\footnotesize\raggedright cited publication, lines 3268-\allowbreak{}3274\par}
\paragraph{Verbatim source input}
\begin{ReviewVerbatim}
1 ≥ CM

Proof. The result follows from noting that tM
1 with uniform matching lower bounds the first value
with any other monotonically non-decreasing g, which is a direct application of Lemma B.1 – scale
g such that g1 = 1. Then, gj ≥ 1, j > 1 and g0 ≤ 1. Then, the condition of the lemma holds.
Lemma C.9. The run-time of NestedBisection is O(M log2 1δ ), where δ is the bisection grid width
\end{ReviewVerbatim}


\paragraph{Expanded PaperInterface specification}
\begin{ReviewVerbatim}
∀ {m : ℕ},
  0 < m →
    ∀ {levels sampleRate : Fin (m + 2) → ℝ},
      GJ19OptimalBinaryRatingSystems.BinaryEndpointLevelVector levels →
        GJ19OptimalBinaryRatingSystems.BinaryEndpointAwareAdjacentRatesEqualize levels sampleRate →
          (∀ (idx : Fin (m + 2)), 0 < sampleRate idx) →
            (∀ {a b : Fin (m + 2)}, ↑a ≤ ↑b → sampleRate a ≤ sampleRate b) →
              sampleRate
                    (GJ19OptimalBinaryRatingSystems.adjacentHighIndex
                      GJ19OptimalBinaryRatingSystems.firstAdjacentIndex) =
                  1 →
                (1 / ↑(m + 1)) ^ 2 / 2 ≤
                  levels
                    (GJ19OptimalBinaryRatingSystems.adjacentHighIndex GJ19OptimalBinaryRatingSystems.firstAdjacentIndex)
\end{ReviewVerbatim}

\paragraph{Lean proof endpoint (built separately)}\begin{ReviewVerbatim}
GJ19OptimalBinaryRatingSystems.PaperInterface.corollaryC3_monotone_scaled_first_level_ge_half_inv_adjacent_count_sq
\end{ReviewVerbatim}

\paragraph{Recorded source-to-Spec screening}
\textbf{Verdict:} \texttt{not recorded}
\\
\textbf{Reason:} not recorded
\reviewmatch{Matches source input}
\noindent\textbf{Reviewer annotation}\par
\reviewerbox
\clearpage
\hypertarget{source-claim-8cc85076f8ab0b2b}{}
\section*{Review row 53}
\textbf{Source locator:} {\footnotesize\raggedright cited publication, lines 3275-\allowbreak{}3279\par}
\paragraph{Verbatim source input}
\begin{ReviewVerbatim}
and M is the number of intervals.
Proof. The outer bisection, in main, runs at most log2 2δ + 1 iterations. Each outer iteration calls
BisectNextLevel M − 3 times, and the inner bisection in each call runs for at most log2 2δ iterations.
Thus the run-time of algorithm is O(M log2 1δ ).
\end{ReviewVerbatim}


\paragraph{Expanded PaperInterface specification}
\begin{ReviewVerbatim}
∀ {M : ℕ},
  0 < M →
    ∀ {delta : ℝ},
      0 < delta →
        ∃ L,
          1 ≤ delta * 2 ^ L ∧
            ↑(L + 1) ≤ Real.logb 2 (max 1 (1 / delta)) + 2 ∧
              ↑(GJ19OptimalBinaryRatingSystems.nestedBisectionOperationCount M (L + 1) L) ≤
                ↑M * (Real.logb 2 (max 1 (1 / delta)) + 2) ^ 2
\end{ReviewVerbatim}

\paragraph{Lean proof endpoint (built separately)}\begin{ReviewVerbatim}
GJ19OptimalBinaryRatingSystems.PaperInterface.source_lemmaC9_nested_bisection_runtime_log_squared
\end{ReviewVerbatim}

\paragraph{Recorded source-to-Spec screening}
\textbf{Verdict:} \texttt{not recorded}
\\
\textbf{Reason:} not recorded
\reviewmatch{Matches source input}
\noindent\textbf{Reviewer annotation}\par
\reviewerbox
\clearpage
\hypertarget{source-claim-51c42bc84b84361b}{}
\section*{Review row 54}
\textbf{Source locator:} {\footnotesize\raggedright cited publication, lines 3467-\allowbreak{}3589\par}
\paragraph{Verbatim source input}
\begin{ReviewVerbatim}
Proof of Theorem B.1

w denote the optimal β with M intervals for weight function w, with intervals swM and levels
Let βM
twM . Let qwM (θ) = i/M when θ ∈ [swM
, swM
i
i+1 ), i.e. the quantile of interval item of type θ is in.
Then we have the following convergence result for βM .

Theorem B.1. Let g be uniform. Suppose w such that qwM converges uniformly. Then, ∀C ∈
w uniformly as N → ∞.
w
N, ∃β w s.t. βC2
N +1 → β
hProof. Note that
[U+0011 control character] the condition on q implies that ∃M s.t.
M
M
sbxθ M c , sdxθ M e .

∀M > M , ∀θ, ∃xθ such that θ ∈

h
[U+0011 control character]
M
Let M 0 = 2M − 1, M 00 = 4M − 3, M q = 2q M − 2q + 1. θ ∈ sM
,
s
=⇒ βM (θ) =
bxθ M c dxθ M e
h
i
M
M
tM
bxθ M c ∈ tbxθ M c−1 , tbxθ M c+1 . Then,
0

βM 0 (θ) = tM
bxθ M 0 c
0

= tM
bx (2M −1)c
h θ0
i
M0
∈ tM
,
t
2bxθ M c−2 2bxθ M c+2
i
h
M
⊂ tM
,
t
bxθ M c−1 bxθ M c+1

23

Lemma C.5


[form-feed in source extraction]
And, for general q,
q

βM q (θ) = tM
bx (2q M −2q +1)c
i
h θq
Mq
,
t
∈ tM
q
q
q
bxθ (2 M )c+1
bxθ 2 M c−2
h q
i
Mq
⊂ tM
,
t
q
q
q
2 bxθ M c−2
2 bxθ M c+1
h
i
M
⊂ tM
,
t
bxθ M c−1 bxθ M c+1

Lemma C.5

h
i
M
Then, ∀N 0 > 1, θ: β2N 0 M −2N 0 +1 (θ) ∈ tM
,
t
bxθ M c−1 bxθ M c+1 and
M
|β2N 0 M −2N 0 +1 (θ) − βM (θ)| ≤ tM
bxθ M c+1 − tbxθ M c−1
0

0

K
By Corollary C.2, ∀δ > 0, ∃K s.t. ∀K 0 > K, tK
bxθ K 0 c+1 − tbxθ K 0 c−1 < 2δ.
By the Cauchy criterion, ∃β s.t. β(C−1)2N +1 → β uniformly.
By change of variables, ∃β s.t. βC2N +1 → β uniformly.

Corollary C.4. For Kendall’s tau and Spearman’s rho correlation measures, ∃β s.t. β2N → β
\end{ReviewVerbatim}


\paragraph{Expanded PaperInterface specification}
\begin{ReviewVerbatim}
∀ {M q : ℕ},
  0 < M →
    ∀ {θ : ℝ},
      0 ≤ θ →
        θ ≤ 1 →
          have old := ⌊↑M * θ⌋₊;
          have refined := ⌊↑(GJ19OptimalBinaryRatingSystems.theoremB1SourceDoubledIndexIterate M q) * θ⌋₊;
          refined ≤ 2 ^ q * old + 2 ^ q ∧ 2 ^ q * old ≤ refined + 2 ^ q
\end{ReviewVerbatim}

\paragraph{Lean proof endpoint (built separately)}\begin{ReviewVerbatim}
GJ19OptimalBinaryRatingSystems.PaperInterface.source_theoremB1_proof_selector_nesting
\end{ReviewVerbatim}

\paragraph{Recorded source-to-Spec screening}
\textbf{Verdict:} \texttt{not recorded}
\\
\textbf{Reason:} not recorded
\reviewmatch{Matches source input}
\noindent\textbf{Reviewer annotation}\par
\reviewerbox
\clearpage
\hypertarget{source-claim-b7febbacfb3801fa}{}
\section*{Review row 55}
\textbf{Source locator:} {\footnotesize\raggedright cited publication, lines 3590-\allowbreak{}3593\par}
\paragraph{Verbatim source input}
\begin{ReviewVerbatim}
uniformly as N → ∞.
Proof. For Kendall’s tau and Spearman’s rho, {si } is spaced such that ∀i, j, si − si−1 = sj − sj−1 .
Thus, xθ = θ meets the criterion.
\end{ReviewVerbatim}


\paragraph{Expanded PaperInterface specification}
\begin{ReviewVerbatim}
∀ (C : ℕ),
  (∀ (M : ℕ),
      Nonempty (Fin M) →
        0 < M →
          EconCSLib.Optimization.IsLexicographicMaximizerOn
            (GJ19OptimalBinaryRatingSystems.ProofBridge.theorem31KendallFiniteDesignFeasible M)
            (GJ19OptimalBinaryRatingSystems.ProofBridge.theorem31KendallFiniteDesignValue M)
            (GJ19OptimalBinaryRatingSystems.ProofBridge.theorem31FiniteDesignEndpointRate M)
            (GJ19OptimalBinaryRatingSystems.ProofBridge.theorem31CanonicalUniformEndpointDesign M)) ∧
    (∀ (M : ℕ),
        Nonempty (Fin M) →
          0 < M →
            EconCSLib.Optimization.IsLexicographicMaximizerOn
              (GJ19OptimalBinaryRatingSystems.ProofBridge.theorem31SpearmanFiniteDesignFeasible M)
              (GJ19OptimalBinaryRatingSystems.ProofBridge.theorem31SpearmanFiniteDesignValue M)
              (GJ19OptimalBinaryRatingSystems.ProofBridge.theorem31FiniteDesignEndpointRate M)
              (GJ19OptimalBinaryRatingSystems.ProofBridge.theorem31CanonicalUniformEndpointDesign M)) ∧
      ∃ betaLimit,
        TendstoUniformlyOn
          (fun N θ =>
            GJ19OptimalBinaryRatingSystems.canonicalUniformEqualizedClampedFloorBetaSeq
              (GJ19OptimalBinaryRatingSystems.theoremB1SubsequenceIndex C N) θ)
          betaLimit Filter.atTop (Set.Icc 0 1)
\end{ReviewVerbatim}

\paragraph{Lean proof endpoint (built separately)}\begin{ReviewVerbatim}
GJ19OptimalBinaryRatingSystems.PaperInterface.source_corollaryC4_kendall_spearman_subsequence
\end{ReviewVerbatim}

\paragraph{Recorded source-to-Spec screening}
\textbf{Verdict:} \texttt{not recorded}
\\
\textbf{Reason:} not recorded
\reviewmatch{Matches source input}
\noindent\textbf{Reviewer annotation}\par
\reviewerbox
\clearpage
\hypertarget{source-claim-196d59db7138fcbd}{}
\section*{Review row 56}
\textbf{Source locator:} {\footnotesize\raggedright cited publication, lines 3595-\allowbreak{}3614\par}
\paragraph{Verbatim source input}
\begin{ReviewVerbatim}
Kendall’s tau and Spearman’s rho related proofs

Definition C.1 (see e.g. Embrechts et al. (2003); Nelsen (2007)). The population version of
Kendall-tau correlation between item true quality and rating scores is proportional to
Z
τ
Pk (θ1 , θ2 )dθ1 dθ2
Wk , 2
θ1 >θ2

Similarly, given items with qualities θ1 , θ2 , θ3 , the population version of Spearman’s rho correlation
between item true quality and rating scores is
Z
ρ
Pk (θ1 , θ3 )dθ1 dθ2 dθ3
Wk , 6
θ1 >θ2 ,θ3

Lemma C.10. Spearman’s ρ can also be written as being proportional to
\end{ReviewVerbatim}


\paragraph{Expanded PaperInterface specification}
\begin{ReviewVerbatim}
∀ (M : ℕ) (s : ℕ → ℝ),
  GJ19OptimalBinaryRatingSystems.kendallConstantWeightIntervalObjective M s = (1 - ∑ i, (s (↑i + 1) - s ↑i) ^ 2) / 2 ∧
    GJ19OptimalBinaryRatingSystems.spearmanLinearWeightIntervalObjective M s = (1 - ∑ i, (s (↑i + 1) - s ↑i) ^ 3) / 6
\end{ReviewVerbatim}

\paragraph{Lean proof endpoint (built separately)}\begin{ReviewVerbatim}
GJ19OptimalBinaryRatingSystems.PaperInterface.definitionC1_kendall_spearman_population_objectives
\end{ReviewVerbatim}

\paragraph{Recorded source-to-Spec screening}
\textbf{Verdict:} \texttt{not recorded}
\\
\textbf{Reason:} not recorded
\reviewmatch{Matches source input}
\noindent\textbf{Reviewer annotation}\par
\reviewerbox
\clearpage
\hypertarget{source-claim-3b01d7ad0b5d48f9}{}
\section*{Review row 57}
\textbf{Source locator:} {\footnotesize\raggedright cited publication, lines 3615-\allowbreak{}3697\par}
\paragraph{Verbatim source input}
\begin{ReviewVerbatim}
i.e. with w(θ1 , θ2 ) = (θ1 − θ2 ).
Proof. Recall Pk (θ1 , θ3 ) =
P r((θ1 − θ2 )(xk1 − xk3 ) > 0)
Z
Z
=
P r(xk1 − xk3 > 0)dθ1 dθ2 dθ3 +
θ1 >θ2 ,θ3

Z
=

P r(xk1 − xk3 > 0)

R

θ1 >θ2 (θ1 −θ2 )Pk (θ1 , θ2 )dθ1 dθ2 ,

P r(xk1 − xk3 < 0)dθ1 dθ2 dθ3

θ1 <θ2 ,θ3

[U+0014 control character]Z θ1

[U+0015 control character]

Z

P r(xk1 − xk3 < 0)

[U+0014 control character]Z 1

[U+0015 control character]

dθ2 dθ1 dθ3 +
dθ2 dθ1 dθ3
θ ,θ
θ2 =0
θ1 ,θ3
θ2 =θ1
Z 1 3h
i
=
P r(xk1 − xk3 > 0)θ1 ] + P r(xk1 − xk3 < 0)(1 − θ1 ) dθ1 dθ3
θ ,θ
Z 1 3h
h
ii
=
P r(xk1 − xk3 < 0) + θ1 P r(xk1 − xk3 > 0) − P r(xk1 − xk3 < 0) dθ1 dθ3
θ1 ,θ3

24


[form-feed in source extraction]
Similarly,
P r((θ1 − θ2 )(xk1 − xk3 ) < 0) =
Z
h
h
ii
=
P r(xk1 − xk3 > 0) + θ1 P r(xk1 − xk3 < 0) − P r(xk1 − xk3 > 0) dθ1 dθ3
θ ,θ
Z 1 3h
h
ii
=
P r(xk3 − xk1 > 0) + θ3 P r(xk3 − xk1 < 0) − P r(xk3 − xk1 > 0) dθ1 dθ3
θ1 ,θ3

Where the second equality follows from θ1 , θ3 interchangeable. Then
Z
Wkρ = 3
(θ1 − θ2 )Pk (θ1 , θ2 )dθ1 dθ2
θ1 ,θ2
Z
6(θ1 − θ2 )Pk (θ1 , θ2 )dθ1 dθ2
=
θ1 >θ2

Note that Spearman’s ρ is similar to Kendall’s τ with an additional weighting for how far apart
the two values that are flipped are.
Lemma C.11. When w is constant, i.e. for Kendall’s τ rank correlation, the intervals s that
\end{ReviewVerbatim}


\paragraph{Expanded PaperInterface specification}
\begin{ReviewVerbatim}
∀ (M : ℕ) (s : ℕ → ℝ),
  s 0 = 0 →
    s M = 1 →
      GJ19OptimalBinaryRatingSystems.spearmanLinearWeightOrderedPairIntervalObjective M s =
        GJ19OptimalBinaryRatingSystems.spearmanLinearWeightIntervalObjective M s
\end{ReviewVerbatim}

\paragraph{Lean proof endpoint (built separately)}\begin{ReviewVerbatim}
GJ19OptimalBinaryRatingSystems.PaperInterface.paper_lemmaC10_spearman_linear_weight_ordered_pair_interval_objective_eq_gap_objective
\end{ReviewVerbatim}

\paragraph{Recorded source-to-Spec screening}
\textbf{Verdict:} \texttt{not recorded}
\\
\textbf{Reason:} not recorded
\reviewmatch{Matches source input}
\noindent\textbf{Reviewer annotation}\par
\reviewerbox
\clearpage
\hypertarget{source-claim-75f334c9ed5f4c67}{}
\section*{Review row 58}
\textbf{Source locator:} {\footnotesize\raggedright cited publication, lines 3698-\allowbreak{}3712\par}
\paragraph{Verbatim source input}
\begin{ReviewVerbatim}
maximize (20),
X
X Z
w(θ1 , θ2 )d(θ1 , θ2 ) =
(si+1 − si )(sj+1 − sj )
(27)
0≤i<j<M

θ2 ∈Si ,θ1 ∈Sj

0≤i<j<M

, are {si = Mi }M
i=0 .
Lemma C.12. When w is (θ1 − θ2 ), i.e. for Spearman’s ρ rank correlation, the intervals s that
\end{ReviewVerbatim}


\paragraph{Expanded PaperInterface specification}
\begin{ReviewVerbatim}
∀ {M : ℕ} [Nonempty (Fin M)] (s : ℕ → ℝ),
  s 0 = 0 → s M = 1 → (∑ i, ∑ j, if i < j then (s (↑i + 1) - s ↑i) * (s (↑j + 1) - s ↑j) else 0) ≤ (1 - (↑M)⁻¹) / 2
\end{ReviewVerbatim}

\paragraph{Lean proof endpoint (built separately)}\begin{ReviewVerbatim}
GJ19OptimalBinaryRatingSystems.PaperInterface.paper_lemmaC11_kendall_constant_weight_ordered_pair_interval_objective_le_equispaced
\end{ReviewVerbatim}

\paragraph{Recorded source-to-Spec screening}
\textbf{Verdict:} \texttt{not recorded}
\\
\textbf{Reason:} not recorded
\reviewmatch{Matches source input}
\noindent\textbf{Reviewer annotation}\par
\reviewerbox
\clearpage
\hypertarget{source-claim-9e5abbd5563b8516}{}
\section*{Review row 59}
\textbf{Source locator:} {\footnotesize\raggedright cited publication, lines 3713-\allowbreak{}3760\par}
\paragraph{Verbatim source input}
\begin{ReviewVerbatim}
maximize (20),
X Z
w(θ1 , θ2 )d(θ1 , θ2 )
(28)
0≤i<j<M

θ2 ∈Si ,θ1 ∈Sj

are {si = Mi }M
i=0 , i.e. the same as those for Kendall’s τ .
Proof.
X
0≤i<j<M

Z
w(θ1 , θ2 )d(θ1 , θ2 ) =
θ2 ∈Si ,θ1 ∈Sj

Z

X

(θ1 − θ2 )d(θ1 , θ2 )

0≤i<j<M

=

θ2 ∈Si ,θ1 ∈Sj

[U+0012 control character]

X
0<i<j≤M

sj + sj−1 si + si−1
−
2
2

[U+0013 control character]
(si − si−1 )(sj − sj−1 )

Finding an asymptotically optimal {si } then is a constrained third order polynomial maximization
i=M , as for Kendall’s tau
problem with M variables. The maximum is achieved at {si = Mi }i=0
correlation.
\end{ReviewVerbatim}


\paragraph{Expanded PaperInterface specification}
\begin{ReviewVerbatim}
∀ {M : ℕ} [Nonempty (Fin M)] (s : ℕ → ℝ),
  Monotone s →
    s 0 = 0 →
      s M = 1 →
        GJ19OptimalBinaryRatingSystems.spearmanLinearWeightOrderedPairIntervalObjective M s ≤ (1 - (↑M)⁻¹ ^ 2) / 6
\end{ReviewVerbatim}

\paragraph{Lean proof endpoint (built separately)}\begin{ReviewVerbatim}
GJ19OptimalBinaryRatingSystems.PaperInterface.paper_lemmaC12_spearman_linear_weight_ordered_pair_interval_objective_le_equispaced
\end{ReviewVerbatim}

\paragraph{Recorded source-to-Spec screening}
\textbf{Verdict:} \texttt{not recorded}
\\
\textbf{Reason:} not recorded
\reviewmatch{Matches source input}
\noindent\textbf{Reviewer annotation}\par
\reviewerbox

\end{document}
